% Generated self-contained reading/editing copy. Modular master: proofs_compendium.tex
% !TeX program = pdflatex
\documentclass[11pt]{article}
\usepackage[T1]{fontenc}
\usepackage[utf8]{inputenc}
\usepackage{amsmath,amssymb,amsthm}
\usepackage[margin=1in]{geometry}
\usepackage{longtable,array,booktabs,enumitem}
\usepackage[hidelinks,hypertexnames=false]{hyperref}
% Shared notation and theorem infrastructure only. Load amsmath, amssymb,
% and amsthm in the root document before this file. No document class here.
\newcommand{\R}{\mathbb{R}}
\newcommand{\C}{\mathbb{C}}
\newcommand{\Z}{\mathbb{Z}}
\newcommand{\Zn}{\mathbb{Z}_n}
\newcommand{\one}{\mathbf{1}}
\newcommand{\ii}{\mathrm{i}}
\newcommand{\eps}{\varepsilon}
\newcommand{\DM}{\mathrm{DM}}
\newcommand{\MB}{\mathrm{MB}}
\newcommand{\LB}{\mathrm{LB}}
\newcommand{\Sync}{\mathcal{S}}
\newcommand{\Orth}{\R_{\geq 0}^{n}}
\newcommand{\PosOrth}{\R_{>0}^{n}}
\newcommand{\Domain}{\mathcal{U}}
\newcommand{\param}{\boldsymbol{\pi}}
\newcommand{\Fourier}{\mathsf{Q}}
\newcommand{\conj}[1]{\overline{#1}}
\newcommand{\inner}[2]{\langle #1,#2\rangle}
\newcommand{\norm}[1]{\left\lVert #1\right\rVert}
\DeclareMathOperator{\spec}{spec}
\DeclareMathOperator{\sa}{sa}
\DeclareMathOperator{\diag}{diag}
\DeclareMathOperator{\Ree}{Re}
\DeclareMathOperator{\Imm}{Im}
\DeclareMathOperator*{\argmax}{arg\,max}
\newtheorem{theorem}{Theorem}[section]
\newtheorem{lemma}[theorem]{Lemma}
\newtheorem{proposition}[theorem]{Proposition}
\newtheorem{corollary}[theorem]{Corollary}
\theoremstyle{definition}
\newtheorem{definition}[theorem]{Definition}
\theoremstyle{remark}
\newtheorem{remark}[theorem]{Remark}

\hypersetup{pdftitle={Directed Cyclic Networks: Shortened Proof Compendium},pdfauthor={}}
\setcounter{tocdepth}{1}
\setlength{\emergencystretch}{2em}
\title{Directed Cyclic Networks\\Shortened Proof Compendium}
\author{}
\date{4 October 2026}
\begin{document}
\maketitle
\noindent This edition replaces repeated standard arguments with citations
and retains the model-specific coefficients, parameter constructions,
spectral gaps and realizability proofs. P01--P05 and P07--P11 contain the
completed arguments supplied in the source project. Earlier independent
audits apply to that source edition, not automatically to this edited text.
The supplied P08 global-interior extension also awaited independent audit.

\medskip\noindent P06 and P12--P18 were outlines in the source.
They are retained as concise, explicitly conditional frameworks with
their outstanding calculations identified. In particular P11 establishes
a spectral Hopf candidate, not a completed nonlinear Hopf theorem.

\medskip\noindent All ten cited works were checked against primary
publisher or author records. Full-text access and theorem-location limits
are recorded in \path{CITATION_VERIFICATION.txt}. Research articles,
reference articles, a preprint and one standard textbook are distinguished
there. No unverified article-specific theorem number is used.
\clearpage
\tableofcontents
\clearpage
% STATUS: CANONICAL SPECIFICATION; no theorem is asserted or proved here.
\section{Canonical model contract}
\label{sec:canonical}
This section specifies the equations and parameter domain used throughout.

\subsection{Main demand-modulated model}
For $i\in\Zn$, define the vector field $F^{\DM}$ by
\begin{equation}
\label{eq:canonical-dm}
\dot{x}_i=a x_i-kx_i^2+\frac{c x_{i-1}}{1+x_i}
             -\frac{g x_i}{1+x_{i+q}}-h x_i x_{i-q}.
\tag{$\DM_{n,q}$}
\end{equation}
The variables are dimensionless normalized activities. The preferred term is
\emph{demand-modulated model}. The positive input is called demand-modulated
predecessor input; no standard Holling type-II classification is adopted.
The intended initial conditions are $x(0)\in\Orth$. For local statements and
boundary bifurcations, use the specified open extension domain
$\Domain=(-1,\infty)^n$. This domain choice asserts no invariance theorem.

\subsection{Primary matched bilinear comparator}
Define $F^{\MB}$ by
\begin{equation}
\label{eq:canonical-mb}
\dot{y}_i=a y_i-k y_i^2+c y_i y_{i-1}
                 -g y_i y_{i+q}-h y_i y_{i-q},\qquad i\in\Zn.
\tag{$\MB_{n,q}$}
\end{equation}
This is the topology- and parameter-count-matched comparator in the main
program. Its intended initial data are nonnegative; the polynomial vector
field is specified on $\R^n$. No existence, invariance, boundedness or
stability conclusion is included in this contract. Matching parameter names
does not settle a quantitative calibration between the two laws.

\subsection{Parameter and indexing contract}
\[
n\in\Z,\quad n\geq5,\qquad q\in\{1,\ldots,n-1\},\qquad
\param=(a,k,c,g,h)\in(0,\infty)^5.
\]
All indices are modulo $n$, with representatives $0,\ldots,n-1$.
Arbitrary $q$ is allowed. The interpretation as one connected step-$q$ cycle
requires $\gcd(n,q)=1$; this condition is not imposed on the general analytical
program. In particular, even $n$ is not excluded when $q=2$. Coincident offsets
do not remove terms from either defining equation.
The intended exact wave-number-selection program focuses on $q=2$.
The five-node Wu-Xing-inspired specialization is $n=5,q=2$.

\subsection{Archived secondary six-term comparator}
The following is a transcription of the blueprint's proposed arbitrary-$n,q$
extension, not a verified attribution to the original Lin--Han publication:
\begin{equation}
\label{eq:archived-lb}
\dot{y}_i=a y_i-b y_i^2+c y_i y_{i-1}-d y_i y_{i-q}
                    -g_{\mathrm c}y_i y_{i+1}-h y_i y_{i+q}.
\tag{$\LB_{n,q}$ archive}
\end{equation}
Its archived five-node equation is obtained by specifying $n=5,q=2$ in this
transcription. The six coefficient names are archival and are not a parameter
identification with $\DM$ or $\MB$. The original source, coefficient domain,
index orientation, and attribution remain to be checked before analysis or
use as a literature application.

\subsection{Scope boundary}
The main program covers homogeneous $\DM$, primary comparator $\MB$, their
synchronous equilibria and spectra, the planned bifurcations, and the planned
small-heterogeneity response theory. All requested P01--P18 tasks remain in
the project; final main-text versus supplement placement is open.
The heterogeneous coefficient assignment is not frozen by this contract.
The six-term model is archived for a possible secondary literature application,
never a replacement for $\MB$. Source-saturating replacements, local motif
theory, pathology extensions, adaptive set points, interventions, noise,
delays, data calibration, and prior Lyapunov-spectrum reanalysis are postponed
to other work. 

% STATUS: CANONICAL NOTATION AND ASSUMPTIONS; no derived spectral conclusions.
\section{Notation and assumptions}
\label{sec:notation}
\subsection{Standing assumptions}
\begin{description}
\item[A1.] $n\in\Z$, $n\geq5$, and $q\in\{1,\ldots,n-1\}$.
\item[A2.] Homogeneous parameters $\param=(a,k,c,g,h)\in(0,\infty)^5$.
\item[A3.] Dimensionless real activities; initial states in the nonnegative
orthant when forward biological-admissibility statements are sought.
For $\DM$ the local extension domain is $\Domain=(-1,\infty)^n$.
\item[A4, conditional.] Impose $\gcd(n,q)=1$ only for the connected step-$q$
cycle interpretation; specify $q=2$ or $(n,q)=(5,2)$ when required by a module.
\item[A5, conditional.] Positivity of an equilibrium, invertibility,
hyperbolicity, stability, Hopf nondegeneracy, and perturbation regularity are
additional hypotheses to be stated and checked in the appropriate future proof.
None is a standing established result.
\end{description}

\subsection{Indices and cyclic shifts}
Define $\Zn=\Z/n\Z$, using representatives $0,\ldots,n-1$.
For $z\in\C^n$, define the predecessor and successor shift matrices by
\[
(Pz)_j=z_{j-1},\qquad(P^{-1}z)_j=z_{j+1},\qquad
P_{j\ell}=\begin{cases}1,&\ell\equiv j-1\pmod n,\\0,&\text{otherwise.}\end{cases}
\]
Thus the symbols $P^q$ and $P^{-q}$ denote repeated predecessor and successor
shifts, respectively. The group-action convention is $\rho(r)z=P^r z$.
For transcription of the original five-node labels, the original $x_1,\ldots,x_5$
are relabeled as $x_0,\ldots,x_4$ in the same cyclic order. This is a label
translation only; no historical or physiological statement is certified.

\subsection{Fourier vectors, transform and pairing}
Use $\ii$ for the imaginary unit and $j$ for a node index. Define
\[
\theta_m=\frac{2\pi m}{n},\qquad
v^{(m)}_j=\frac{1}{\sqrt n}\exp(\ii\theta_m j),\qquad
\Fourier_{jm}=v^{(m)}_j,\quad 0\leq j,m<n.
\]
The canonical transform convention is
\[
\widehat z_m=\frac{1}{\sqrt n}\sum_{j=0}^{n-1}
    z_j\exp(-\ii\theta_m j),\qquad
z_j=\frac{1}{\sqrt n}\sum_{m=0}^{n-1}
    \widehat z_m\exp(\ii\theta_m j).
\]
The normalization is the unitary DFT convention. The blueprint's displayed
unscaled Fourier vectors are not copied as normalized eigenvectors.
Use $A^*$ for conjugate transpose, $A^{\mathsf T}$ for transpose, and
\[
\inner{u}{v}=u^*v=\sum_j\overline{u_j}v_j,
\qquad \norm{z}_2=(z^*z)^{1/2}.
\]
The pairing is conjugate-linear in its first argument. Matrix $2$-norms are
induced Euclidean norms. The adjoint eigenvector equation, its phase/scale,
Taylor tensors and first-Lyapunov-coefficient convention must be reconciled
in P13 before any blueprint normal-form formula is imported.

\subsection{State sets and equilibrium notation}
\[
\Orth=\{x\in\R^n:x_j\geq0\ \text{for all }j\},\qquad
\PosOrth=\{x\in\R^n:x_j>0\ \text{for all }j\},
\]
\[
\one=(1,\ldots,1)^{\mathsf T},\qquad
\Sync=\{s\one:s\in\R\}.
\]
The phrase positive orthant means the strict set $\PosOrth$; nonnegative
orthant means $\Orth$. A synchronous state is a member of $\Sync$; a positive
synchronous equilibrium, if established, is denoted $x^*\one$ with $x^*>0$.
The admissible synchronous domain for $\DM$ is $\Sync\cap\Domain$.
No invariance or equilibrium classification is asserted by these definitions.

\subsection{Spectrum, hyperbolicity and mode terminology}
For a square matrix $J$, define
\[
\sa(J)=\max\{\Ree\lambda:\lambda\in\spec(J)\}.
\]
An equilibrium of a differentiable vector field is hyperbolic when its
Jacobian has no eigenvalue with zero real part. A matrix is Hurwitz if all
its eigenvalues have negative real part, and normal if $JJ^*=J^*J$.
These definitions assert neither stability nor normality for a model.
Mode $m=0$ is called the uniform or synchronous mode; modes $m\neq0$ are
called nonuniform or symmetry-breaking modes. The mode-pair notation is
$\{m,(-m)\bmod n\}$, usually written $\{m,n-m\}$. A pair representative is
chosen in $0\leq m\leq\lfloor n/2\rfloor$. Keep the uniform mode and, for
even $n$, $m=n/2$ as singleton labels. Any eigenvalue-conjugacy or simplicity
claim belongs to P04 or P11.
When a spectrum has been derived, the selected-mode set will mean
$\argmax_m\Ree\lambda_m$, including every tie. A first unstable mode additionally
requires a specified equilibrium branch and oriented parameter path.

\subsection{Reserved scalar notation}
For continuity with the archive, reserve $A=k+h$ and $\mu=a+c-g$ as
abbreviations, without asserting an equilibrium equation or bifurcation.
Use $\nu$ for a future scalar continuation parameter and $\eps$ for a
heterogeneity amplitude. Reserve $\omega>0$ for a chosen positive Hopf
frequency. Use $\sa(J)$ for spectral abscissa and $\delta=-\sa(J)>0$ only
under an explicitly stated Hurwitz hypothesis. This avoids reusing
$\alpha$ as both a stability margin and a parameter perturbation.
The spectral coefficient called $T_0$ in the blueprint will be named
$\tau$ if needed, keeping $T_{\mathrm{scale}}$ available for a physical
time scale. Heterogeneity profiles will be denoted collectively by
$\boldsymbol\zeta$ if adopted; no coefficient-assignment convention is implied.


% EDITION: citation-based shortening, 4 October 2026.
% Earlier audits refer to the source edition; this revision has not been independently audited.
\section{P01: Well-posedness and bounded absorbing sets}
\label{sec:P01}

\subsection{Domain, hypotheses, and set terminology}
Use exactly the demand-modulated vector field in
\eqref{eq:canonical-dm}, denoted here by $F=(F_i)_{i\in\Zn}$.
The canonical assumptions imply the weaker hypotheses
\begin{equation}
\label{eq:P01-weak}
 n\ge1,\quad q\in\Z\pmod n,\quad
 a\in\R,\quad k>0,\quad c,g,h\ge0.
\tag{P01-W}
\end{equation}
We prove the main results under \eqref{eq:P01-weak}, without changing the
canonical contract for later modules. Coincident offsets retain every
term of the defining equation. Write
\[
 \Domain=(-1,\infty)^n,\qquad K_r=[0,r]^n\ (r\ge0),\qquad
 X(t)=\max_{i\in\Zn}x_i(t),\qquad s_+=\max\{s,0\}.
\]
For a continuous real function $v$, define
\[
 D^+v(t)=\limsup_{\tau\downarrow0}
             \frac{v(t+\tau)-v(t)}{\tau}.
\]
Distances use the maximum norm; for $x\in\Orth$,
\begin{equation}
\label{eq:P01-distance}
 \operatorname{dist}_\infty(x,K_r)=(\max_i x_i-r)_+.
\end{equation}
A set is \emph{positively invariant} if solutions starting there remain
there for all forward times for which they exist. It \emph{absorbs bounded
subsets of $\Orth$} if every bounded $E\subset\Orth$ has a finite $T_E$
such that all solutions starting in $E$ lie in the set for every
$t\ge T_E$. Attraction means convergence of the distance to zero;
uniform attraction of bounded sets means convergence uniform over each
such $E$. We will also establish stronger absorption bounds with a
single time valid for the entire orthant.

\subsection{Finite maxima, local solutions, and coordinate boundaries}
\label{slot:P01-1}

\begin{lemma}[Finite maxima, including ties]
\label{lem:P01-max}
For finitely many continuously differentiable functions $y_j$, set
$v(t)=\max_j y_j(t)$ and $I(t)=\{j:y_j(t)=v(t)\}$.
Then $v$ is locally Lipschitz, its right derivative exists at every
time with a right neighborhood in the interval, and
\begin{equation}
\label{eq:P01-max}
 D^+v(t)=\max_{j\in I(t)}\dot y_j(t).
\end{equation}
\end{lemma}
\begin{proof}
Bounded derivatives on compact time intervals make the finite maximum
Lipschitz. At a fixed $t$, every nonmaximizing index has a positive gap
from $v(t)$. Finiteness and continuity ensure that none of these indices
maximizes at $t+\tau$ for sufficiently small $\tau>0$. Uniformly over
$I(t)$,
$y_j(t+\tau)=v(t)+\tau\dot y_j(t)+o(\tau)$.
Taking the maximum and dividing by $\tau$ proves the assertion.
\end{proof}

\begin{proposition}[Local well-posedness and nonnegative invariance]
\label{prop:P01-local}
For any finite real coefficients, $F$ is $C^\infty$ on $\Domain$.
Every $x^0\in\Domain$ has a unique maximal classical forward solution
$x:[0,T_{\max})\to\Domain$, with continuous dependence on initial data.
If $c\ge0$ and $x^0\in\Orth$, then $x(t)\in\Orth$ for $t<T_{\max}$,
regardless of the signs of the other coefficients.
\end{proposition}
\begin{proof}
The denominators $1+x_i$ and $1+x_{i+q}$ are strictly positive on
$\Domain$. Each component of $F$ is a sum of polynomials and rational
functions with nonvanishing denominators, hence is smooth and locally
Lipschitz. Standard local ODE existence, uniqueness, and maximal
continuation theory on an open domain applies.

On each coordinate face,
\begin{equation}
\label{eq:P01-face}
 x_i=0,\quad x\in\Orth\quad\Longrightarrow\quad F_i(x)=c x_{i-1}\ge0.
\end{equation}
To treat tangent faces rigorously, fix $T<T_{\max}$ and let $x^+(t)$
be the coordinatewise positive part of $x(t)$. The curves $x,x^+$ and
their connecting line segments lie in a compact rectangle in $\Domain$.
Choose a Lipschitz constant $L$ for $F$ on this rectangle, in the
maximum norm, and put
\[
 V(t)=\max\{0,-x_0(t),\ldots,-x_{n-1}(t)\}
      =\norm{x(t)-x^+(t)}_\infty.
\]
If $V>0$ and $-x_i=V$, then $x_i^+=0$, so
$F_i(x^+)\ge0$ and $-F_i(x)\le |F_i(x)-F_i(x^+)|\le LV$.
If $V=0$, any active coordinate has $x_i=0$ and $-F_i(x)\le0$;
the constant zero is also active. Lemma~\ref{lem:P01-max} gives
$D^+V\le LV$. Since $V$ is Lipschitz, this inequality holds for its
ordinary derivative almost everywhere. Integration of
$(e^{-Lt}V)' \le0$, with $V(0)=0$, gives $V=0$ on $[0,T]$.
The choice of $T$ was arbitrary.
\end{proof}

\begin{proposition}[Persistence and propagation of positivity]
\label{prop:P01-positive}
For $c\ge0$ and nonnegative initial data, every initially positive
coordinate remains strictly positive on $[0,T_{\max})$. In particular,
strictly positive initial data stay strictly positive. If $c>0$ and
$x^0\ne0$, all coordinates are strictly positive at every
$0<t<T_{\max}$. If $c=0$, initially zero coordinates remain zero.
The zero initial state remains identically zero.
\end{proposition}
\begin{proof}
Along a nonnegative solution write
\[
 \dot x_i=b_i(t)x_i+d_i(t),\qquad
 b_i=a-kx_i-\frac{g}{1+x_{i+q}}-h x_{i-q},\qquad
 d_i=\frac{c x_{i-1}}{1+x_i}\ge0.
\]
These coefficients are continuous on every compact time interval in
the maximal interval. The integrating-factor formula gives
\begin{equation}
\label{eq:P01-positive}
 x_i(t)=x_i^0 e^{\int_0^t b_i(s)\,ds}
       +\int_0^t e^{\int_s^t b_i(r)\,dr}d_i(s)\,ds.
\end{equation}
This proves persistence of an initially positive coordinate.
If $c>0$, a predecessor positive at every $s>0$ supplies a strictly
positive integrand to its successor at every $s>0$, making that successor
positive at every $t>0$. At most $n-1$ iterations along the step-one
cycle reach all nodes, independently of $q$. For $c=0$ the integral
vanishes, proving preservation of zero coordinates. Finally $F(0)=0$
and uniqueness proves the zero-data assertion.
\end{proof}

\subsection{Maximum comparison and the exact comparison radius}
\label{slot:P01-2}
Under \eqref{eq:P01-weak}, define
\[
 b(s)=a-ks+\frac{c}{1+s},\qquad H(s)=s b(s)\quad(s\ge0),
\]
and set
\begin{equation}
\label{eq:P01-radius}
 X_{\max}=R:=\max\left\{0,\frac{a-k+\sqrt{(a+k)^2+4kc}}{2k}\right\}.
\end{equation}
This radius is defined by a scalar upper comparison; it is not asserted
to be the smallest invariant cube of the full vector field.

\begin{lemma}[Dini inequality and scalar root]
\label{lem:P01-bound}
Every nonnegative solution under \eqref{eq:P01-weak} satisfies
\begin{equation}
\label{eq:P01-dini}
 D^+X=\max_{i:x_i=X}F_i(x)
       \le H(X)-\frac{gX}{1+X}\le H(X).
\end{equation}
If $a+c>0$, then $R>0$ is the unique positive zero of $b$, and
\begin{equation}
\label{eq:P01-polynomial}
 kR^2+(k-a)R-(a+c)=0.
\end{equation}
If $a+c\le0$, then $R=0$ and $H(s)\le-ks^2$ for $s\ge0$.
For $R>0$ and $s\ge R$,
\begin{equation}
\label{eq:P01-factor}
 H(s)=s\left[-k(s-R)
    +c\left(\frac1{1+s}-\frac1{1+R}\right)\right]\le-ks(s-R).
\end{equation}
\end{lemma}
\begin{proof}
Lemma~\ref{lem:P01-max} includes every tied maximizing coordinate.
At $x_i=X$, nonnegative invariance gives
$0\le x_{i-1},x_{i+q},x_{i-q}\le X$. Consequently
\[
 \frac{c x_{i-1}}{1+X}\le\frac{cX}{1+X},\qquad
 -\frac{gX}{1+x_{i+q}}\le-\frac{gX}{1+X},\qquad
 -hXx_{i-q}\le0.
\]
The middle inequality uses the negative prefactor and $g\ge0$.
Adding proves \eqref{eq:P01-dini}; no independence of neighbor
coordinates is required.

Now $b'(s)=-k-c/(1+s)^2<0$, $b(0)=a+c$, and $b(s)\to-\infty$.
If $a+c>0$, multiply $b(R)=0$ by $1+R$ to obtain
\eqref{eq:P01-polynomial}. Its constant term is negative and its
discriminant is $(a+k)^2+4kc$, so its positive root is
\eqref{eq:P01-radius}. If $a+c\le0$, then $k-a>0$ and
\[
 (a+k)^2+4kc-(k-a)^2=4k(a+c)\le0.
\]
Thus the raw root in \eqref{eq:P01-radius} is nonpositive, $R=0$,
and $b(s)\le a+c-ks\le-ks$. Subtracting $b(R)=0$ when $R>0$
gives \eqref{eq:P01-factor}; its inequality follows from $c\ge0$
and $s\ge R$.
\end{proof}

\begin{theorem}[Global continuation and invariant attracting cubes]
\label{thm:P01-global}
Under \eqref{eq:P01-weak}, every $x^0\in\Orth$ has a unique classical
solution for all $t\ge0$. The scalar solution
\begin{equation}
\label{eq:P01-scalar}
 \dot u=H(u),\qquad u(0)=X(0)
\end{equation}
is global and nonnegative, and
\begin{equation}
\label{eq:P01-globalbound}
 0\le X(t)\le u(t)\le\max\{X(0),R\}.
\end{equation}
Every $K_r$ with $r\ge R$ is positively invariant. Furthermore,
\begin{align}
 \operatorname{dist}_\infty(x(t),K_R)
   &\le e^{-kRt}(X(0)-R)_+ &&(R>0), \label{eq:P01-distance-rate}\\
 X(t)&\le\frac{X(0)}{1+kX(0)t} &&(R=0). \label{eq:P01-zero-rate}
\end{align}
There are also initial-state-independent bounds, for $t>0$:
\begin{align}
 X(t)&\le\frac{R}{1-e^{-kRt}} &&(R>0), \label{eq:P01-comingdown}\\
 X(t)&\le\frac1{kt} &&(R=0). \label{eq:P01-comingdown-zero}
\end{align}
Thus $K_R$ attracts all nonnegative solutions uniformly on bounded
initial sets, and indeed uniformly over the entire orthant as
$t\to\infty$. Every larger cube $K_r$, $r>R$, absorbs the entire
orthant by the common time
\begin{equation}
\label{eq:P01-outer-time}
 T_r=\begin{cases}
 \displaystyle\frac1{kR}\log\frac{r}{r-R},&R>0,\\[4pt]
 \displaystyle\frac1{kr},&R=0.
 \end{cases}
\end{equation}
\end{theorem}
\begin{proof}
The scalar field $H$ is smooth on $(-1,\infty)$, $H(0)=0$, and
$H(s)\le0$ for $s\ge R$. Its solution starting nonnegative stays in
$[0,M]$, where $M=\max\{u(0),R\}$: uniqueness prevents crossing zero,
and the upper endpoint has nonpositive field. The latter barrier also
follows by the positive-part argument below applied to $(u-M)_+$,
using $H(M)\le0$. Compact containment gives global scalar continuation.

We justify the vector comparison without assuming differentiability of
$X$ everywhere. On a compact subinterval of the vector maximal interval,
let $L$ be a Lipschitz constant for $H$ on an interval containing $X$
and $u$. The Lipschitz function
\[
 W=(X-u)_+=\max\{0,x_i-u:i\in\Zn\}
\]
satisfies $D^+W\le H(X)-H(u)\le LW$ where $W>0$.
At $W=0$, any active coordinate has $x_i=u=X$, unless the constant
zero is the only active function; Lemma~\ref{lem:P01-max} and
\eqref{eq:P01-dini} give $D^+W\le0$. Integrating the almost-everywhere
inequality, as in Proposition~\ref{prop:P01-local}, yields $W=0$.
This proves \eqref{eq:P01-globalbound} on the maximal interval.

If $T_{\max}<\infty$, the vector trajectory stays in $K_M$, a compact
subset of $\Domain$ at distance at least one from every singular
hyperplane. The field is bounded on a compact neighborhood in $\Domain$,
hence $\dot x$ is bounded. Thus $x(t)$ has a limit in $K_M$ as
$t\uparrow T_{\max}$. Local existence at that limit and uniqueness
extend the solution past $T_{\max}$, a contradiction. This proves
global continuation and absence of finite-time blow-up. If
$X(0)\le r$ and $r\ge R$, \eqref{eq:P01-globalbound} proves invariance.

For $R>0$, \eqref{eq:P01-factor} gives
$\dot u\le-ku(u-R)\le-kR(u-R)$ where $u\ge R$. It follows that
$(u(t)-R)_+\le e^{-kRt}(u(0)-R)_+$, including solutions initially at
or below $R$, and hence \eqref{eq:P01-distance-rate}. If $u(0)>R$,
set $w=1-R/u$ while $u>R$. Then
\[
 \dot w=\frac{R\dot u}{u^2}\le-kR w,\qquad 0<w(0)<1.
\]
Solving this inequality gives \eqref{eq:P01-comingdown} for $u$ and
therefore for $X$. If $u$ reaches $R$, invariance supplies the same
bound afterwards; if $u(0)\le R$, it is immediate.
For $R=0$, integrating $(1/u)'\ge k$ where $u>0$, using
$\dot u\le-ku^2$, proves \eqref{eq:P01-zero-rate} and
\eqref{eq:P01-comingdown-zero}; zero initial data follow by uniqueness.

Attraction follows from these estimates and \eqref{eq:P01-distance}.
For $R>0$ the distance is at most $R/(e^{kRt}-1)$ independently of
initial data; for $R=0$ use $1/(kt)$. Substitution of $T_r$ into
\eqref{eq:P01-comingdown} or \eqref{eq:P01-comingdown-zero} gives
$X(T_r)\le r$, and invariance supplies all subsequent times.
\end{proof}

\subsection{Finite-time absorption at the comparison radius}
\label{slot:P01-3}

\begin{theorem}[Exact absorption criterion for $K_R$]
\label{thm:P01-absorption}
Assume \eqref{eq:P01-weak}. If $R=0$, no initially nonzero nonnegative
solution enters $K_R=\{0\}$ at finite time. This cube is invariant and
globally attracting but not absorbing.

If $R>0$, the following are equivalent:
\begin{enumerate}[label=(\roman*),leftmargin=*]
\item $K_R$ absorbs every nonnegative solution in finite time.
\item $K_R$ absorbs bounded subsets of $\Orth$.
\item One finite absorption time is valid for the entire orthant.
\item The coefficients and offset satisfy
\begin{equation}
\label{eq:P01-absorption-criterion}
 g>0\quad\hbox{or}\quad
 [\,h>0\ \hbox{and}\ (c>0\ \hbox{or}\ q\equiv0\pmod n)\,].
\end{equation}
\end{enumerate}
Canonically $R>0$, $g>0$, and the entire orthant enters $K_R$ by
\begin{equation}
\label{eq:P01-canonical-time}
 T_*=\frac{\log2}{kR}+\frac{1+R}{g}.
\end{equation}
If $g>0$ and $X(0)\le M$, an alternative entry-time bound is
\begin{equation}
\label{eq:P01-bounded-time}
 T_M=\frac{(M-R)_+(1+R)}{gR}.
\end{equation}
\end{theorem}
\begin{proof}
When $R=0$, Proposition~\ref{prop:P01-positive} ensures that some
coordinate of any initially nonzero solution stays positive at each
finite time.

Now suppose $R>0$ and set $t_0=(\log2)/(kR)$. By
\eqref{eq:P01-comingdown}, every trajectory is in the invariant cube
$K_{2R}$ by $t_0$. If $g>0$, then whenever $X\ge R$,
\begin{equation}
\label{eq:P01-strict-gap}
 D^+X\le-\frac{gX}{1+X}\le-\frac{gR}{1+R}=:-\delta<0,
\end{equation}
using $H(X)\le0$. Absolute continuity permits integration before
the first entry into $X\le R$. If $X(0)\le M$, a longer time outside
than $(M-R)_+/\delta$ would force $X\le R$, a contradiction.
Invariance gives all later times. Applying this after $t_0$ with
$M=2R$ proves \eqref{eq:P01-canonical-time}.

If $g=0$, $h>0$, and $q\equiv0\pmod n$, a maximizing coordinate
has $x_{i-q}=X$. Thus $D^+X\le H(X)-hX^2\le-hR^2$ when $X\ge R$,
giving common absorption by $t_0+1/(hR)$.

For the remaining sufficient case $g=0$, $h>0$, $c>0$, we give a
coordinatewise quantitative argument. Put $M=2R$ and fix any
$\tau>0$. On $K_M$,
\begin{equation}
\label{eq:P01-propagation-bound}
 \dot x_i\ge-\beta x_i+\eta x_{i-1},\qquad
 \beta=(-a)_++(k+h)M+g,\quad \eta=\frac{c}{1+M}>0.
\end{equation}
Indeed $a\ge-(-a)_+$, $kx_i\le kM$,
$g/(1+x_{i+q})\le g$, $h x_{i-q}\le hM$, and
$c/(1+x_i)\ge c/(1+M)$. On $[s,s+\tau]$ in this cube, the transformed
variables $z_i(r)=e^{\beta(r-s)}x_i(r)$ satisfy
\[
 z_i(r)\ge x_i(s)+\eta\int_s^r z_{i-1}(v)\,dv.
\]
Repeated substitution, keeping the remaining nested integral
nonnegative, yields
\begin{equation}
\label{eq:P01-propagation-series}
 x_i(s+\tau)\ge e^{-\beta\tau}
       \sum_{j=0}^{n-1}\frac{(\eta\tau)^j}{j!}x_{i-j}(s).
\end{equation}
A $j$-fold integral of a constant over the time simplex equals
$\tau^j/j!$, which gives the displayed coefficients.

For a trajectory that has not entered $K_R$ by time $t\ge t_0+\tau$,
we have $X(t-\tau)>R$. The indices $i-j$, $0\le j\le n-1$, cover
the cycle once. Thus \eqref{eq:P01-propagation-series} gives, for every $i$,
\begin{equation}
\label{eq:P01-positive-floor}
 x_i(t)\ge d:=R e^{-\beta\tau}
        \min_{0\le j\le n-1}\frac{(\eta\tau)^j}{j!}>0.
\end{equation}
At such times, every maximizing coordinate satisfies
$F_i(x)\le H(X)-hXx_{i-q}\le-hRd$.
Starting from $X\le2R$ at $t_0+\tau$, integration proves entry by
$t_0+\tau+1/(hd)$, or earlier. We have proved (iv)$\Rightarrow$(iii).

The implications (iii)$\Rightarrow$(ii)$\Rightarrow$(i) are immediate.
If (iv) fails and $g=h=0$, choose all initial coordinates equal to
$s_0>R$. Direct substitution verifies the solution $x_i(t)=s(t)$
with $\dot s=H(s)$. It stays above $R$ at each finite time:
a first meeting with the constant scalar solution $R$ would violate
uniqueness, also backwards near that meeting. Hence it never enters
$K_R$. The other failure case is $g=0$, $h>0$, $c=0$,
$q\not\equiv0\pmod n$. Here $R=a/k>0$. Choose one initial coordinate
$x_j(0)>R$ and all others zero. The zero coordinates stay zero and
$j-q\ne j$, so $\dot x_j=a x_j-kx_j^2$. It likewise stays strictly
above $a/k$ at every finite time. These cases exhaust failure of (iv)
and contradict (i), completing the equivalence.
\end{proof}

\subsection{Canonical conclusions and weaker-assumption limits}
\label{slot:P01-4}

\begin{corollary}[Canonical foundation]
\label{cor:P01-canonical}
For canonical $n\ge5$, $1\le q\le n-1$ and $a,k,c,g,h>0$, every
nonnegative initial state has a unique global nonnegative solution.
Every initially positive coordinate stays positive, and every nonzero
initial state becomes strictly positive in every coordinate at each
positive time. The cube $[0,X_{\max}]^n$, with
\[
 X_{\max}=\frac{a-k+\sqrt{(a+k)^2+4kc}}{2k}>0,
\]
is positively invariant, globally attracting, and finite-time absorbing
for the entire orthant with the common bound
\eqref{eq:P01-canonical-time}. All larger cubes are also invariant
and absorbing. No condition on $\gcd(n,q)$ is required.
\end{corollary}
\begin{proof}
The canonical assumptions imply \eqref{eq:P01-weak}, $a+c>0$,
$c>0$ and $g>0$. Apply the preceding results.
\end{proof}

\begin{remark}[Archived radius and the meaning of stability]
The blueprint's candidate radius
\[
 B=\frac{a+\sqrt{a^2+4kc}}{2k}
\]
is the larger real root of $as-ks^2+c=0$, and is nonnegative. If $c>0$, then $B>0$ and
$H(B)=-c/(1+B)<0$, so $B>R$. Its cube is therefore also invariant
and absorbing by Theorem~\ref{thm:P01-global}, even when $g=h=0$.
The smaller radius $R$ retains $cX/(1+X)$ instead of replacing it
by $c$. If $c=0$, then $B=R$ and the exact-radius criterion applies.
Neither radius is claimed optimal after all $g,h,q$ information is used.

Invariance concerns remaining inside; absorption concerns finite-time
entry; attraction alone permits approaching from outside forever.
The distance estimates also show Lyapunov stability of $K_R$ as a set,
relative to $\Orth$, because distance never exceeds its initial value.
Under the convention that asymptotic stability of a set means Lyapunov
stability plus attraction, $K_R$ is globally asymptotically stable as
a set. This is not convergence to a particular point, nor a claim that
the cube is a minimal global attractor or is invariant in both time
directions.
\end{remark}

\begin{proposition}[Global existence when $k=0$ is allowed]
\label{prop:P01-kzero}
Global nonnegative existence and the positivity conclusions still hold
for $k\ge0$, $a\in\R$, $c,g,h\ge0$, any $n\ge1$, and any residue $q$.
For $k=0$, no bounded absorbing cube is guaranteed over this parameter
class.
\end{proposition}
\begin{proof}
Local existence and nonnegative invariance did not require $k>0$.
Discarding nonpositive losses at a maximizing coordinate gives
\[
 D^+X\le aX+\frac{cX}{1+X}\le(a_++c)X.
\]
The Lipschitz/Dini argument yields $X(t)\le X(0)e^{(a_++c)t}$.
If the maximal time were finite, this would confine the trajectory
to a compact subset of $\Domain$, and the earlier continuation
argument would extend it. Positivity follows from
Proposition~\ref{prop:P01-positive}. For the last assertion take
$a>0$, $k=c=g=h=0$: an initially positive coordinate equals
$x_i(0)e^{at}$ and is unbounded.
\end{proof}


% EDITION: citation-based shortening, 4 October 2026.
% Earlier audits refer to the source edition; this revision has not been independently audited.
\section{P02: Synchrony and synchronous equilibria}
\label{sec:P02}

The standing assumptions are the canonical A1--A3. In particular,
$n\ge5$, $1\le q\le n-1$, $a,k,c,g,h>0$, and the open extension
domain is $\Domain=(-1,\infty)^n$. The nonnegative model state space
is $\Orth$. Only P01's local uniqueness and, for nonnegative initial data, global continuation are needed here.
The calculations below start from \eqref{eq:canonical-dm}.
No assertion about nonsynchronous equilibria, equilibrium stability,
or bifurcation type is made.

\subsection{Invariant synchronous subspace}
\label{slot:P02-1}

\begin{proposition}[Invariance of the synchronous subspace]
\label{prop:P02-invariance}
The set
\[
 \Sync\cap\Domain=\{s\one:s\in(-1,\infty)\}
\]
is invariant for the canonical vector field for the entire maximal
solution interval in $\Domain$. Its scalar restriction is
\begin{equation}
\label{eq:P02-scalar-original}
 \dot s=f(s):=as-ks^2+\frac{cs}{1+s}
                     -\frac{gs}{1+s}-hs^2,
 \qquad s\in(-1,\infty).
\end{equation}
This conclusion needs no coprimality or distinct-offset assumption.
Indeed, the substitution and local invariance remain valid for any
integer $n\ge1$, any residue $q$ modulo $n$, and any finite real
coefficients. The canonical restrictions are retained for the
equilibrium classification below.
\end{proposition}

\begin{proof}
Put $x=s\one$ with $s>-1$. For every $i\in\Zn$, each of
$x_i,x_{i-1},x_{i+q},x_{i-q}$ equals $s$. Direct substitution gives
\[
 F_i(s\one)=as-ks^2+\frac{cs}{1+s}
                        -\frac{gs}{1+s}-hs^2=f(s).
\]
Every component is identical, so $F(s\one)=f(s)\one$.
Coincident offsets retain all the displayed terms; no distinctness
of neighbors is used. Nor is the size or connectedness of a
step-$q$ cycle used.

The scalar field is smooth on $(-1,\infty)$. A scalar solution
$s(t)$ lifts to the vector solution $s(t)\one$. Uniqueness from
Proposition~\ref{prop:P01-local} makes this lift agree with the
vector solution having initial value $s(0)\one$ wherever both exist.
If the scalar maximal interval ended earlier while the vector
solution continued in $\Domain$, the latter would supply a limit
$s_*\one\in\Domain$ at that time. Local scalar continuation at
$s_*>-1$ would extend the lift, a contradiction. Conversely, a longer
scalar interval would extend the vector solution. Thus the maximal
intervals coincide and synchrony persists throughout. The same local
uniqueness argument applies on either side of an initial time.
The scalar domain inherited from $\Domain$ is exactly $(-1,\infty)$.
\end{proof}

\subsection{Scalar synchronous dynamics}
\label{slot:P02-2}
Use the reserved abbreviations and one local auxiliary abbreviation
\begin{equation}
\label{eq:P02-abbreviations}
 A=k+h>0,\qquad \mu=a+c-g,\qquad b=a-A.
\end{equation}
Here $b$ is used only within P02; it is not the archived comparator's
coefficient. In particular $b>-A$ because $a>0$. The exact restriction is
\begin{equation}
\label{eq:P02-scalar}
 f(s)=s\left(a-As+\frac{c-g}{1+s}\right)
     =\frac{s}{1+s}\bigl(\mu+bs-As^2\bigr),\qquad s>-1.
\end{equation}
For $s(0)\ge0$, P01 and Proposition~\ref{prop:P02-invariance}
give a global nonnegative scalar solution. For $-1<s(0)<0$, only
the maximal interval within this open domain is asserted. Even when
$c=g$ cancels the two rational terms along synchrony, this does not
enlarge the canonical domain or admit $s=-1$.

\subsection{The origin and the nonzero equilibrium equation}

\begin{proposition}[Equilibrium equation with domain restrictions]
\label{prop:P02-equation}
The origin $0\one$ is an equilibrium for every canonical parameter
tuple. Every other synchronous equilibrium in $\Domain$ is exactly
$s\one$, where $s\ne0$, $s>-1$, and
\begin{align}
 0=Q(s)&:=As^2-bs-\mu \label{eq:P02-polynomial}\\
       &=(k+h)s^2+(k+h-a)s-(a+c-g).\nonumber
\end{align}
Define the discriminant
\begin{align}
 D&:=b^2+4A\mu \label{eq:P02-discriminant}\\
  &=(a-A)^2+4A(a+c-g)
   =(a+A)^2+4A(c-g).\nonumber
\end{align}
The complete list of real polynomial roots, when $D\ge0$, is
\begin{equation}
\label{eq:P02-roots}
 s_\pm=\frac{b\pm\sqrt D}{2A}
 =\frac{a-(k+h)\pm
 \sqrt{(a+k+h)^2+4(k+h)(c-g)}}{2(k+h)}.
\end{equation}
They are distinct for $D>0$ and coincide for $D=0$. When $D<0$
there are no real roots of $Q$.
\end{proposition}

\begin{proof}
Directly in \eqref{eq:P02-scalar-original}, all five terms vanish
at $s=0$. Thus the origin is always an equilibrium, before any
division is performed. Now take $s\ne0$ with $s>-1$.
Dividing \eqref{eq:P02-scalar} by $s$ and multiplying by
the strictly positive number $1+s$ gives
\[
 (a-As)(1+s)+c-g=\mu+(a-A)s-As^2=0.
\]
Both operations are reversible on this stated set, giving
\eqref{eq:P02-polynomial}. Completing the square yields
\begin{equation}
\label{eq:P02-square}
 Q(s)=A\left(s-\frac b{2A}\right)^2-\frac D{4A}.
\end{equation}
Expansion of $b^2+4A\mu$ proves both discriminant forms and the
quadratic formula gives \eqref{eq:P02-roots}. Finally,
\begin{equation}
\label{eq:P02-minusone}
 Q(-1)=A+b-\mu=g-c.
\end{equation}
Thus $s=-1$ is a polynomial root precisely when $g=c$, but it is
never a point of the canonical domain. Multiplication by $1+s$
does not authorize retaining that root. A root $s=0$ of $Q$, when
present, is the already counted origin rather than another equilibrium.
\end{proof}

\subsection{Complete classification in the canonical parameter regime}
\label{slot:P02-3}

In this subsection, a \emph{mathematical root} means a real root of
$Q$, including a zero or an inadmissible root. A \emph{DM-domain root}
additionally satisfies $s>-1$. A \emph{nonnegative model equilibrium}
has $s\ge0$, and a \emph{positive model equilibrium} has $s>0$.
The ever-present origin is added once to the admissible nonzero roots.
All counts below count distinct values, not algebraic multiplicity.

\begin{theorem}[Complete synchronous-equilibrium classification]
\label{thm:P02-complete}
Under the canonical assumptions, the synchronous equilibrium set is
\begin{equation}
\label{eq:P02-equilibrium-set}
 \{0\one\}\ \cup\
 \{s\one:Q(s)=0,\ s>-1,\ s\ne0\}.
\end{equation}
The following eleven cases are exhaustive. In the table, $r$ counts
real nonzero polynomial roots, $d$ counts those with $s>-1$, $N_0$
counts all nonnegative synchronous equilibria (including the origin),
and $N_+$ counts the strictly positive ones. The total number of
synchronous equilibria in $\Domain$ is $1+d$.

\begin{center}
\small
\renewcommand{\arraystretch}{1.18}
\begin{tabular}{@{}c p{0.28\textwidth} p{0.35\textwidth} c c c c@{}}
\toprule
Case & Parameter conditions & Real roots of $Q$ and location
 & $r$ & $d$ & $N_0$ & $N_+$\\
\midrule
I & $\mu>0,\ g<c$ & $s_-<-1<0<s_+$ & 2&1&2&1\\
II & $\mu>0,\ g=c$ & $s_-=-1,\ s_+=a/A>0$ & 2&1&2&1\\
III & $\mu>0,\ g>c$ & $-1<s_-<0<s_+$ & 2&2&2&1\\
IV & $\mu=0,\ b>0$ & $s_-=0,\ s_+=b/A>0$ & 1&1&2&1\\
V & $\mu=0,\ b=0$ & $s_-=s_+=0$ (double) & 0&0&1&0\\
VI & $\mu=0,\ b<0$ & $-1<s_-=b/A<0=s_+$ & 1&1&1&0\\
VII & $\mu<0,\ D<0$ & No real roots; any sign of $b$ & 0&0&1&0\\
VIII & $\mu<0,\ D>0,\ b>0$ & $0<s_-<s_+$ & 2&2&3&2\\
IX & $\mu<0,\ D=0,\ b>0$ & $s_-=s_+=b/(2A)>0$ & 1&1&2&1\\
X & $\mu<0,\ D>0,\ b<0$ & $-1<s_-<s_+<0$ & 2&2&1&0\\
XI & $\mu<0,\ D=0,\ b<0$ & $-\tfrac12<s_-=s_+=b/(2A)<0$ & 1&1&1&0\\
\bottomrule
\end{tabular}
\end{center}
When $\mu<0$ and $b=0$, necessarily $D=4A\mu<0$, so case VII
applies. Each appearance of $s_\pm$ refers to \eqref{eq:P02-roots}.
The only repeated polynomial roots are in cases V, IX and XI.

As zeros of the original smooth scalar field $f$ on $(-1,\infty)$,
all nonzero equilibria are simple except the double roots in IX and XI.
The origin is simple if $\mu\ne0$, double if $\mu=0$, $b\ne0$,
and triple if $\mu=b=0$. The algebraic continuation of the nonzero
equilibrium curve meets the origin exactly at $\mu=0$: through $s_-$
in case IV, through $s_+$ in case VI, and through both coincident
root labels in case V. No negative root represents a state in $\Orth$.
\end{theorem}

\begin{proof}
Proposition~\ref{prop:P02-equation} proves completeness of
\eqref{eq:P02-equilibrium-set}. The Vieta relations are
\begin{equation}
\label{eq:P02-vieta}
 s_-+s_+=\frac bA,\qquad s_-s_+=-\frac\mu A.
\end{equation}
We first establish the domain test for every real pair of roots.
Because $a>0$ and $A>0$,
\[
 s_+\ge\frac b{2A}=\frac a{2A}-\frac12>-\frac12>-1.
\]
For the smaller root,
\begin{equation}
\label{eq:P02-domain-test}
 s_->-1
 \ \Longleftrightarrow\ \sqrt D<a+A
 \ \Longleftrightarrow\ D<(a+A)^2
 \ \Longleftrightarrow\ g>c.
\end{equation}
The middle equivalence uses $a+A>0$ and $D\ge0$. Equality gives
$s_-=-1$ exactly when $g=c$; the reverse inequality gives $s_-<-1$
exactly when $g<c$. These are polynomial-root statements on the
real-discriminant set, not assertions that such roots are admissible.

If $\mu>0$, then $D>b^2$, so $s_-<0<s_+$ and both roots are
nonzero and simple. The domain test gives I--III. When $g=c$,
$D=(a+A)^2$ and the stated roots $-1,a/A$ follow directly.

If $\mu=0$, then $Q(s)=s(As-b)$. If $b\ne0$, its distinct roots
are $0$ and $b/A$. Since $b/A=a/A-1>-1$, the latter is always
domain-admissible and has the sign of $b$. Their ordering gives
IV and VI. If $b=0$, $Q(s)=As^2$ gives V. The origin is counted
only once in the equilibrium set in all three cases.

If $\mu<0$, real roots exist precisely when $D\ge0$. For $D<0$
we obtain VII. For $D\ge0$, the product in \eqref{eq:P02-vieta}
is strictly positive and $b\ne0$ (otherwise $D<0$). Thus both
roots are nonzero with the sign of their sum, namely the sign of $b$.
They are distinct for $D>0$ and double for $D=0$. Also
$g>a+c>c$, so both are admissible by \eqref{eq:P02-domain-test}.
This gives VIII--XI, including the sharper negative repeated-root
bound from $b>-A$. These arguments yield every count in the table.

On $s>-1$, the identity $f(s)=-sQ(s)/(1+s)$ shows that the
nonzero zeros have the same multiplicity as their polynomial roots.
At the origin, use \eqref{eq:P02-scalar} without dividing:
the numerator is $\mu s+bs^2-As^3$ and $1+s$ does not vanish.
This gives the stated orders one, two and three. Finally $Q(0)=-\mu$,
so the nonzero equation's algebraic continuation contains $s=0$
exactly when $\mu=0$. The formulas at that locus give its root
labels, and the local branch description below confirms the approach
by nonzero roots. The counting and multiplicity assertions are purely
algebraic and carry no stability conclusion.
\end{proof}

\subsection{Positive synchronous equilibria: necessary and sufficient conditions}

\begin{theorem}[Exact positive-equilibrium alternatives]
\label{thm:P02-positive}
In the canonical parameter regime the following four alternatives
partition parameter space:
\begin{enumerate}[label=(\roman*),leftmargin=*]
\item There is no positive synchronous equilibrium if and only if
\begin{equation}
\label{eq:P02-none}
 [\,\mu=0,\ b\le0\,]
 \quad\hbox{or}\quad
 [\,\mu<0,\ (b\le0\ \hbox{or}\ D<0)\,].
\end{equation}
\item There is exactly one positive synchronous equilibrium and it is
simple if and only if $\mu>0$, or $\mu=0$ and $b>0$.
It is $s_+\one$; in the second case $s_+=b/A$.
\item There are exactly two distinct positive synchronous equilibria
if and only if $\mu<0$, $b>0$, and $D>0$. They are
$s_-\one,s_+\one$, both simple.
\item There is exactly one positive synchronous equilibrium and it
is repeated if and only if $\mu<0$, $b>0$, and $D=0$.
It is a double root, $s_-\one=s_+\one=(b/(2A))\one$.
\end{enumerate}
Consequently, exactly one distinct positive equilibrium (without a
simplicity qualification) occurs if and only if (ii) or (iv) holds.
There are never more than two positive synchronous equilibria.
\end{theorem}

\begin{proof}
The cases of Theorem~\ref{thm:P02-complete} prove every implication,
including the necessity and all equalities. One can also reconstruct
the positive-root count directly from
$H(s)=\mu+bs-As^2$ for $s>0$. If $b\le0$, then
$H'(s)=b-2As<0$ on this interval: it has one positive zero exactly
when $\mu>0$. If $b>0$, its maximum on $s\ge0$ is at
$s=b/(2A)>0$ and has value $D/(4A)$. For $\mu<0$, that maximum
is positive, zero, or negative precisely in alternatives (iii),
(iv), or (i). For $\mu=0$, the zero at the endpoint is excluded
from the positive count and the other root is $b/A>0$. For
$\mu>0$, the opposite-sign Vieta product leaves one positive root.
This gives an independent sign-and-vertex justification of the list.
\end{proof}

\subsection{Algebraic branch loci and degeneracies}

\begin{proposition}[Branch-intersection and repeated-root loci]
\label{prop:P02-loci}
For fixed positive $a,k,c,h$, define
\begin{equation}
\label{eq:P02-gf}
 g_f:=a+c+\frac{(a-A)^2}{4A}
      =c+\frac{(a+A)^2}{4A},\qquad
 s_f:=\frac{a-A}{2A}.
\end{equation}
The following statements concern only the scalar equilibrium equation.
\begin{enumerate}[label=(\roman*),leftmargin=*]
\item The \emph{branch-intersection locus} is $\mu=0$, or $g=a+c$.
It is the set where the algebraic continuation of the nonzero
equilibrium equation meets the origin branch.
\item The \emph{repeated-root locus} of $Q$ is $D=0$, or $g=g_f$,
and its root is $s_f$. Indeed,
\begin{equation}
\label{eq:P02-discriminant-factor}
 D=4A(g_f-g),\qquad
 g_f-(a+c)=\frac{(a-A)^2}{4A}\ge0.
\end{equation}
The root is positive for $a>A$, negative for $0<a<A$, and zero
for $a=A$. It always lies in $(-1/2,\infty)$ and hence in the
DM domain. It is a nonzero double equilibrium only when $a\ne A$.
\item The two parameter loci meet exactly at the \emph{algebraic
degeneracy} $a=A$, $g=a+c$. There $Q(s)=As^2$ and
\begin{equation}
\label{eq:P02-degenerate-field}
 f(s)=-\frac{As^3}{1+s}.
\end{equation}
The origin is the sole synchronous equilibrium in the domain at
that parameter tuple; it is a triple zero of $f$ and a double root
of $Q$, not an additional nonzero equilibrium.
\end{enumerate}
\end{proposition}

\begin{proof}
The nonzero equation can be written, without approximation, as
\begin{equation}
\label{eq:P02-branch-graph}
 \mu=As^2-bs,\qquad
 g=a+c+bs-As^2
   =g_f-A(s-s_f)^2.
\end{equation}
This is its algebraic continuation, to be filtered by $s>-1$,
$g>0$, and $s\ne0$ when counting nonzero model equilibria.
At $s=0$ it gives exactly $g=a+c$. Completing the square or using
\eqref{eq:P02-square} gives $D=0$, $g=g_f$, and $s=s_f$ as the
unique repeated-root condition. The difference of the two parameter
values in \eqref{eq:P02-discriminant-factor} vanishes exactly at
$a=A$. Since $a>0$, $s_f=a/(2A)-1/2>-1/2$. At the common
locus, substitution into \eqref{eq:P02-scalar} proves
\eqref{eq:P02-degenerate-field} and all the remaining assertions.
\end{proof}

For additional local algebra at the origin, fix $a,A,c$ and vary $g$
(equivalently $\mu$). If $b\ne0$, the root that tends to zero is
\begin{equation}
\label{eq:P02-local-branch}
 s_0(\mu)=\frac{b-\operatorname{sgn}(b)\sqrt{b^2+4A\mu}}{2A}
        =-\frac{\mu}{b}+\frac{A\mu^2}{b^3}+O(\mu^3).
\end{equation}
This follows by expansion of the explicit square root near $\mu=0$;
$g=a+c-\mu$ remains positive in a sufficiently small neighborhood.
Thus nonzero domain-admissible roots approach the origin from the
neighboring parameter regions. At the intersection itself,
\begin{equation}
\label{eq:P02-quadratic-coefficient}
 f(s)\big|_{\mu=0}=\frac{s^2(b-As)}{1+s}
                     =bs^2-as^3+O(s^4).
\end{equation}
The coefficient of the usual quadratic scalar term is $b=a-A$.
It vanishes when $a=A$; this is precisely the degeneracy above.

More explicitly, for $a=A$ the nonzero polynomial equation is
$As^2=\mu$. For $\mu>0$ its roots are
\begin{equation}
\label{eq:P02-square-root-branches}
 s_\pm=\pm\sqrt{\mu/A}.
\end{equation}
The negative one is admissible exactly when $0<\mu<A$ (equivalently
$c<g<a+c$); it is $-1$ when $\mu=A$, and below $-1$ when $\mu>A$.
At $\mu=0$ both labels give the origin; for $\mu<0$ no real
nonzero roots exist. These are algebraic branch descriptions only.
Any bifurcation theorem requires the separate work assigned to P09
and P10; none is asserted here.

\subsection{Canonical parameter corollary}

\begin{corollary}[Positive branches in the original canonical parameters]
\label{cor:P02-canonical}
Let $a,k,c,g,h>0$, $A=k+h$, and let $g_f$ be as in
\eqref{eq:P02-gf}. If $a\le k+h$, there is one simple positive
synchronous equilibrium exactly for $0<g<a+c$, and none for
$g\ge a+c$. If $a>k+h$, the exact alternatives are
\begin{center}
\begin{tabular}{@{}ll@{}}
\toprule
Parameter range & Positive synchronous equilibria\\
\midrule
$0<g\le a+c$ & One simple equilibrium, $s_+\one$\\
$a+c<g<g_f$ & Two distinct simple equilibria, $s_-\one,s_+\one$\\
$g=g_f$ & One double equilibrium, $s_f\one$\\
$g>g_f$ & None\\
\bottomrule
\end{tabular}
\end{center}
In particular, $a+c<g$ does not by itself preclude positive
synchrony: two distinct positive equilibria exist if and only if
\begin{equation}
\label{eq:P02-two-positive-original}
 a>k+h,\qquad
 a+c<g<a+c+\frac{(a-k-h)^2}{4(k+h)}.
\end{equation}
The upper equality gives one repeated positive equilibrium when
$a>k+h$; the lower equality gives one simple positive equilibrium
and the origin. The origin exists in every row, independently of
the displayed positive count. There is no dependence on $n$, $q$,
$\gcd(n,q)$, or distinctness of the offsets.
\end{corollary}

\begin{proof}
The sign of $b$ is the sign of $a-(k+h)$; the sign of $\mu$ is
the sign of $a+c-g$; and $D=4A(g_f-g)$ with $A>0$.
Substitute these identities into Theorem~\ref{thm:P02-positive}.
\end{proof}

\subsection{Weaker assumptions and limits of the classification}
The synchronous substitution and the scalar identity require only
homogeneity, finite real coefficients, and nonvanishing denominators.
They therefore apply under P01's weaker assumptions as well.
For $A>0$, the polynomial formula, discriminant, Vieta relations,
and the positive-root alternatives of Theorem~\ref{thm:P02-positive}
remain valid for real $a,c,g$: their positive-root proofs use only
$A>0$. This observation is not a generalization of the canonical
negative-root table. In particular, the domain test
\eqref{eq:P02-domain-test} uses $a+A>0$, and the bound $s_f>-1/2$
uses $a>0$. P01 allows real $a$ and does not supply these inequalities.
The complete domain-filtered classification in this module is
intentionally restricted to the canonical positive parameter regime.
No classification for $A=0$ or $A<0$ is claimed.

\subsection{Checks and possible failure points}
\label{subsec:P02-checks}

\paragraph{Root substitution and discriminant reconstruction.}
For $D\ge0$, put $d=\sqrt D$ and
$r=(b+\varepsilon d)/(2A)$, $\varepsilon\in\{-1,1\}$.
Direct expansion gives
\[
 \mu+br-Ar^2
 =\mu+\frac{2b(b+\varepsilon d)-(b+\varepsilon d)^2}{4A}
 =\mu+\frac{b^2-d^2}{4A}=0.
\]
For $r\ne-1$, substitution in the original scalar expression gives
\[
 ar-kr^2+\frac{cr}{1+r}-\frac{gr}{1+r}-hr^2
   =\frac{r}{1+r}(\mu+br-Ar^2)=0.
\]
Only $r>-1$ is retained as a DM equilibrium. If $r=-1$, the two
original rational terms are undefined even when $c=g$, so no
substitution or canceled extension admits that point.
At $\mu=0$, the specialized root $r=b/A$ also gives
$br-Ar^2=0$. At $D=0$, $r=s_f=b/(2A)$ gives
$\mu+br-Ar^2=\mu+b^2/(4A)=0$. At $g=c$, the retained root
$a/A$ gives $ar-Ar^2=0$. At $b=0$, the square-root branches give
$\mu-Ar^2=0$. Finally, direct substitution at $r=0$ gives zero
for every parameter tuple. These checks include every exact root
formula displayed above, with $s_0$ in \eqref{eq:P02-local-branch}
being one of $s_\pm$; its truncated expansion is not used as a root.

Independently expanding the polynomial discriminant gives
\[
 (A-a)^2-4A[-(a+c-g)]
 =a^2+2aA+A^2+4A(c-g)
 =(a+A)^2+4A(c-g).
\]
Also the sum of the two explicit roots is $b/A$, and their product
is $(b^2-D)/(4A^2)=-\mu/A$, confirming both Vieta relations.

\paragraph{Exact representative cells and equality boundaries.}
The following rational tuples realize every row of the classification.
Set $k=h=1/2$, hence $A=1$, and $c=1$ throughout.
All original coefficients are strictly positive. Roots are obtained
by the formulas already proved; the counts refer to the preceding table.
\begin{center}
\small
\renewcommand{\arraystretch}{1.12}
\begin{tabular}{@{}c c c c c c@{}}
\toprule
Case & $a$ & $g$ & $\mu$ & $D$ & $(r,d,N_0,N_+)$\\
\midrule
I & $2$ & $1/2$ & $5/2$ & $11$ & $(2,1,2,1)$\\
II & $2$ & $1$ & $2$ & $9$ & $(2,1,2,1)$\\
III & $2$ & $2$ & $1$ & $5$ & $(2,2,2,1)$\\
IV & $2$ & $3$ & $0$ & $1$ & $(1,1,2,1)$\\
V & $1$ & $2$ & $0$ & $0$ & $(0,0,1,0)$\\
VI & $1/2$ & $3/2$ & $0$ & $1/4$ & $(1,1,1,0)$\\
VII & $1$ & $3$ & $-1$ & $-4$ & $(0,0,1,0)$\\
VIII & $2$ & $25/8$ & $-1/8$ & $1/2$ & $(2,2,3,2)$\\
IX & $2$ & $13/4$ & $-1/4$ & $0$ & $(1,1,2,1)$\\
X & $1/2$ & $49/32$ & $-1/32$ & $1/8$ & $(2,2,1,0)$\\
XI & $1/2$ & $25/16$ & $-1/16$ & $0$ & $(1,1,1,0)$\\
\bottomrule
\end{tabular}
\end{center}
Cases II, IV--VI, IX and XI test all relevant equality surfaces,
including the simultaneous equality $a=A$, $g=a+c$.
For $b<0,0,>0$ inside the $\mu>0$ cases, one may respectively
take $a=1/2,1,2$, and $g=1/2,1,1+a/2$ for I, II, III.
For all three signs of $b$ in VII, take the same three $a$ values
and $g=g_f+1$. These choices also check that omitted sign subcases
do not alter the counts. Internal exact-arithmetic checks supplement
these displayed examples; they are not numerical dependencies or
independent verification of the module.

For reproducibility of that finite check, take
$V=\{1/4,1/2,1,2,3,4\}$ and all $(a,A,c,g)\in V^4$.
For each $(a,A,c)\in V^3$, also take
$g\in\{c,a+c,g_f\}+\{-\eta,0,\eta\}$, where
$\eta=\min\{a,A,c,1\}/64$. The union contains 2,868 distinct
positive tuples. Set $k=A/3$, $h=2A/3$ for this grid; use the
displayed $k=h=1/2$ for the eleven table examples. Exact arithmetic
in $u+v\sqrt D$, with rational $u,v,D$, checks root residuals,
Vieta relations, domain comparisons, and every positive-count
equivalence and its converse. Including the eleven table examples,
4,438 original scalar residuals vanish exactly and 217 occurrences
of $s=-1$ are explicitly excluded. All eleven cells and all equality
surfaces are represented. These finite checks found no counterexample;
they supplement rather than establish the exhaustive proof.

\paragraph{Boundary continuity.}
The ordered formulas $s_\pm$ are continuous on the parameter set
$D\ge0$, since $A>0$. At $D\downarrow0$ they both tend to $s_f$;
for $D<0$ their absence as real roots is recorded, not bridged by
a real continuation. At $\mu\to0$ with fixed $b>0$,
$(s_-,s_+)\to(0,b/A)$; with fixed $b<0$,
$(s_-,s_+)\to(b/A,0)$. With $b=0$ and $\mu\downarrow0$, both labels
tend to zero as in \eqref{eq:P02-square-root-branches}. At $g\to c$,
$(s_-,s_+)\to(-1,a/A)$; the first root is excluded from the DM domain
at the boundary.
At $b\to0$ all limits within $D\ge0$ are again given by the same
continuous formulas. The sign-specific table is therefore compatible
with every boundary, although admissible root counts can change there.

\clearpage
\paragraph{Logical completeness and attempted counterexamples.}
The origin was retained before division, and every other equilibrium
is equivalent to a nonzero root of $Q$ in $s>-1$. This proves that
none was lost or introduced. In particular, the excluded root $-1$
at $g=c$ and the zero root at $\mu=0$ cannot be counted as additional
nonzero equilibria. Positivity, discriminant sign, and root order
partition all parameter tuples as in the theorem; each implication
in the necessary-and-sufficient statements is checked in every cell.
The independent positive-half-line argument supplies a second check.
Case VIII is an explicit canonical counterexample to the rejected
claim that $a+c\le g$ rules out positive synchrony. Case IV is a
counterexample even at equality. Case X shows that two admissible
nonzero roots need not be positive; case XI shows that a repeated
domain root need not be positive. Case V prevents counting a repeated
zero as a positive equilibrium, and case II prevents admitting a
canceled denominator root. All these distinctions are essential.


% EDITION: citation-based shortening, 4 October 2026.
% Earlier audits refer to the source edition; this revision has not been independently audited.
\section{P03: Origin and homogeneous linearization}
\label{sec:P03}
Use the canonical demand-modulated field \eqref{eq:canonical-dm},
with $n\ge5$, $1\le q\le n-1$, $a,k,c,g,h>0$, and
$\Domain=(-1,\infty)^n$. All indices and equalities of indices are
modulo $n$. Write $\mu=a+c-g$ and $A=k+h$ as in the project notation.
P01 supplies smoothness on $\Domain$; P02 supplies the synchronous
restriction and, only where explicitly invoked below, its equilibrium
relation. No formula is imported from the archived blueprint.
The stability equivalences below concern the \emph{linearized system}.
Nonlinear behavior at a zero spectral abscissa and bifurcation type
are reserved for P09. The positive-state Fourier spectrum and the
homogeneous normality theorem are reserved for P04.

\subsection{General Jacobian}

Let $\delta^{(n)}_{j\ell}$ be one if $j\equiv\ell\pmod n$ and
zero otherwise. These are modular Kronecker symbols, not mutually
exclusive case labels.

\begin{lemma}[Jacobian entries]
\label{lem:P03-jacobian}
For every $x\in\Domain$ and $i,j\in\Zn$, $J(x)=DF(x)$ has entries
\begin{align}
 J_{ij}(x)={}&\left(a-2kx_i-\frac{c x_{i-1}}{(1+x_i)^2}
             -\frac{g}{1+x_{i+q}}-h x_{i-q}\right)
                 \delta^{(n)}_{ji}\notag\\
 &+\frac{c}{1+x_i}\delta^{(n)}_{j,i-1}
   +\frac{g x_i}{(1+x_{i+q})^2}\delta^{(n)}_{j,i+q}
   -h x_i\delta^{(n)}_{j,i-q}.
 \label{eq:P03-general-entries}
\end{align}
All contributions are added when indices coincide; all other entries
are zero. The three off-diagonal indices differ from $i$, but need
not differ from each other.
\end{lemma}
\begin{proof}
Differentiate the five terms of \eqref{eq:canonical-dm} by the product
and chain rules. In particular the denominator derivative of
$-gx_i/(1+x_{i+q})$ is $+gx_i/(1+x_{i+q})^2$.
When indices coincide, differentiation adds every contribution,
as encoded by the modular Kronecker symbols.
\end{proof}

For a second formulation, put
\[
 D_i=a-2kx_i-\frac{c x_{i-1}}{(1+x_i)^2}
            -\frac{g}{1+x_{i+q}}-h x_{i-q},\quad
 C_i=\frac{c}{1+x_i},\quad
 G_i=\frac{g x_i}{(1+x_{i+q})^2},\quad H_i=h x_i.
\]
The directional derivative is
\begin{equation}
 (DF(x)z)_i=D_i z_i+C_i z_{i-1}+G_i z_{i+q}-H_i z_{i-q}.
 \label{eq:P03-directional}
\end{equation}
Since $(P^r z)_i=z_{i-r}$ for every integer $r$, this gives
\begin{equation}
 J(x)=\diag(D_i)+\diag(C_i)P+\diag(G_i)P^{-q}
                           -\diag(H_i)P^q.
 \label{eq:P03-variable-matrix}
\end{equation}
The diagonal factors are on the left because their coefficients
belong to receiving row $i$. This general matrix need not be circulant.

\subsection{Linearization at the origin}
\label{slot:P03-1}

Direct substitution gives $F(0)=0$. In
\eqref{eq:P03-general-entries}, $D_i(0)=a-g$, $C_i(0)=c$, and
$G_i(0)=H_i(0)=0$. Both derivatives of the bilinear $h$ term
contain a state coordinate and vanish. The denominator-coordinate
(target) derivative of the input contains $x_{i-1}$ and vanishes;
that of the rational loss contains $x_i$ and vanishes as well.
The numerator derivative of the loss survives as $-g$ on the
 diagonal. Thus, including every collision case,
\begin{equation}
 \boxed{J(0)=(a-g)I+cP.}
 \label{eq:P03-origin-jacobian}
\end{equation}

\begin{proposition}[Origin spectrum]
\label{prop:P03-origin-spectrum}
With the project's DFT convention $\theta_m=2\pi m/n$ and
$v^{(m)}_j=n^{-1/2}e^{\ii\theta_m j}$, the complete spectrum is
\begin{equation}
 \boxed{\lambda_m(0)=a-g+c e^{-\ii\theta_m},\qquad 0\le m<n.}
 \label{eq:P03-origin-spectrum}
\end{equation}
Every eigenvalue is algebraically and geometrically simple over $\C$.
Its real and imaginary parts are
\begin{equation}
 \Ree\lambda_m(0)=\mu-c(1-\cos\theta_m),\qquad
 \Imm\lambda_m(0)=-c\sin\theta_m.
 \label{eq:P03-origin-parts}
\end{equation}
The unique mode attaining the largest real part is $m=0$, and
\begin{equation}
 \boxed{\sa(J(0))=\lambda_0(0)=\mu=a+c-g.}
 \label{eq:P03-abscissa}
\end{equation}
\end{proposition}
\begin{proof}
Apply Gray's circulant diagonalization theorem
\cite[Theorem~3.1]{Gray2006} to \eqref{eq:P03-origin-jacobian}.
With our positive-exponent columns, $Pv^{(m)}=e^{-\ii\theta_m}v^{(m)}$;
Gray's opposite exponent is reconciled by relabelling $m\mapsto -m$.
Distinct roots of unity and $c>0$ give simple eigenvalues. In particular
\begin{equation}
 \det(\lambda I-J(0))=(\lambda-(a-g))^n-c^n.
 \label{eq:P03-characteristic}
\end{equation}
Taking real and imaginary parts gives the stated formulas. Since
$1-\cos\theta_m>0$ for every $m\ne0$, only $m=0$ attains $\mu$.
\end{proof}

\subsection{Origin stability and the nonhyperbolic boundary}
\label{slot:P03-2}

\begin{proposition}[Exact linear stability and hyperbolicity conditions]
\label{prop:P03-stability}
For $\dot z=J(0)z$, zero is asymptotically (indeed exponentially)
stable if and only if $\mu<0$, and unstable if and only if $\mu>0$.
At $\mu=0$ it is Lyapunov stable but not asymptotically stable.
The origin of the nonlinear field is nonhyperbolic if and only if
\begin{equation}
 g=a+c\cos\frac{2\pi m}{n}
 \quad\hbox{for some }m\in\{0,\ldots,n-1\},
 \label{eq:P03-all-nonhyperbolic}
\end{equation}
with $g>0$ as always. Nonhyperbolicity is therefore not restricted
to the stability boundary $\mu=0$: every locus in
\eqref{eq:P03-all-nonhyperbolic} with $m\ne0$ has $\mu>0$.
\end{proposition}
\begin{proof}
In the orthonormal basis of Proposition~\ref{prop:P03-origin-spectrum},
each coefficient solves $\dot{\widehat z}_m=\lambda_m(0)\widehat z_m$.
For $\mu<0$, all coefficients decay exponentially, with
$\norm{z(t)}_2\le e^{\mu t}\norm{z(0)}_2$ for $t\ge0$.
For $\mu>0$, the real synchronous solution
$z(t)=\varepsilon e^{\mu t}\one$ proves instability for arbitrarily
small $\varepsilon>0$. For $\mu=0$, every nonsynchronous coefficient
 decays and the synchronous coefficient is constant. Therefore
$\norm{z(t)}_2\le\norm{z(0)}_2$, while nonzero constant synchronous
solutions prevent attraction. This also proves both converses
without a nonlinear stability theorem. The definition of hyperbolicity
and \eqref{eq:P03-origin-parts} give
\eqref{eq:P03-all-nonhyperbolic}. For $m\ne0$ its value of
$\mu$ is $c(1-\cos\theta_m)>0$.
\end{proof}

The spectral-abscissa sign and the linear stability threshold depend
on neither $n$, $q$, nor $\gcd(n,q)$. Individual eigenvalues and the
\emph{full set} of nonhyperbolic loci depend on $n$, but not on $q$
or a coprimality condition. For example, $n=6$, $a=2$, $c=1$,
$g=5/2$ gives $\mu=1/2>0$ and purely imaginary eigenvalues in
modes $1,5$; mode zero is already unstable. This is not a Hopf claim.
More generally, using $0\le m\le\lfloor n/2\rfloor$ as pair
representatives, each interior representative in
\eqref{eq:P03-all-nonhyperbolic} contributes exactly the two
simple nonzero imaginary eigenvalues in modes $m,n-m$.
Strict decrease of the cosine on $[0,\pi]$ excludes any other
simultaneous zero-real-part mode. The endpoint cases are treated below.

\begin{proposition}[Simple zero at the stability boundary]
\label{prop:P03-boundary}
At $g=a+c$ the only zero eigenvalue is $\lambda_0(0)=0$.
Its algebraic and geometric multiplicities are both one, with
kernel $\operatorname{span}\{\one\}$. Every nonzero mode has
strictly negative real part; none can simultaneously have zero real part.
\end{proposition}
\begin{proof}
At $g=a+c$, \eqref{eq:P03-origin-spectrum} becomes
$c(e^{-\ii\theta_m}-1)$. Only $m=0$ is zero, and every other real
part is negative. Simplicity follows from Proposition~\ref{prop:P03-origin-spectrum};
$Pz=z$ is equivalent to $z\in\operatorname{span}\{\one\}$.
\end{proof}

\begin{remark}[A zero eigenvalue away from the stability boundary]
\label{rem:P03-other-zero}
Over \emph{all} canonical parameter values, a nonzero-mode zero
is possible. Its imaginary part requires $\sin\theta_m=0$;
the only nonzero possibility is even $n$ and $m=n/2$.
The real part then vanishes exactly at $g=a-c$, which is
canonical precisely when $a>c$. That zero is also simple, but
$\mu=2c>0$, so the origin is already linearly unstable. For example,
$a=2$, $c=g=1$ works in every even canonical dimension.
Thus $c>0$ excludes nonzero-mode zeros at $g=a+c$, not at every
parameter value. For odd $n$ there is no nonzero-mode zero.
\end{remark}

These are spectral and linear-system conclusions. ``Nonhyperbolic''
and ``linearly unstable'' are not disjoint categories. The Hurwitz
criterion is not an assertion that nonlinear asymptotic stability
holds if and only if $\mu<0$. Linearization alone does not decide
nonlinear behavior at $\mu=0$, either in $\Domain$ or relative to
$\Orth$. The simple zero does not establish a transcritical
bifurcation or a nonlinear nondegeneracy condition.

\subsection{Jacobian at a synchronous state}
\label{slot:P03-3}
Let $x=x^*\one$ be any synchronous state with $x^*>-1$ and set
$s=1+x^*>0$. Here $s$ denotes the positive denominator, not the
scalar state variable denoted by $s$ in P02.

\begin{proposition}[Homogeneous Jacobian]
\label{prop:P03-homogeneous}
At every such state, whether or not it is an equilibrium, define
\begin{align}
 d&=a-2kx^*-\frac{c x^*}{s^2}-\frac{g}{s}-h x^*,
       \label{eq:P03-d}\\
 C&=\frac{c}{s},\qquad G=\frac{g x^*}{s^2},\qquad H=h x^*.
       \label{eq:P03-CGH}
\end{align}
Then
\begin{equation}
 \boxed{J(x^*\one)=dI+CP+GP^{-q}-HP^q.}
 \label{eq:P03-homogeneous}
\end{equation}
Its entries are
$d\delta^{(n)}_{ji}+C\delta^{(n)}_{j,i-1}
+G\delta^{(n)}_{j,i+q}-H\delta^{(n)}_{j,i-q}$,
with sums retained at coincident indices.
\end{proposition}
\begin{proof}
Substitute $x_i=x^*$ in \eqref{eq:P03-general-entries}.
The index $i+q$ corresponds to $P^{-q}$ and $i-q$ to $P^q$;
coincident terms still add.
\end{proof}

Always $C>0$. If $x^*>0$, then $G,H>0$, and the separate
neighbor contributions have signs $+C,+G,-H$. If $x^*=0$,
then $G=H=0$; if $-1<x^*<0$, then $G,H<0$.
No fixed sign for $d$ at an arbitrary synchronous state is assumed.
Signs at colliding offsets must be assessed after addition.

\paragraph{Optional equilibrium simplification using P02.}
Only now suppose $x^*\ne0$ and $x^*\one$ is an equilibrium.
P02's \eqref{eq:P02-scalar}, divided by this nonzero state, gives
\begin{equation}
 a-(k+h)x^*+\frac{c-g}{1+x^*}=0.
 \label{eq:P03-equilibrium-relation}
\end{equation}
Substitution of $a=(k+h)x^*-(c-g)/s$ in the unsimplified
 diagonal coefficient yields
\begin{equation}
 d=-k x^*-\frac{c}{s}-\frac{c x^*}{s^2}
   =-k x^*-\frac{c(1+2x^*)}{s^2}.
 \label{eq:P03-equilibrium-d}
\end{equation}
This applies in particular at a positive synchronous equilibrium.
It is not available at an arbitrary state or at the origin merely
because the origin is an equilibrium. No stability conclusion is drawn.

\subsection{Circulant representation and offset collisions}

Since $P^n=I$, \eqref{eq:P03-homogeneous} is the ordinary polynomial
$dI+CP+GP^{n-q}-HP^q$. Each power of $P$ has entries depending
only on the cyclic difference $i-j$; their linear combination has
rows obtained by cyclic shifts of the preceding row. This proves
that the homogeneous Jacobian is circulant without requiring four
distinct cyclic diagonals. The standard consequence is Fourier
 diagonalizability; the actual homogeneous spectrum and normality
 theorem are left to P04.

The diagonal cannot collide with a neighbor since $1\not\equiv0$
and $q\not\equiv0$ modulo $n$. Neighbor collisions are exhausted by
\[
\begin{aligned}
 i-1\equiv i-q&\iff q=1,\\
 i-1\equiv i+q&\iff q=n-1,\\
 i+q\equiv i-q&\iff 2q\equiv0\pmod n.
\end{aligned}
\]
The last case means even $n$ and $q=n/2$. For $n\ge5$, these
three cases are mutually exclusive; no triple collision occurs.
A nontrivial $\gcd(n,q)$ creates no further coincident offsets.
The resulting matrices are
\begin{align}
 q=1:&\quad J=dI+(C-H)P+GP^{-1},
    \label{eq:P03-collision-one}\\
 q=n-1:&\quad J=dI+(C+G)P-HP^{-1},
    \label{eq:P03-collision-last}\\
 n\text{ even},\ q=n/2:&\quad J=dI+CP+(G-H)P^{n/2}.
    \label{eq:P03-collision-half}
\end{align}
For other canonical $q$ the three neighbor offsets are distinct.
Even then a coefficient may vanish; the representation lists
contributions, not a claim that every entry is nonzero.

At nonsynchronous states the same rule gives, for $q=1$, the
entry $c/(1+x_i)-h x_i$ at $j=i-1$; for $q=n-1$, the entry
$c/(1+x_i)+g x_i/(1+x_{i-1})^2$ at $j=i-1$; and for even $n$,
$q=n/2$, the entry $g x_i/(1+x_{i+q})^2-h x_i$ at
$j=i+q=i-q$. The actual repeated coordinate is identified before
combining contributions. At homogeneous states $C-H$ or $G-H$
can vanish through cancellation. The summed polynomial remains valid.

Differentiating P02's synchronous field also gives the consistency identity
\begin{equation}
 f'(r)=a-2(k+h)r+\frac{c-g}{(1+r)^2}
       =d+C+G-H\quad(x^*=r).
 \label{eq:P03-row-sum-check}
\end{equation}


% EDITION: citation-based shortening, 4 October 2026.
% Earlier audits refer to the source edition; this revision has not been independently audited.
\section{P04: Arbitrary-\texorpdfstring{$(n,q)$}{(n,q)} Fourier spectrum and normality}
\label{sec:P04}
The standing assumptions are $n\ge5$, $1\le q<n$ and
$a,k,c,g,h>0$. Use the canonical domain $\Domain=(-1,\infty)^n$,
indices in $\Zn$, predecessor shift $(Pz)_j=z_{j-1}$ and the
unitarily normalized, positive-exponent Fourier vectors from
Section~\ref{sec:notation}. No coprimality condition is imposed.
This is supporting, known Fourier machinery specialized to the
verified homogeneous Jacobian, not a novelty claim.

For any synchronous state $x^*\one\in\Domain$, set $s=1+x^*>0$.
Here $s$ is the denominator, whereas the scalar state variable below
is $r$. Import Proposition~\ref{prop:P03-homogeneous}, without
re-deriving its Jacobian:
\begin{align}
 J&:=DF(x^*\one)=dI+CP+GP^{-q}-HP^q,\label{eq:P04-input}\\
 d&=a-2kx^*-\frac{cx^*}{s^2}-\frac{g}{s}-hx^*,
 &C&=\frac{c}{s},\quad G=\frac{gx^*}{s^2},\quad H=hx^*.
 \label{eq:P04-coefficients}
\end{align}
All coefficients are real; $C>0$, while $G,H$ have the sign of
$x^*$ (and vanish at zero). The spectral statements through the
synchronous-mode identity apply even when the state is not an
equilibrium. At a nonstationary state they describe its Jacobian,
not an autonomous linearization about an equilibrium.

\subsection{Fourier diagonalization}
\label{slot:P04-1}

\begin{theorem}[Complete homogeneous Fourier spectrum]
\label{thm:P04-spectrum}
Let $\theta_m=2\pi m/n$, $v^{(m)}_j=n^{-1/2}e^{\ii\theta_m j}$,
and $Q:=\Fourier=[v^{(0)}\ \cdots\ v^{(n-1)}]$.
Then $Q^*Q=QQ^*=I$, and the full spectrum, listed with algebraic
multiplicity, is
\begin{equation}
 \boxed{\lambda_m=d+Ce^{-\ii\theta_m}
                      +Ge^{\ii q\theta_m}-He^{-\ii q\theta_m},
           \qquad m=0,\ldots,n-1.}
 \label{eq:P04-spectrum}
\end{equation}
In particular,
\begin{equation}
 Jv^{(m)}=\lambda_m v^{(m)},\qquad
 Q^*JQ=\Lambda:=\diag(\lambda_0,\ldots,\lambda_{n-1}),\qquad
 J=Q\Lambda Q^*.
 \label{eq:P04-diagonalization}
\end{equation}
These statements do not require distinct shift offsets, distinct
numerical eigenvalues, nonzero combined coefficients, or
$\gcd(n,q)=1$.
\end{theorem}
\begin{proof}
Equation~\eqref{eq:P04-input} is a polynomial in the cyclic shift,
so Gray's theorem \cite[Theorem~3.1]{Gray2006} gives a unitary Fourier
diagonalization, also when offsets or numerical eigenvalues coincide.
The project's convention fixes the phases as
\begin{equation}
 Pv^{(m)}=e^{-\ii\theta_m}v^{(m)},\quad
 P^{-q}v^{(m)}=e^{\ii q\theta_m}v^{(m)},\quad
 P^qv^{(m)}=e^{-\ii q\theta_m}v^{(m)}.
 \label{eq:P04-shifts}
\end{equation}
Substitution gives \eqref{eq:P04-spectrum} and
\eqref{eq:P04-diagonalization}, hence
\begin{equation}
 \det(\zeta I-J)=\prod_{m=0}^{n-1}(\zeta-\lambda_m).
 \label{eq:P04-characteristic}
\end{equation}
We reuse this diagonalization throughout.
\end{proof}

\subsection{Real and imaginary dispersion relations}

Writing $\lambda_m=r_m+\ii\omega_m$, Euler's identities give
\begin{align}
 \boxed{r_m=d+C\cos\theta_m+(G-H)\cos(q\theta_m),}
   \label{eq:P04-real}\\
 \boxed{\omega_m=-C\sin\theta_m+(G+H)\sin(q\theta_m).}
   \label{eq:P04-imag}
\end{align}
Indeed, the imaginary contributions of the three shift terms are
$-C\sin\theta_m$, $+G\sin(q\theta_m)$ and
$+H\sin(q\theta_m)$: the last plus sign is the product of
$-H$ and the negative imaginary part of $e^{-\ii q\theta_m}$.
Their real contributions are $C\cos\theta_m$,
$G\cos(q\theta_m)$ and $-H\cos(q\theta_m)$.
Here $\omega_m$ is a signed imaginary part, possibly zero; it is not
the reserved positive frequency $\omega$ of a future Hopf theorem.

\subsection{Conjugacy and multiplicities}

For a label $m$, write $\bar m=(-m)\bmod n$ in $\{0,\ldots,n-1\}$;
in particular $\bar0=0$, not an additional label $n$.
The phases at $\bar m$ are conjugates of those at $m$, so
\begin{equation}
 v^{(\bar m)}=\overline{v^{(m)}},\qquad
 \lambda_{\bar m}=\overline{\lambda_m},\qquad
 r_{\bar m}=r_m,\quad \omega_{\bar m}=-\omega_m.
 \label{eq:P04-conjugacy}
\end{equation}
This verifies conjugacy explicitly from the formula and real
coefficients, consistently with $J$ being real.

A label is self-conjugate precisely when $2m=0\pmod n$.
For odd $n$ there is the singleton $0$ and the $(n-1)/2$ pairs
$\{m,n-m\}$ with $1\le m\le(n-1)/2$.
For even $n$ the singletons are $0,n/2$, with $n/2-1$ pairs
$1\le m<n/2$. Each singleton necessarily has a real eigenvalue.
The converse is false: a non-singleton pair can have $\omega_m=0$,
and then its two labels give the same real eigenvalue.

\begin{proposition}[Numerical eigenvalue multiplicities]
\label{prop:P04-multiplicity}
For a numerical eigenvalue $\eta$, define its label set
$M_\eta=\{m:0\le m<n,\ \lambda_m=\eta\}$.
Its algebraic and geometric multiplicities over $\C$ both equal
$|M_\eta|$, and
\[
 \ker(J-\eta I)=\operatorname{span}_{\C}
                  \{v^{(m)}:m\in M_\eta\}.
\]
Every eigenvalue is semisimple. Each Fourier label contributes one
one-dimensional Fourier coordinate, but this does not imply that
its numerical eigenvalue is simple.
\end{proposition}
\begin{proof}
Equation~\eqref{eq:P04-characteristic} counts the algebraic
multiplicity. Write $z=Qw$; then $(J-\eta I)z=0$ if and only if
$(\lambda_m-\eta)w_m=0$ for every $m$. Exactly the coordinates in
$M_\eta$ are free, proving the eigenspace formula and dimension.
Diagonalization also excludes nontrivial Jordan blocks.
\end{proof}

Thus ``Fourier-mode multiplicity'' means the number of labels
producing a given value, if that term is used at all; it is not a
second source of multiplicity. A nonreal conjugate pair consists
of two different numerical eigenvalues, each of which can also
coincide with eigenvalues at other labels. A tie $r_m=r_\ell$
alone says nothing about equality of $\lambda_m,\lambda_\ell$.
Singleton labels can likewise share their numerical value with
other labels; singleton refers to label conjugacy, not simplicity.

For a concrete canonical example, take $q=1$, $x^*=1$ and
$(a,k,c,g,h)=(2,1,4,4,1)$. P02's scalar equation vanishes,
and P03 gives $d=-4,C=2,G=H=1$. Hence
$J=-4I+P+P^{-1}$ and $\lambda_m=-4+2\cos\theta_m$;
each non-singleton pair has a repeated real eigenvalue.
At the origin, by contrast, $G=H=0$, $C=c>0$ and $d=a-g$,
so \eqref{eq:P04-spectrum} reduces to P03's
$a-g+ce^{-\ii\theta_m}$. The distinct roots of unity and $c>0$
make all those origin eigenvalues simple, as established in
Proposition~\ref{prop:P03-origin-spectrum}.

\subsection{Offset collisions and self-conjugate modes}
\label{slot:P04-extra-1}

P03's exhaustive collision list requires no replacement spectrum.
Direct substitution in the master formula gives the following
checks, with phases understood modulo $2\pi$.
\begin{enumerate}[label=(\roman*),leftmargin=*]
\item If $q=1$, then
\[
 J=dI+(C-H)P+GP^{-1},\qquad
 \lambda_m=d+(C-H)e^{-\ii\theta_m}+Ge^{\ii\theta_m}.
\]
The same expression follows by combining the $C$ and $-H$ terms
in \eqref{eq:P04-spectrum}. Here
$r_m=d+(C+G-H)\cos\theta_m$ and
$\omega_m=(-C+G+H)\sin\theta_m$.
\item If $q=n-1$, then $P^{-q}=P$ and $P^q=P^{-1}$.
Since $e^{\ii(n-1)\theta_m}=e^{-\ii\theta_m}$ and
$e^{-\ii(n-1)\theta_m}=e^{\ii\theta_m}$,
\[
 J=dI+(C+G)P-HP^{-1},\qquad
 \lambda_m=d+(C+G)e^{-\ii\theta_m}-He^{\ii\theta_m}.
\]
Thus $r_m=d+(C+G-H)\cos\theta_m$ and
$\omega_m=-(C+G+H)\sin\theta_m$, also obtained from
\eqref{eq:P04-real}--\eqref{eq:P04-imag}.
\item If $n$ is even and $q=n/2$, then
$e^{\pm\ii q\theta_m}=e^{\pm\ii\pi m}=(-1)^m$ and
$P^{-q}=P^q=P^{n/2}$. Therefore
\[
 J=dI+CP+(G-H)P^{n/2},\qquad
 \lambda_m=d+Ce^{-\ii\theta_m}+(G-H)(-1)^m.
\]
Here $r_m=d+C\cos\theta_m+(G-H)(-1)^m$ and
$\omega_m=-C\sin\theta_m$.
\end{enumerate}
These equalities include cancellations such as $C=H$ in the first
case or $G=H$ in the third. Adding contributions before or after
Fourier transformation gives identical results; coincident entries
must never be overwritten.

The singleton formulas, for every admissible $q$, are
\begin{align}
 \lambda_0&=d+C+G-H\in\R,\label{eq:P04-zero-mode}\\
 \lambda_{n/2}&=d-C+(G-H)(-1)^q\in\R
       \quad(n\text{ even}).\label{eq:P04-half-mode}
\end{align}
For the second formula use $e^{-\ii\pi}=-1$ and
$e^{\pm\ii q\pi}=(-1)^q$; all sine terms vanish.
It equals $d-C+G-H$ for even $q$ and $d-C-G+H$ for odd $q$.
When also $q=n/2$, the factor is $(-1)^{n/2}$, agreeing with the
half-offset formula at $m=n/2$. Thus no real singleton critical
mode is omitted from subsequent all-mode tests.

\subsection{Normality}
\label{slot:P04-2}

\begin{proposition}[Normality]
\label{prop:P04-normality}
Every homogeneous Jacobian \eqref{eq:P04-input} is normal over
$\C^n$: $JJ^*=J^*J$.
\end{proposition}
\begin{proof}
The unitary diagonalization in Theorem~\ref{thm:P04-spectrum}
makes $J$ normal, as in \cite[\S3.1]{Gray2006}.
\end{proof}

Normality here concerns homogeneous states, including negative ones in the domain; it is not asserted for heterogeneous or nonsynchronous Jacobians.

\subsection{Synchronous mode}

P02's scalar field, using a fresh state variable $r$, is
\[
 f(r)=ar-(k+h)r^2+\frac{(c-g)r}{1+r},\qquad r>-1.
\]

Differentiate $F(r\one)=f(r)\one$ and use $v^{(0)}=\one/\sqrt n$:
\begin{equation}
 \boxed{\lambda_0=d+C+G-H=f'(x^*).}
 \label{eq:P04-synchronous-derivative}
\end{equation}
This holds at any synchronous state, without the equilibrium hypothesis.

\paragraph{Positive synchronous equilibrium specialization.}
Now assume additionally $x^*>0$ and $F(x^*\one)=0$.
Then $C,G,H>0$. Put $A=k+h$ as in the canonical notation.
P02 and P03 give
\[
 a-Ax^*+\frac{c-g}{s}=0,\qquad
 d=-kx^*-\frac{c}{s}-\frac{cx^*}{s^2}
   =-kx^*-\frac{c(1+2x^*)}{s^2}.
\]
This is an optional equilibrium-only replacement for $d$ in
\eqref{eq:P04-spectrum}; the unsimplified expression
\eqref{eq:P04-coefficients} remains valid at all synchronous states.
Writing $f(r)=rB(r)$ with
$B(r)=a-Ar+(c-g)/(1+r)$ and $B(x^*)=0$ yields
\begin{equation}
 \boxed{\lambda_0=-x^*\left(A+\frac{c-g}{s^2}\right)
          =\frac{x^*}{s}\bigl(a-A-2Ax^*\bigr).}
 \label{eq:P04-equilibrium-zero-mode}
\end{equation}
The first equality uses $f'(x^*)=x^*B'(x^*)$; for the second,
substitute $c-g=(Ax^*-a)s$.
The equilibrium identities require a nonzero equilibrium; they also
hold at a negative nonzero equilibrium in $\Domain$, but $G,H>0$
is then false. At the origin use \eqref{eq:P04-coefficients}, which
gives $\lambda_0=a+c-g$. No comparison with other $r_m$ or selection
of a dominant mode is made here.

\subsection{Linear stability and hyperbolicity}
\label{slot:P04-3}

In this subsection $x^*\one$ is a synchronous \emph{equilibrium}
in $\Domain$. Adopt the project's spectral-abscissa notation
\begin{equation}
 \sa(J)=\max_{0\le m<n}r_m.
 \label{eq:P04-abscissa}
\end{equation}
For $\dot z=Jz$, exponential asymptotic stability means there exist
$M\ge1$ and $\beta>0$ such that
$\norm{z(t)}_2\le Me^{-\beta t}\norm{z(0)}_2$ for all $t\ge0$.

\begin{corollary}[Exact linear stability criterion]
\label{cor:P04-stability}
The linearized system is exponentially asymptotically stable if and
only if
\begin{equation}
 \boxed{\sa(J)<0\quad\Longleftrightarrow\quad
   d+C\cos\theta_m+(G-H)\cos(q\theta_m)<0
      \quad\text{for every }m.}
 \label{eq:P04-stability}
\end{equation}
For all $t\ge0$, its exact induced Euclidean propagator norm is
\begin{equation}
 \boxed{\norm{e^{tJ}}_2=e^{t\sa(J)}.}
 \label{eq:P04-semigroup}
\end{equation}
If $\sa(J)=0$, the linear system is Lyapunov stable but not
asymptotically stable. If $\sa(J)>0$, it is unstable.
\end{corollary}
\begin{proof}
Theorem~\ref{thm:P04-spectrum} gives
$e^{tJ}=Q\diag(e^{t\lambda_m})Q^*$, whose induced Euclidean norm
is $\max_m|e^{t\lambda_m}|=e^{t\sa(J)}$ for $t\ge0$.
This proves the strict criterion. At zero abscissa a maximizing real
Fourier block preserves nonzero norm; at positive abscissa it grows
exponentially. Real and complex induced norms agree for a real matrix,
so these conclusions apply to the real system as well.
\end{proof}

\begin{proposition}[Hyperbolicity]
\label{prop:P04-hyperbolicity}
The equilibrium is hyperbolic if and only if $r_m\ne0$ for every
Fourier label. It is nonhyperbolic whenever any $r_m=0$, including
when $\omega_m\ne0$. Invertibility requires instead
$\lambda_m\ne0$ for every label, and is strictly weaker than
hyperbolicity.
\end{proposition}
\begin{proof}
The complete spectrum is \eqref{eq:P04-spectrum}; substitute it in
the definitions of hyperbolicity and invertibility. A nonzero purely
imaginary eigenvalue meets the invertibility test but fails the
hyperbolicity test. For an instance within the canonical family,
P03's origin example $n=6$, $a=2,c=1,g=5/2$ (any positive $k,h$,
any canonical $q$) has $r_1=r_5=0$ and nonzero imaginary parts,
whereas no eigenvalue is zero: $ce^{-\ii\theta_m}=1/2$ is impossible
by modulus. Thus the weaker criterion is strictly weaker here as well.
\end{proof}
Hyperbolicity alone is not stability: hyperbolic equilibria may have
positive real parts. Nonhyperbolicity and instability may coexist
when one mode has zero real part and another has positive real part.

\paragraph{Nonlinear implications and their limits.}
The following conclusions concern the nonlinear equilibrium in the
open domain $\Domain$, using P01's $C^\infty$ smoothness there:
\begin{enumerate}[label=(\roman*),leftmargin=*]
\item If $\sa(J)<0$, the equilibrium is locally asymptotically stable;
in fact it is locally exponentially stable.
\item If $\sa(J)>0$, it is unstable, even if some other eigenvalues
have zero real part.
\item If $\sa(J)=0$, linearization alone is generally inconclusive
about nonlinear stability or attraction. Normality only settles the
linear system's boundedness in this case.
\end{enumerate}

These are the standard linearized stability and instability results
\cite[Ch.~1]{GuckenheimerHolmes1983}: P01 supplies smoothness on the
open domain, and the complete spectrum verifies the respective sign
condition. The instability theorem allows other eigenvalues on the
imaginary axis. At zero abscissa, a normal linearization alone does not
decide nonlinear stability (compare $\dot y=-y^3$ and $\dot y=y^3$).

All stability implications above are local. They establish neither
global attraction nor a global basin, even given P01's separate
boundedness result. At a positive equilibrium one may choose the
local ball inside $\PosOrth$, so the local assertions apply there
as well. At the origin the conclusions above are stated in
$\Domain$; instability relative to a restricted state set is a
separate question unless its admissible perturbations are checked.
A Hurwitz Jacobian is sufficient for nonlinear local asymptotic
stability, not a necessary condition in general. Thus the rejected
blanket nonlinear equivalence C040 and global assertion C153 are
not reinstated.

\paragraph{Dominant modes and first loss of stability.}
At a fixed parameter point and specified synchronous equilibrium,
the dominant-label set is $\argmax_{0\le m<n}r_m$, retaining all
ties. This is the fixed-point ``selected-mode set'' reserved in
Section~\ref{sec:notation}; it need not contain an unstable mode.
A \emph{first unstable mode along a path} additionally needs a
specified equilibrium branch $x^*(\nu)\one$, a specified oriented
parameter path $\param(\nu)$, a prior interval of stability, and a
loss-of-stability crossing event. For a continuous branch/path a
finite boundary event from $\sa(J)<0$ must have $\sa(J)=0$;
calling it an instability crossing also requires positive real part
on the unstable side, not merely tangency at zero. Several labels
or pairs can participate, so uniqueness and transverse crossing
cannot be presumed. No such path, crossing, or $q=2$ selection
analysis is supplied by this module.

\subsection{Arbitrary offset and the step-\texorpdfstring{$q$}{q} subgraph}
\label{prop:P04-step-cycles}
Let $\kappa=\gcd(n,q)$, $n=\kappa n'$ and $q=\kappa q'$,
where $\gcd(n',q')=1$. Repeated addition of $q$ returns an index
to itself after the least positive $L$ for which $n\mid Lq$.
This is equivalent to $n'\mid Lq'$, hence to $n'\mid L$,
so $L=n/\kappa$. The step-$q$ permutation therefore has exactly
$\kappa$ disjoint cycles, each of length $n/\kappa$; they are the
residue classes modulo $\kappa$. Reversing the step gives the same
partition. Thus coprimality is required only to call this step-$q$
subgraph a single cycle. The full canonical coupling also includes
the step-one predecessor cycle, and the step-$q$ partition is not
an assertion of full-network disconnection or of a block decomposition
of $J$ by those residue classes.

The Fourier sampling grid and label count depend on $n$; parity of
$n$ determines whether the $n/2$ singleton exists. The general phases
and possible eigenvalue coincidences depend on $q$ modulo $n$ as
well as $n$ and the coefficients. The value of $\lambda_0$ is
independent of $n,q$ at fixed coefficients and state; the half-mode
value depends on the parity of $q$. No proof of diagonalization,
conjugacy, multiplicity, normality, or the all-mode stability and
hyperbolicity criteria requires coprimality. The actual spectral
abscissa generally depends on the grid and phases; P03's
$n,q$-independent origin threshold is a special case.


% EDITION: citation-based shortening, 4 October 2026.
% Earlier audits refer to the source edition; this revision has not been independently audited.
\section{P05: Exact step-two mode selection}
\label{sec:P05}
Fix the canonical demand-modulated model with $q=2$, $n\ge5$ and
$a,k,c,g,h>0$. Let $x^*\one$ be a specified positive synchronous
equilibrium and put $s=1+x^*>1$. Import only the P04 coefficients
and real dispersion relation:
\begin{equation}
\begin{gathered}
 C=\frac cs>0,\qquad G=\frac{gx^*}{s^2}>0,\qquad H=hx^*>0,\\
 r_m=d+C\cos\theta_m+(G-H)\cos(2\theta_m),\qquad
 \theta_m=\frac{2\pi m}{n}.
\end{gathered}
 \label{eq:P05-input}
\end{equation}
Here $r_m=\Ree\lambda_m$ and $d$ is P04's diagonal coefficient.
All selection conclusions below are derived from \eqref{eq:P05-input}.
The maximization is over all Fourier labels, including mode zero, and
selection means the complete set maximizing the real part at this
equilibrium, retaining all ties. No coprimality condition is used.
P01 supplies smoothness, P02 supplies the equilibrium relation, and P03--P04
supply the homogeneous coefficients. No later proof module is used.


We follow the finite Fourier-dispersion selection method used by
van Kekem--Sterk for Lorenz--96 \cite[\S2.1, Proposition~5]{vanKekemSterk2018}.
Their model-specific thresholds do not transfer to DM. Here put
\begin{equation}
 T=G-H,\qquad \chi=\frac{4(H-G)}C=-\frac{4T}C.
 \label{eq:P05-control}
\end{equation}
The constructions below verify realizability with positive DM coefficients.

\subsection{Quadratic dispersion relation}
\label{slot:P05-1}
Set $u_m=\cos\theta_m$. The identity $\cos(2\theta)=2\cos^2\theta-1$
gives
\begin{equation}
 r_m=R(u_m),\qquad
 \boxed{R(u)=d-T+Cu+2Tu^2.}
 \label{eq:P05-quadratic}
\end{equation}
The additive constant is $d-T$; the part controlling selection is
$Cu+2Tu^2$. In particular, differences of growth rates and the complete
maximizer set depend only on $C,T$, not on $d$. This is an algebraic
selection statement, not permission to vary $d$ independently in the model.

Put $N=\lfloor n/2\rfloor$ and $u_j=\cos(2\pi j/n)$ for $0\le j\le N$.
Since cosine is strictly decreasing on $[0,\pi]$,
\begin{equation}
 1=u_0>u_1>\cdots>u_N,\qquad
 u_N=\begin{cases}-1,&n\text{ even},\\-\cos(\pi/n),&n\text{ odd}.
 \end{cases}
 \label{eq:P05-grid}
\end{equation}
Equation~\eqref{eq:P05-input} gives $r_{n-m}=r_m$ for $1\le m<n$.
Define the distinct label sets
\begin{equation}
 L_j=\begin{cases}
 \{0\},&j=0,\\
 \{n/2\},&n\text{ even and }j=N,\\
 \{j,n-j\},&\text{otherwise}.
 \end{cases}
 \label{eq:P05-labels}
\end{equation}
For a maximizing representative set $S$, the full dominant-label set is
$\bigcup_{j\in S}L_j$, a disjoint union.


\subsection{Continuous maximizer}
If $T\ge0$, then
$R(1)-R(u)=(1-u)[C+2T(1+u)]>0$ for $-1\le u<1$.
If $T<0$, $R$ is strictly concave with vertex
\begin{equation}
 u_*:=\frac{C}{4(H-G)}=\frac1\chi>0.
 \label{eq:P05-vertex}
\end{equation}
Thus zero is uniquely selected when $u_*\ge1$; otherwise selection
is the nearest-grid problem below. The continuous boundary $\chi=1$
need not be a discrete switch.

\subsection{Exact discrete selection rule}
\label{slot:P05-2}
For $T<0$, direct completion of the square gives
\begin{equation}
 R(u)=d+(H-G)+\frac{C^2}{8(H-G)}
           -2(H-G)(u-u_*)^2.
 \label{eq:P05-square}
\end{equation}
Consequently maximizing the growth rate is exactly minimizing
$|u_j-u_*|$. This holds for every $u_*>0$, including $u_*\ge1$.
Distances are in cosine space; the grid in that space is nonuniform.
No angular rounding rule is invoked.

Define the strictly decreasing midpoint sequence
\begin{equation}
 \beta_j=\frac{u_j+u_{j+1}}2
 =\cos\frac\pi n\cos\frac{(2j+1)\pi}{n},\qquad 0\le j<N.
 \label{eq:P05-midpoints}
\end{equation}
Strict decrease follows from
$\beta_j-\beta_{j+1}=(u_j-u_{j+2})/2>0$.

\begin{theorem}[A: Exact selection, including endpoints and all ties]
\label{thm:P05-selection}
Under the stated canonical equilibrium assumptions, if $T\ge0$ then
$S=\{0\}$. If $T<0$, the selected representatives are exactly those
whose closed cells below contain $u_*$:
\begin{align}
 V_0&=[\beta_0,\infty),\label{eq:P05-cells}\\
 V_j&=[\beta_j,\beta_{j-1}]\quad(1\le j<N),\nonumber\\
 V_N&=(-\infty,\beta_{N-1}].\nonumber
\end{align}
A representative is the unique maximizer precisely when $u_*$ lies in
the corresponding open cell. At $u_*=\beta_j$ the selected set is
exactly $\{j,j+1\}$. No nonadjacent or triple representative tie can
be maximal. Only positive midpoints can be attained when $T<0$.
The full Fourier-label set is given by \eqref{eq:P05-labels}.
\end{theorem}
\begin{proof}
The endpoint inequality handles $T\ge0$. For $T<0$, square completion
\eqref{eq:P05-square} reduces the problem to the nearest cosine.
Adjacent squared distances differ by
$2(u_j-u_{j+1})(\beta_j-u_*)$, so the midpoint cells are exactly
\eqref{eq:P05-cells}. In an ordered one-dimensional grid the nearest
points are either one point or two adjacent points at their midpoint.
Thus no other maximal tie occurs; $u_*>0$ excludes nonpositive midpoints.
\end{proof}

\subsection{Mode ties and selection cells}
\label{slot:P05-3}
The real-part comparison also has the useful exact form
\begin{equation}
 R(u_j)-R(u_{j+1})
       =C(u_j-u_{j+1})(1-\chi\beta_j).
 \label{eq:P05-adjacent-gap}
\end{equation}
For $T<0$ an adjacent \emph{maximal} tie therefore occurs if and only
if $u_*=\beta_j>0$, equivalently $\chi=1/\beta_j$.
The word maximal is essential: adjacent equalities outside the concave
regime need not be ties at the spectral abscissa.

Let $w_j=|L_j|$. A strict cell has $w_j$ selected Fourier labels; a
representative tie has $w_j+w_{j+1}$. Thus two nonsingleton
representatives give four labels, a tie involving $0$ gives three,
and an even-$n$ endpoint would contribute only one label. Below we prove
that the $N=n/2$ endpoint is never selected here. In the canonical
reachable cells every nonzero selected representative gives exactly two
labels, and every reachable nonzero/nonzero switch gives four.
These counts are counts of labels, not of distinct numerical eigenvalues.

\subsection{Reachable Fourier representatives}

\begin{theorem}[B: Reachable cells and parity]
\label{thm:P05-reachable}
The full list of representatives whose nearest-cosine cell intersects
$(0,1)$, and also the full list attainable at some canonical positive
synchronous equilibrium, is
\begin{equation}
 \boxed{0,1,\ldots,J,\qquad J=j_{\max}(n)=\left\lfloor\frac{n+1}{4}\right\rfloor.}
 \label{eq:P05-jmax}
\end{equation}
Each is attainable as the unique representative. Precisely the midpoints
$\beta_0,\ldots,\beta_{J-1}$ are positive and give attainable switches.
Every $j>J$, in particular the even-$n$ singleton $n/2$, is excluded.
Canonical attainability and robustness are proved in Theorem~\ref{thm:P05-realization}.
\end{theorem}
\begin{proof}
For $1\le j\le N$, the cell intersects $(0,1)$ if and only if
$\beta_{j-1}>0$. Necessity follows from its upper endpoint. Sufficiency
follows by taking a point strictly between
$\max(0,\beta_j)$ and $\beta_{j-1}$ when $j<N$, with the endpoint
case interpreted using $V_N$. The upper endpoint is less than $1$.
The cell for $0$ intersects $(0,1)$ since $0<\beta_0<1$.
By \eqref{eq:P05-midpoints} and $\cos(\pi/n)>0$,
\[
 \beta_j>0\ \Longleftrightarrow\ 2j+1<n/2
 \ \Longleftrightarrow\ j<(n-2)/4.
\]
Indeed the relevant midpoint angles lie in $(0,\pi)$, where cosine
changes sign only at $\pi/2$. Hence
$\beta_{j-1}>0$ is equivalent to $j<(n+2)/4$, giving exactly
\eqref{eq:P05-jmax}. In particular $J<N$ for every $n\ge5$,
$\beta_{J-1}>0$, and $\beta_J\le0$. All retained cell intersections
have nonempty interiors; there is no additional representative reachable
only through an isolated tie. When $T\ge0$ or $u_*\ge1$ only $0$
is selected, so neither regime enlarges the list.
\end{proof}

The parity information, with $n=4\ell+r$ in each row, is exact:
\begin{center}
\renewcommand{\arraystretch}{1.2}
\begin{tabular}{@{}c c c c@{}}
\toprule
$r$ & $J$ & $u_J$ & Nearest representatives at the formal point $u_*=0$\\
\midrule
$0$ & $\ell$ & $0$ & $\ell$\\
$1$ & $\ell$ & $\sin(\pi/(2n))>0$ & $\ell$\\
$2$ & $\ell$ & $\sin(\pi/n)>0$ & $\ell,\ell+1$\\
$3$ & $\ell+1$ & $-\sin(\pi/(2n))<0$ & $\ell+1$\\
\bottomrule
\end{tabular}
\end{center}
To verify the last column independently, the two cosines straddling zero
in the odd cases are
\[
 \begin{array}{ll}
 n=4\ell+1:&u_\ell=\sin(\pi/(2n)),\quad
             u_{\ell+1}=-\sin(3\pi/(2n)),\\
 n=4\ell+3:&u_\ell=\sin(3\pi/(2n)),\quad
             u_{\ell+1}=-\sin(\pi/(2n)).
 \end{array}
\]
All their arguments are in $(0,\pi/2)$, where sine increases strictly.
All remaining cosines are farther from zero by monotonicity of the grid.
For $4\mid n$ there is an exact zero cosine at $j=n/4$, and its cell
contains a positive interval. This representative is a conjugate pair,
not the real singleton $n/2$. For $n=4\ell+2$, the straddling values
are $\pm\sin(\pi/n)$, and $\beta_\ell=0$.

\paragraph{Large restraint at fixed finite dimension.}
For fixed $C>0$, as $H-G\to+\infty$, the vertex decreases to zero
and $J$ is the unique dominant representative for all sufficiently large
finite $H-G$, in every parity class. The formal zero-vertex tie for
$n=4\ell+2$ is never attained with finite positive $C,H-G$.
In fact the two competing real parts have the strictly positive,
$H-G$-independent difference
\[
 r_\ell-r_{\ell+1}=2C\sin(\pi/n)>0.
\]
Their squared cosines are equal, so this follows directly from
\eqref{eq:P05-quadratic}. Only a leading-order normalization by $H-G$
loses that difference. Thus a formal limiting nearest-grid tie must not
be added to the finite-parameter dominant set. No $n\to\infty$ limit is taken.
For $n\equiv3\pmod4$, $J$ is the sole reachable representative with a
negative cosine; for the other residue classes no negative-cosine
representative is reachable.

\paragraph{Complete regions in one dimensionless parameter.}
Put $\alpha_j=1/\beta_j$ for $0\le j<J$. Then
$1<\alpha_0<\cdots<\alpha_{J-1}$, and for \emph{all} real $\chi$
allowed by the canonical coefficients the entire selection rule is
\begin{equation}
 \begin{array}{c|c}
 \text{Condition}&S\\\hline
 \chi<\alpha_0&\{0\}\\
 \alpha_{j-1}<\chi<\alpha_j\quad(1\le j<J)&\{j\}\\
 \chi>\alpha_{J-1}&\{J\}\\
 \chi=\alpha_j\quad(0\le j<J)&\{j,j+1\}.
 \end{array}
 \label{eq:P05-regions}
\end{equation}
The middle row is empty when $J=1$. Negative and zero $\chi$ are covered
by the first row. For positive $\chi$, the inequalities follow by inverting
positive midpoint bounds, reversing their order. Thus the continuous
vertex boundary $\chi=1$ is strictly inside the uniform selection region;
the first discrete switch is at $\alpha_0>1$.

\subsection{Canonical parameter realization}

The concave condition and vertex in original variables are
\begin{equation}
 T<0\ \Longleftrightarrow\ h>\frac{g}{s^2},\qquad
 u_*=\frac{c/s}{4x^*(h-g/s^2)}
     =\frac{cs}{4x^*(hs^2-g)},\qquad
 \chi=\frac{4x^*(hs^2-g)}{cs}.
 \label{eq:P05-canonical-control}
\end{equation}
All denominators in the vertex formula are strictly positive in this
regime. For $0\le j<J$, the switch equations are exactly
\begin{equation}
 H-G=\frac{C}{4\beta_j},\qquad
 \boxed{h=\frac{g}{(1+x^*)^2}
       +\frac{c}{4(1+x^*)x^*\beta_j}.}
 \label{eq:P05-switch}
\end{equation}
They are subject to P02's equilibrium condition
\begin{equation}
 a-(k+h)x^*+\frac{c-g}{1+x^*}=0.
 \label{eq:P05-equilibrium}
\end{equation}
For $\beta_j\le0$ there is no switch in the $T<0$ regime: division
by zero is not meaningful and negative $\beta_j$ would require
$H-G<0$. Such negative-midpoint equalities in the convex regime, if
realizable as equal lower values, are not dominant-mode switches.

\begin{theorem}[C: Canonical realizability, open cells and switch geometry]
\label{thm:P05-realization}
Every strict cell and every switch in \eqref{eq:P05-regions} is nonempty
in the canonical positive-equilibrium parameter space. Each strict cell
has a realization persisting on a nonempty open set of positive physical
parameters $(a,k,c,g,h)$ with a specified smooth positive synchronous
branch. Both $T>0$ and $T<0$, and the boundary $T=0$, are realizable.
Every positive-midpoint switch is locally a smooth codimension-one
hypersurface in physical parameter space along its specified branch;
all its positive-equilibrium points are regular branch points. No
nonpositive-midpoint dominant switch is realizable. These are local
branch statements and assert no global identification of distinct branches.
\end{theorem}
\begin{proof}[Proof of realizability]
Choose any $x>0$, any $c,g>0$, and any desired vertex $v>0$. Define
\begin{equation}
 s=1+x,\qquad h=\frac{g}{s^2}+\frac{c}{4sxv},\qquad
 k>\max\left\{0,\frac{c-g}{sx}-h\right\},\qquad
 a=(k+h)x-\frac{c-g}{s}.
 \label{eq:P05-construction}
\end{equation}
Then $h,k,a$ are strictly positive: the displayed strict lower bound on
$k$ is exactly what is required for $a>0$, together with $k>0$.
There is no upper bound on $k$, so this choice is always possible.
Equation~\eqref{eq:P05-equilibrium} vanishes identically at $x^*=x$;
substitution into the original vector field gives, in every component,
\[
 ax-kx^2+\frac{cx}{s}-\frac{gx}{s}-hx^2
 =x\left[a-(k+h)x+\frac{c-g}{s}\right]=0.
\]
Also $H-G=c/(4sv)>0$ and $u_*=v$. Hence every retained open
Voronoi cell and every positive midpoint can be prescribed exactly.
The construction does not treat $C,G,H$ as independent of an equilibrium.

A particularly simple exact family, used below, is
\begin{equation}
 x^*=1,\qquad c=g=4,\qquad h=1+\chi/2,\qquad
 k>0,\qquad a=k+1+\chi/2,\qquad \chi>-2.
 \label{eq:P05-simple-family}
\end{equation}
Here $C=2$, $G=1$, $H=1+\chi/2$, so the dimensionless parameter
is exactly the displayed $\chi$. Every positive $\chi$ is realized.
At $k=1$, choosing $\chi=-1,0,1$ realizes respectively $T>0$,
$T=0$, and $T<0$ with $u_*=1$, all selecting $0$.
Thus the uniform region has examples in both sign regimes.
The geometry and openness assertions are proved next.
\end{proof}

\subsection{Realizable switch hypersurfaces}

Use the equilibrium coordinates $(x,k,c,g,h)$ on the open set
\begin{equation}
 \mathcal U_{\rm eq}=
 \left\{(x,k,c,g,h)\in(0,\infty)^5:
     (k+h)x-\frac{c-g}{1+x}>0\right\}.
 \label{eq:P05-equilibrium-chart}
\end{equation}
The map
\[
 \Phi(x,k,c,g,h)=\left((k+h)x-\frac{c-g}{1+x},\ k,c,g,h\right)
\]
projects this parameterization of the positive equilibrium graph onto
physical parameters. With $s=1+x$ its Jacobian determinant is
\begin{equation}
 \det D\Phi=k+h+\frac{c-g}{s^2}.
 \label{eq:P05-chart-determinant}
\end{equation}
At every point with $H-G>0$, this is
$k+(h-g/s^2)+c/s^2>0$. Equivalently, for the scalar equilibrium
function $B(x)=a-(k+h)x+(c-g)/(1+x)$,
$B_x=-\det D\Phi<0$. The inverse function theorem makes $\Phi$ a
local diffeomorphism there, and the implicit function theorem gives a
specified smooth positive branch $x=x^*(a,k,c,g,h)$. This regularity
uses only the scalar relation, not full-system invertibility or a Hopf result.

Strict inequalities in \eqref{eq:P05-regions} define open subsets of
$\mathcal U_{\rm eq}$. Every strict cell has a $\chi>0$ point from
\eqref{eq:P05-construction}, including the uniform cell, so local
projection supplies a nonempty open set of physical parameters for each.
In fact every concave-regime point of a strict cell has this local property.
For a $T>0$ uniform example, the family at $k=1,\chi=-1$ has
$\det D\Phi=k+h=3/2>0$; strict $T>0$ persists in a neighborhood.
No global openness assertion at an arbitrary scalar fold in the uniform
region is required or inferred.

For $\beta_j>0$ define in equilibrium coordinates
\begin{equation}
 \Psi_j(x,k,c,g,h)=h-\frac{g}{(1+x)^2}
                  -\frac{c}{4(1+x)x\beta_j}.
 \label{eq:P05-level-set}
\end{equation}
Its partial derivative with respect to $h$, holding $x,k,c,g$ fixed,
is exactly $1$. Hence its zero set is a smooth codimension-one regular
level set within $\mathcal U_{\rm eq}$. It is nonempty by the construction,
and every one of its points has $H-G>0$, so the local diffeomorphism
$\Phi$ transfers this conclusion to physical parameter space on the
chosen branch. Moreover, at fixed $x,c,g,k$,
\[
 \frac{\partial\chi}{\partial h}=\frac{4(1+x)x}{c}>0,
 \qquad
 \frac{\partial}{\partial h}(r_j-r_{j+1})
       =-4x(u_j-u_{j+1})\beta_j<0.
\]
Here $a$ varies according to \eqref{eq:P05-equilibrium}; this is a
legitimate direction in physical parameter space. Distinct midpoint
thresholds and the strict gap to all other representatives at the tie
show that the two sides locally select $j$ and $j+1$, respectively.
This proves genuine local changes of dominance, and finishes
Theorem~\ref{thm:P05-realization}. No global manifold structure after
forgetting branch labels is claimed.

\subsection{Exact examples and small-dimension checks}

In \eqref{eq:P05-simple-family} take $k=1$. The tuple
$(a,k,c,g,h)=(2,1,4,4,1)$ selects $0$ in every dimension.
For $n=8$, $(3,1,4,4,2)$ has $\chi=2$, $u_*=1/2$, and selects
representative $1$, since
$\sqrt2/4<1/2<(2+\sqrt2)/4$.
For $n=6$, $(8/3,1,4,4,5/3)$ has $\chi=4/3=1/\beta_0$ and
selects exactly labels $\{0,1,5\}$.
For $n=8$, $(2+\sqrt2,1,4,4,1+\sqrt2)$ has
$\chi=2\sqrt2=1/\beta_1$ and selects $\{1,2,6,7\}$.
Every example has $x^*=1$, all five coefficients positive, and original
DM residual $a-k-h+(c-g)/2=0$.

For exact highest-mode examples across all four residue classes, use
$k=1$ in the same family with $\chi=2/\beta_{J-1}$. Then
$u_*=\beta_{J-1}/2\in(0,\beta_{J-1})$, while $\beta_J\le0$,
so $J$ is strict. In particular the following are exact specifications:
\begin{center}
\begin{tabular}{@{}c c c c@{}}
\toprule
$n$ & $J$ & $\chi$ & $u_J$\\\midrule
$8$ & $2$ & $4\sqrt2$ & $0$\\
$9$ & $2$ & $4/\cos(\pi/9)$ & $\sin(\pi/18)$\\
$10$ & $2$ & $2/[\cos(\pi/10)\cos(3\pi/10)]$ & $\sin(\pi/10)$\\
$11$ & $3$ & $2/[\cos(\pi/11)\cos(5\pi/11)]$ & $-\sin(\pi/22)$\\
\bottomrule
\end{tabular}
\end{center}
The corresponding tuples are exactly
$(2+\chi/2,1,4,4,1+\chi/2)$. No numerical search is used.

The next table records the complete representative grid and positive
midpoints for $5\le n\le12$. Within each row only, write
$\gamma_r=\cos(r\pi/n)$. All listed midpoint values are exact.
Their reciprocal values are the switch thresholds in
\eqref{eq:P05-regions}; the reachable set is $\{0,\ldots,J\}$.
The terminal midpoint $\beta_J$ is negative except at $n=6,10$,
where it is zero and unreachable. The five-node row is solely a check;
the golden-ratio specialization is reserved for P07.
\begin{center}
\small
\renewcommand{\arraystretch}{1.3}
\begin{tabular}{@{}c l c l@{}}
\toprule
$n$ & $(u_0,\ldots,u_N)$ & $J$ & $(\beta_0,\ldots,\beta_{J-1})$\\
\midrule
$5$ & $(1,\gamma_2,-\gamma_1)$ & $1$ & $(\gamma_1^2)$\\
$6$ & $(1,1/2,-1/2,-1)$ & $1$ & $(3/4)$\\
$7$ & $(1,\gamma_2,-\gamma_3,-\gamma_1)$ & $2$ & $(\gamma_1^2,\gamma_1\gamma_3)$\\
$8$ & $(1,\sqrt2/2,0,-\sqrt2/2,-1)$ & $2$ & $((2+\sqrt2)/4,\sqrt2/4)$\\
$9$ & $(1,\gamma_2,\gamma_4,-1/2,-\gamma_1)$ & $2$ & $(\gamma_1^2,\gamma_1/2)$\\
$10$ & $(1,\gamma_2,\gamma_4,-\gamma_4,-\gamma_2,-1)$ & $2$ & $(\gamma_1^2,\gamma_1\gamma_3)$\\
$11$ & $(1,\gamma_2,\gamma_4,-\gamma_5,-\gamma_3,-\gamma_1)$ & $3$ & $(\gamma_1^2,\gamma_1\gamma_3,\gamma_1\gamma_5)$\\
$12$ & $(1,\sqrt3/2,1/2,0,-1/2,-\sqrt3/2,-1)$ & $3$ & $((2+\sqrt3)/4,(\sqrt3+1)/4,1/4)$\\
\bottomrule
\end{tabular}
\end{center}

\subsection{Dominant modes versus first instability}

The preceding results concern $\argmax_m r_m$ at a fixed equilibrium.
The scalar $d$ cancels from that comparison but affects whether the
maximum is negative, zero, or positive. A first loss of stability would
require a specified oriented physical path, a positive equilibrium branch,
a prior interval with $\max_m r_m<0$, and a crossing to a positive
maximum. Neither a selection switch nor a real-part tie establishes that.

Indeed every reachable switch can occur on a path that stays strictly
stable. Fix its threshold $\alpha_j$, and choose a compact interval
$[b_-,b_+]\subset(0,\infty)$ with $\alpha_j$ in its interior and no
other threshold in it. In \eqref{eq:P05-simple-family} let $\chi$
increase across this interval, hold $k>b_+/2$ fixed, and keep $x^*=1$.
The equilibrium is regular, all parameters remain positive, and P04's
equilibrium diagonal gives $d=-k-3$. Thus directly from
\eqref{eq:P05-input},
\begin{equation}
 r_m=-k-3+2\cos\theta_m-\frac\chi2\cos(2\theta_m)
       \le-k-1+\frac{b_+}{2}<-1.
 \label{eq:P05-stable-switch}
\end{equation}
The dominant representative changes from $j$ to $j+1$, with the exact
tie at $\alpha_j$, while every mode remains strictly stable. This also
shows that every strict cell has a stable canonical realization, since
$k$ can be enlarged without changing $\chi$ or $x^*=1$ in this family.
No first-instability subresult is asserted; that separate crossing question
is deferred to P11/P12 under their branch/path gates.

\paragraph{Real-part ties are not eigenvalue coincidences.}
At every switch P05 proves $r_j=r_{j+1}$ only. P04 supplies the signed
imaginary parts
\[
 \omega_m=-C\sin\theta_m+(G+H)\sin(2\theta_m).
\]
For particular labels, equality of the complex eigenvalues would additionally
require equality of these signed imaginary parts. For conjugate labels the
sign changes. None of the selection, realization, or geometry proofs needs
an imaginary-part condition. Numerical eigenvalue simplicity, nonzero
critical frequency, simultaneous critical pairs and Hopf exclusions remain
P11's tasks. Counting two or four selected labels does not certify a simple
pair or a multiple-eigenvalue collision.


% STATUS: conditional framework; outstanding model-specific work is explicit.
\section{P06: Large-size limit (unfinished)}
\label{sec:P06}
The supplied file contains an outline, not a completed limit proof.
\subsection{Selected-wave-number limit}\label{slot:P06-1}
Use P05's nearest-cosine rule. The method is analogous to the
Lorenz--96 limit in \cite[Proposition~6]{vanKekemSterk2018}, but
its model-specific limit cannot be imported as a DM theorem.
\subsection{Discretization and uniformity}\label{slot:P06-2}
Still required: a specified parameter regime, control of the maximizing
grid point, endpoint and tie cases, and any uniform-in-$n$ error estimate.


% EDITION: citation-based shortening, 4 October 2026.
% Earlier audits refer to the source edition; this revision has not been independently audited.
\section{P07: Five-node golden-ratio specialization}
\label{sec:P07}
This is a corollary package for the P03--P05 demand-modulated
results. The direct dependency remains $\mathrm{P05}\to\mathrm{P07}$;
P03--P04 and the P01--P02 equilibrium foundation are inherited.
P08 is present in the current project but supplies no theorem used here.

Fix $a,k,c,g,h>0$ in the canonical model
\begin{equation}
 \dot x_i=ax_i-kx_i^2+\frac{cx_{i-1}}{1+x_i}
       -\frac{gx_i}{1+x_{i+2}}-hx_ix_{i-2},\qquad i\in\Z_5.
 \label{eq:P07-model}
\end{equation}
Let $x^*\one$, $x^*>0$, be a specified synchronous equilibrium. Write
\begin{equation}
 s=1+x^*>1,\quad C=\frac cs>0,\quad G=\frac{gx^*}{s^2}>0,
 \quad H=hx^*>0,\quad T=G-H,\quad S=G+H.
 \label{eq:P07-coefficients}
\end{equation}
Here $s$ denotes the denominator, not P02's scalar state variable;
$S$ is local to P07 and denotes effective DM coefficients, not the
physical sum $g+h$ used locally in P08. We use $D_5:=H-G=-T$,
leaving P02's discriminant $D$ unchanged. The symbols $\omega_1,\omega_2$
below are signed imaginary parts, not a reserved positive Hopf frequency.
By P02 and P03,
\begin{align}
 a-(k+h)x^*+\frac{c-g}{s}&=0,\label{eq:P07-equilibrium}\\
 d&=a-2kx^*-\frac{cx^*}{s^2}-\frac gs-hx^*
   =-kx^*-C-\frac{Cx^*}{s}.\label{eq:P07-d}
\end{align}
The last equality uses equilibrium and will only be used when stated.

\subsection{Exact fifth-root constants and Fourier spectrum}
\label{slot:P07-1}
\begin{lemma}[Golden-ratio trigonometry]
\label{lem:P07-trig}
Set
\begin{equation}
 \phi=\frac{1+\sqrt5}{2},\qquad
 u_1=\cos\frac{2\pi}{5},\qquad u_2=\cos\frac{4\pi}{5},\qquad
 \sigma=\sin\frac\pi5.
 \label{eq:P07-constants}
\end{equation}
Then
\begin{gather}
 \phi^2=\phi+1,\quad \phi^{-1}=\phi-1,\quad
 \phi^{-2}=2-\phi=\frac{3-\sqrt5}{2},\quad
 \phi+\phi^{-1}=\sqrt5,\label{eq:P07-phi}\\
 u_1=\frac{\sqrt5-1}{4}=\frac1{2\phi},\qquad
 u_2=-\frac{\sqrt5+1}{4}=-\frac\phi2,\qquad
 u_1-u_2=\frac{\sqrt5}{2},\label{eq:P07-cosines}\\
 \sigma=\frac{\sqrt{10-2\sqrt5}}4>0,\qquad
 \sin\frac{2\pi}{5}=\phi\sigma
       =\frac{\sqrt{10+2\sqrt5}}4.\label{eq:P07-sines}
\end{gather}
\end{lemma}
\begin{proof}
These are the exact fifth-root trigonometric values; substituting
$\phi=(1+\sqrt5)/2$ verifies the displayed identities.
\end{proof}


There are three representative classes, $\{0\}$, $\{1,4\}$ and
$\{2,3\}$. Substitute the constants above into P04 with
$\theta_m=2\pi m/5$.

\begin{theorem}[Exact five-node Fourier spectrum]
\label{thm:P07-spectrum}
The five eigenvalues, with their Fourier labels and algebraic
multiplicities retained, are
\begin{align}
 \lambda_0&=d+C+T,\label{eq:P07-spectrum}\\
 \lambda_1&=d+\frac{C}{2\phi}-\frac{\phi T}{2}
                +\ii\sigma(S-\phi C),\nonumber\\
 \lambda_2&=d-\frac{\phi C}{2}+\frac{T}{2\phi}
                -\ii\sigma(C+\phi S),\nonumber\\
 \lambda_4&=\overline{\lambda_1},\qquad
 \lambda_3=\overline{\lambda_2}.\nonumber
\end{align}
In particular the three real-part classes are
\begin{equation}
 r_0=d+C+T,\qquad
 r_1=d+\frac{C}{2\phi}-\frac{\phi T}{2},\qquad
 r_2=d-\frac{\phi C}{2}+\frac{T}{2\phi}.
 \label{eq:P07-real}
\end{equation}
Here $r_1=\Ree\lambda_1=\Ree\lambda_4$ and
$r_2=\Ree\lambda_2=\Ree\lambda_3$.
\end{theorem}
\begin{proof}
Specialize \eqref{eq:P04-real}--\eqref{eq:P04-imag} to $n=5,q=2$
using Lemma~\ref{lem:P07-trig}; the remaining labels follow by conjugacy.
\end{proof}

\subsection{Complete ordering and the dominant-mode switch}
\label{slot:P07-2}
\begin{theorem}[Complete ordering of the three real-part classes]
\label{thm:P07-ordering}
Set $t=T/C$, which is well defined since $C>0$. The exact differences are
\begin{align}
 r_1-r_2&=\frac{\sqrt5}{2}(C-T)
              =\frac{\sqrt5 C}{2}(1-t),\label{eq:P07-gaps}\\
 r_0-r_1&=\frac{\sqrt5\phi}{2}
                    \left(T+\frac{C}{\phi^2}\right)
              =\frac{\sqrt5\phi C}{2}(t+\phi^{-2}),\nonumber\\
 r_0-r_2&=\frac{\sqrt5}{2\phi}(T+\phi^2 C)
              =\frac{\sqrt5 C}{2\phi}(t+\phi^2).\nonumber
\end{align}
Consequently the following table is exhaustive, including all equalities:
\begin{center}
\renewcommand{\arraystretch}{1.15}
\begin{tabular}{@{}l l l@{}}
\toprule
Condition&Ordering&Tie type\\
\midrule
$t>1$&$r_0>r_2>r_1$&none\\
$t=1$&$r_0>r_1=r_2$&subdominant\\
$-\phi^{-2}<t<1$&$r_0>r_1>r_2$&none\\
$t=-\phi^{-2}$&$r_0=r_1>r_2$&dominant\\
$-\phi^2<t<-\phi^{-2}$&$r_1>r_0>r_2$&none\\
$t=-\phi^2$&$r_1>r_0=r_2$&subdominant\\
$t<-\phi^2$&$r_1>r_2>r_0$&none\\
\bottomrule
\end{tabular}
\end{center}
Ordering is independent of $d$. No triple tie of real-part classes occurs.
\end{theorem}
\begin{proof}
Subtract the three expressions in \eqref{eq:P07-real} and use
$\phi^2=\phi+1$. The positive prefactors in \eqref{eq:P07-gaps}
and the strict ordering $-\phi^2<-\phi^{-2}<1$ give every row.
Their three distinct zeros exclude a triple tie.
\end{proof}

\begin{corollary}[Exact golden-ratio selection threshold]
\label{cor:P07-selection}
The dominant representative set and the full selected Fourier-label set are
\begin{equation}
\begin{array}{c|c|c}
\text{Condition}&\text{Representatives}&\text{Labels}\\\hline
D_5/C<\phi^{-2}&\{0\}&\{0\}\\
D_5/C=\phi^{-2}&\{0,1\}&\{0,1,4\}\\
D_5/C>\phi^{-2}&\{1\}&\{1,4\}.
\end{array}
\label{eq:P07-selection}
\end{equation}
Pair $\{2,3\}$ is never dominant at a canonical positive synchronous
 equilibrium. The unique dominant-class switch is
\begin{equation}
 \boxed{H-G=\frac{C}{\phi^2},\qquad
         \frac{H-G}{C}=\frac{3-\sqrt5}{2}.}
 \label{eq:P07-switch}
\end{equation}
\end{corollary}
\begin{proof}
Apply Theorem~\ref{thm:P05-selection} with
$\beta_0=\phi^2/4$ and $\beta_1=-1/4$, or read the maxima
in Theorem~\ref{thm:P07-ordering}.
\end{proof}


The five-node midpoint and control values are
\begin{equation}
 \beta_0=\frac{\phi^2}{4},\quad\beta_1=-\frac14,
 \quad\alpha_0=\frac4{\phi^2}=6-2\sqrt5.
 \label{eq:P07-midpoints}
\end{equation}
\begin{equation}
 \chi=\frac{4(H-G)}C,\qquad
 \chi\lessgtr\frac4{\phi^2}\ \Longleftrightarrow\
 \text{strict selection of representative }0\text{ or }1,
 \label{eq:P07-chi}
\end{equation}
respectively; equality selects both, as in \eqref{eq:P07-selection}.

\paragraph{Canonical switch surface and the equilibrium constraint.}
Since $D_5=x^*(h-g/s^2)$, solving \eqref{eq:P07-switch} gives
\begin{equation}
 \boxed{h=\frac{g}{(1+x^*)^2}
            +\frac{c}{\phi^2(1+x^*)x^*}.}
 \label{eq:P07-canonical-switch}
\end{equation}
The divisions are valid because $x^*,s,\phi,C$ are positive.
The equilibrium relation \eqref{eq:P07-equilibrium} must hold
simultaneously. Thus this equation describes a switch on the positive
branch-state manifold; $x^*$ cannot be prescribed independently of
a fixed physical tuple. The corresponding strict inequalities in $h$
at that equilibrium give uniform or pair-1 selection, respectively.
P05 already supplies regularity and local realizability of this switch.

\paragraph{The other golden-ratio equality surfaces.}
At $t=1$, equivalently $G-H=C$, one has $r_1=r_2$ but
$r_0-r_1=r_0-r_2=5C/2>0$. At $t=-\phi^2$, equivalently
$H-G=\phi^2C$, one has $r_0=r_2$ but
$r_1-r_0=r_1-r_2=5\phi C/2>0$.
These values follow from \eqref{eq:P07-gaps} and $\phi^2=\phi+1$.
Both are subdominant real-part ties, not dominant-mode switches.
No equality of complex numerical eigenvalues is implied by either tie.

\subsection{Frequencies, intersection, and numerical eigenvalue equality}
\begin{proposition}[Signed imaginary parts and frequency degeneracy]
\label{prop:P07-frequency}
At positive synchrony,
\begin{equation}
 \omega_1:=\Imm\lambda_1=\sigma(S-\phi C),\qquad
 \omega_2:=\Imm\lambda_2=-\sigma(C+\phi S)<0.
 \label{eq:P07-frequency}
\end{equation}
Pair $\{2,3\}$ is always nonreal. Pair $\{1,4\}$ is nonreal
if and only if $S\ne\phi C$. At the spectral locus
\begin{equation}
 \boxed{G+H=\phi C,}\label{eq:P07-frequency-locus}
\end{equation}
one has $\lambda_1=\lambda_4=r_1$, a repeated real numerical
 eigenvalue represented by two Fourier labels. All numerical eigenvalues
are semisimple; any further coincident label is counted separately below.
\end{proposition}
\begin{proof}
The imaginary terms were derived directly in Theorem~\ref{thm:P07-spectrum}.
Since $\sigma,C,S,\phi>0$, $\omega_2$ is strictly negative, while
$\omega_1=0$ is equivalent to $S=\phi C$. Conjugacy then proves
all reality assertions. Semisimplicity is
Proposition~\ref{prop:P04-multiplicity}, which does not assume distinct
 numerical eigenvalues. These assertions classify spectra only.
\end{proof}

\begin{proposition}[Intersection with the dominant-mode switch]
\label{prop:P07-intersection}
The switch and the frequency-degeneracy locus intersect in the strictly
positive coefficient domain exactly at
\begin{equation}
 \boxed{\frac GC=\frac1\phi,\qquad \frac HC=1.}
 \label{eq:P07-intersection}
\end{equation}
They define a nonempty codimension-two spectral locus in $(C,G,H)>0$,
and locally on the canonical positive equilibrium branch. At that locus
\begin{equation}
 \lambda_0=\lambda_1=\lambda_4=d+\frac C\phi
   =-kx^*-\frac{Cx^*}{s}-\frac{C}{\phi^2}<0.
 \label{eq:P07-triple}
\end{equation}
This numerical eigenvalue is semisimple of multiplicity three, and
$r_2<r_0$. In particular the full linearization is Hurwitz there.
This is a spectral codimension statement, not a bifurcation classification.
\end{proposition}
\begin{proof}
Adding and subtracting $H-G=C\phi^{-2}$ and $G+H=\phi C$ gives
\[
 \frac GC=\frac{\phi-\phi^{-2}}2=\phi-1=\phi^{-1},\qquad
 \frac HC=\frac{\phi+\phi^{-2}}2=1.
\]
Both are positive. The gradients in $(G,H)$ of
$H-G-C/\phi^2$ and $G+H-\phi C$ have determinant $-2$, so the
two linear conditions are independent. For a canonical witness take
\begin{equation}
 x^*=1,\qquad (a,k,c,g,h)=(\sqrt5,1,2,4/\phi,1).
 \label{eq:P07-intersection-witness}
\end{equation}
The equilibrium bracket equals $\sqrt5-2+(2-4/\phi)/2=0$,
using $1+2/\phi=\sqrt5$. The effective coefficients are
$C=1,G=\phi^{-1},H=1$.

For the local branch assertion use P05's verified equilibrium chart
$(x,k,c,g,h)\mapsto((k+h)x-(c-g)/(1+x),k,c,g,h)$.
Its determinant is positive wherever $H-G>0$; the intersection lies
in that regime. At fixed $x,k,c$, the two constraint gradients in
$(g,h)$ have determinant $-2x^2/(1+x)^2\ne0$.
The chart therefore transfers their independent codimension-two
conditions locally to physical parameters on the specified branch.
There is no claim about a parameter-path crossing or an unfolding.

At the intersection, conjugacy, $\omega_1=0$, and $r_0=r_1$
give the triple value. Substitute \eqref{eq:P07-d} and
$1-\phi^{-1}=\phi^{-2}$ for its strict negativity.
The ordering table gives $r_2<r_0$; $\omega_2<0$ keeps labels 2 and 3
distinct from the real value and from each other. P04 counts precisely
three independent Fourier eigenvectors for the triple value.
\end{proof}

\begin{proposition}[Complete numerical coincidence classification]
\label{prop:P07-coincidences}
Under the stated positive-equilibrium assumptions, the only numerical
 eigenvalue coincidences are as follows:
\begin{enumerate}[label=(\roman*),leftmargin=*]
\item If $S\ne\phi C$, all five numerical eigenvalues are distinct.
\item If $S=\phi C$ and $T\ne-C/\phi^2$, only
$\lambda_1=\lambda_4$ coincide; their real value has multiplicity two.
\item If $S=\phi C$ and $T=-C/\phi^2$, exactly
$\lambda_0=\lambda_1=\lambda_4$ coincide as in
\eqref{eq:P07-triple}; their value has multiplicity three.
\end{enumerate}
All multiplicities are algebraic and geometric. Labels 2 and 3 always
give distinct, simple numerical eigenvalues, sharing no value with any
other label. Every nonreal pair-1 eigenvalue is simple.
\end{proposition}
\begin{proof}
Within pairs, \eqref{eq:P07-frequency} proves that labels 1 and 4
coincide exactly when $S=\phi C$, while labels 2 and 3 never coincide.
The real $\lambda_0$ cannot equal either nonreal pair-2 eigenvalue.
Equality between $\lambda_0$ and a pair-1 label requires both
$r_0=r_1$ and $\omega_1=0$, exactly the intersection just analyzed.

A cross-pair numerical equality requires $r_1=r_2$, hence $T=C$.
Then $G=H+C$ and $S=C+2H>C$. At this real-part tie,
\begin{align*}
 \omega_1-\omega_2
    &=\sigma[(1+\phi)S-(\phi-1)C]>0,\\
 -\omega_2-\omega_1
    &=\sigma[(1+\phi)C+(\phi-1)S]>0.
\end{align*}
For the first inequality use $S>C$ and
$(1+\phi)-(\phi-1)=2$; the second has strictly positive summands.
Thus $\omega_2<\omega_1<-\omega_2$, so
$\omega_1$ is neither $\omega_2$ nor $-\omega_2$; conjugating
excludes all four possible cross-pair equalities. This exhausts all
unordered label pairs. P04 then gives the asserted multiplicities.
\end{proof}
The frequency-degeneracy locus itself has $-\phi<t<\phi$, since
$G,H>0$ imply $|T|<S$. It can intersect the subdominant tie $t=1$
without a cross-pair coincidence, but cannot intersect $t=-\phi^2$.
Equal real parts, repeated real values, and distinct nonreal conjugates
are consequently kept separate throughout.

\subsection{Exact stability test and critical-mode ledger}
\begin{corollary}[Piecewise five-node linear stability]
\label{cor:P07-stability}
The linearization is exponentially asymptotically stable exactly when
\begin{equation}
 \max\{r_0,r_1,r_2\}<0
 \quad\Longleftrightarrow\quad
 \begin{cases}
 r_0<0,&t>-\phi^{-2},\\
 r_0=r_1<0,&t=-\phi^{-2},\\
 r_1<0,&t<-\phi^{-2}.
 \end{cases}
 \label{eq:P07-stability}
\end{equation}
The middle row means that the common value must be negative. At canonical
positive equilibria on that tie it is automatically negative:
\begin{equation}
 r_0=r_1=-kx^*-\frac{Cx^*}{s}-\frac{C}{\phi^2}<0.
 \label{eq:P07-stable-tie}
\end{equation}
\end{corollary}
\begin{proof}
P04 Corollary~\ref{cor:P04-stability} gives the strict all-mode test;
Theorem~\ref{thm:P07-ordering} gives its piecewise maximum.
At the dominant tie, substitute $T=-C/\phi^2$ and
\eqref{eq:P07-d} into $r_0=d+C+T$.
\end{proof}
If the maximum is zero, P04 implies that the normal linear system is
Lyapunov stable but not asymptotically stable. This does not settle
nonlinear stability or establish a bifurcation. A strictly positive
maximum means linear instability. The dominant tie cannot itself be a
zero-abscissa boundary in the canonical domain, by
\eqref{eq:P07-stable-tie}.

\begin{proposition}[Fixed-state critical-mode ledger]
\label{prop:P07-ledger}
Only the following classes can attain the spectral abscissa; this is a
fixed-state eligibility ledger, not a first-instability theorem.
\begin{center}
\begin{tabular}{@{}l l l@{}}
\toprule
Class&Eligibility for the maximum&Signed imaginary part\\
\midrule
$0$&$t\ge-\phi^{-2}$&$0$\\
$\{1,4\}$&$t\le-\phi^{-2}$&$\pm\sigma(S-\phi C)$\\
$\{2,3\}$&never&$\mp\sigma(C+\phi S)$\\
\bottomrule
\end{tabular}
\end{center}
At a zero-abscissa state, the inequalities in the first two rows must
be strict, because the shared boundary is strictly stable.
If $r_2=0$, some other class has strictly positive real part, so pair 2
cannot be the sole or tied controlling class of a stability boundary.
No oriented parameter path, prior stability, or crossing is asserted.
\end{proposition}
\begin{proof}
The eligibility and exclusions are precisely
Corollary~\ref{cor:P07-selection}. At every point at least one of
$r_0,r_1$ exceeds $r_2$ strictly. The zero-abscissa statements follow
from that strict comparison and \eqref{eq:P07-stable-tie}.
\end{proof}

\subsection{Realizability checks, downstream inputs, and scope}
\paragraph{All ordering regions have canonical representatives.}
For any $C,G,H>0$, the following single construction gives a canonical
positive equilibrium with those exact effective coefficients:
\begin{equation}
 x^*=1,\quad k=C,\quad c=2C,\quad g=4G,\quad h=H,
 \quad a=H+2G>0,\quad d=-\frac{5C}{2}.
 \label{eq:P07-realization}
\end{equation}
Indeed \eqref{eq:P07-equilibrium} becomes
$H+2G-(C+H)+(2C-4G)/2=0$. Every original component equals
$x^*$ times that zero bracket. For any real $t$, take $C=1$,
$G=1+\max(t,0)$ and $H=1+\max(-t,0)$. Then $G-H=t$ and both
coefficients are positive. Thus no row of the ordering table is merely
an unattached-coefficient possibility. This elementary five-node check
uses no MB comparison or new general selection theorem.

\paragraph{Exact inputs available to P11.}
The unique nonuniform dominant class, when present, is $\{1,4\}$.
Its signed imaginary part is $\sigma(S-\phi C)$, and its exact
nonzero-frequency condition is $S\ne\phi C$. Whenever pair 1 is
maximal, $r_2<r_1$ strictly, including the uniform/pair-1 dominant tie.
When pair 1 is uniquely dominant, $r_0<r_1$ as well.
Every nonreal pair-1 eigenvalue is simple; pair 2 never coincides
numerically with another label. The only uniform/pair-1 numerical
coincidence is the strictly stable triple value
\eqref{eq:P07-triple}; more generally the entire dominant tie is
strictly stable. These are exact spectral inputs, without a Hopf
existence or first-crossing conclusion. No transversality, crossing
derivative, Lyapunov coefficient, normal form, periodic orbit, or rotating
wave is analyzed here; Later sections address their own hypotheses.

\begin{remark}[Golden ratio is network geometry]
\label{rem:P07-geometry}
The constants $\phi$ and $\sigma$ arise solely from the fifth roots of
unity and the five-node Fourier grid. They establish no physiological
golden-ratio law, biological optimum, medical significance, or empirical
TCM validation. P07 is a supporting specialization of the general program, not standalone golden-ratio novelty. 
\end{remark}


% EDITION: citation-based shortening, 4 October 2026.
% Earlier audits refer to the source edition; this revision has not been independently audited.
\section{P08: Matched bilinear analysis and exact comparison}
\label{sec:P08}

The global-interior theorem and its comparison corollary were already
present as unaudited extensions in the supplied edition. Their shortened
proofs below retain that status; the earlier A08 audit did not cover them.
All coefficients are strictly positive, $n\ge5$, $1\le q<n$, and indices
are modulo $n$. The comparator is exactly the five-parameter model
\eqref{eq:canonical-mb}. The same tuple $(a,k,c,g,h)$, $n,q$ and time
variable are used in DM and MB; each equilibrium amplitude is derived
for its own equation. The foundation and selection inputs are P01--P05.

The substantive results are open same-tuple stability contrasts in both
directions at $n=7$, a universal forward contrast, and realization of
every ordered pair of reachable dominant representatives with both
models stable. None is a first-instability, Hopf, criticality, or
literature-priority theorem. Here ``strictly stable'' means Hurwitz;
``unstable'' means a strictly positive spectral abscissa. Zero abscissa
is treated separately. Selection maximizes over all labels, including
zero, retaining every maximal tie.

\subsection{MB local flow, synchrony, and synchronous growth}
\label{slot:P08-1}
Put $S=g+h>0$ and $M=k+S-c$; the subspace remains denoted by $\Sync$.
\begin{theorem}[MB foundation]
\label{thm:P08-foundation}
For arbitrary canonical $(n,q)$ the following hold.
\begin{enumerate}[label=(\roman*),leftmargin=*]
\item MB has a unique maximal classical solution through every point of
$\R^n$. Every coordinate face is invariant. Nonnegative initial data
stay nonnegative, and initially positive coordinates remain strictly
positive on the maximal interval. No uniform asymptotic lower bound
is meant by this finite-time positivity statement.
\item $\Sync$ is invariant, with exact restriction
\begin{equation}
 \dot s=as-Ms^2.\label{eq:P08-scalar}
\end{equation}
The real synchronous equilibria are $0,a/M$ if $M\ne0$, and only $0$
if $M=0$. A positive synchronous equilibrium exists if and only if
$M>0$, and is uniquely $y^*=a/M$. For $M<0$ the nonzero root is negative.
\item For $s_0>0,M>0$, synchronous dynamics are global,
$0<s(t)\le\max(s_0,a/M)$ and $s(t)\to a/M$. For $M=0$ they are
$s(t)=s_0e^{at}$, global and unbounded. For $M<0$ they blow up at
\begin{equation}
 t_b=\frac1a\log\left(1+\frac{a}{|M|s_0}\right).
 \label{eq:P08-blowup}
\end{equation}
These classify nonnegative synchronous dynamics, not all MB trajectories.
\end{enumerate}
\end{theorem}
\begin{proof}
The polynomial field is smooth and locally Lipschitz on $\R^n$, so
local existence, uniqueness, continuous dependence and maximal
continuation apply. On the coordinate face $y_i=0$, direct substitution
in the original five terms gives $F_i^{\MB}(y)=0$, even when other
coordinates are negative. Along any solution the integrating-factor
identity is
\begin{equation}
 y_i(t)=y_i(0)\exp\int_0^t
 [a-ky_i(r)+cy_{i-1}(r)-gy_{i+q}(r)-hy_{i-q}(r)]\,dr.
 \label{eq:P08-coordinate}
\end{equation}
The integrand is continuous on each compact subinterval of the maximal
interval. The identity holds also for zero initial coordinates and
proves sign preservation and face invariance without assuming global
existence.

At $y=s\one$, each component of the original field is
$as-ks^2+cs^2-gs^2-hs^2=as-Ms^2$. The scalar solution lifts to a
vector solution, uniqueness identifies the two, and continuation makes
their maximal intervals coincide. Zeros of $s(a-Ms)$ give the complete
synchronous equilibrium list, retaining zero before division. For
$s_0>0$, solve $(1/s)'=-a/s+M$ to obtain
\begin{equation}
 s(t)=\frac{s_0e^{at}}{1+(Ms_0/a)(e^{at}-1)}.\label{eq:P08-logistic}
\end{equation}
For $M>0$ the denominator stays positive and the bound and limit follow;
for $M=0$ it is one; for $M<0$ its first zero is exactly
\eqref{eq:P08-blowup}, approached from above with positive numerator.
The zero initial state stays zero in every case.
\end{proof}

\subsection{Sufficient full-system bounds and their scope}
\label{slot:P08-3}
\begin{proposition}[Maximum comparison]
\label{prop:P08-bounds}
For nonnegative MB solutions and $Y=\max_i y_i$,
\begin{equation}
 D^+Y\le aY-(k-c)Y^2.\label{eq:P08-maximum}
\end{equation}
Thus $k\ge c$ suffices for global nonnegative existence. If $k>c$,
put $R=a/(k-c)$. Then $Y(t)\le\max\{Y(0),R\}$, $[0,R]^n$ is
invariant and uniformly attracts bounded subsets, and every
$[0,R+\varepsilon]^n$, $\varepsilon>0$, absorbs bounded subsets in
finite time. MB is therefore dissipative in this sufficient regime.
The comparison alone asserts neither finite-time absorption into the
exact cube $[0,R]^n$ nor boundedness when $k=c$.
\end{proposition}
\begin{proof}
Apply the generic finite-maximum identity of Lemma~\ref{lem:P01-max}.
At any maximizing coordinate $y_i=Y$, all neighbors lie in $[0,Y]$,
so $cy_i y_{i-1}\le cY^2$ and both inhibitory terms are nonpositive.
This includes tied maxima and coincident offsets. For $k\ge c$, the
inequality gives $Y(t)\le Y(0)e^{at}$, excluding finite-time escape;
polynomial continuation yields global existence. For $k>c$, logistic
comparison bounds $Y$ and proves invariance. A logistic solution
starting from the supremum of any bounded set of initial maxima
converges to $R$ and enters $[0,R+\varepsilon]$ in finite time,
proving the uniform attraction and absorption assertions.
\end{proof}
Consequently $M>0$ is necessary for boundedness of every nonnegative
MB trajectory, and $M\ge0$ is necessary for global existence of every
such trajectory, by the synchronous counterexamples. No converse is
inferred. In particular existence of positive synchrony is not used to
assert full-system boundedness. The sufficient condition $k>c$ is not
necessary, as an explicit example below will prove.

\subsection{Homogeneous MB Jacobian and arbitrary-offset spectrum}
\label{slot:P08-4}
At $y=s\one$, differentiate the original field before using equilibrium.
None of $i-1,i+q,i-q$ equals $i$ in the canonical range, so the diagonal is
\[
 d_{\MB}(s)=a-2ks+cs-gs-hs=a-(k+M)s.
\]
The three off-diagonal contributions are $cs,-gs,-hs$. Therefore
\begin{equation}
 DF^{\MB}(s\one)=[a-(k+M)s]I+s(cP-gP^{-q}-hP^q).
 \label{eq:P08-state-J}
\end{equation}
In full, with row $i$ and column $j$, the general derivative entry is
\[
 [a-2ky_i+cy_{i-1}-gy_{i+q}-hy_{i-q}]\delta^{(n)}_{ji}
 +cy_i\delta^{(n)}_{j,i-1}-gy_i\delta^{(n)}_{j,i+q}
 -hy_i\delta^{(n)}_{j,i-q}.
\]
Coincident columns add; none is overwritten. At $y^*=a/M>0$ use
$a=My^*$ only now, giving $d_{\MB}=-ky^*$.

\begin{theorem}[Full MB spectrum, normality, and uniform direction]
\label{thm:P08-spectrum}
For $M>0$, positive MB synchrony has
\begin{align}
 J_{\MB}&=y^*[-kI+cP-gP^{-q}-hP^q],\label{eq:P08-J}\\
 \lambda_m^{\MB}&=y^*[-k+ce^{-\ii\theta_m}-ge^{\ii q\theta_m}
                           -he^{-\ii q\theta_m}],\label{eq:P08-spectrum}\\
 \Ree\lambda_m^{\MB}&=y^*[-k+c\cos\theta_m-S\cos(q\theta_m)],
 \label{eq:P08-real}\\
 \Imm\lambda_m^{\MB}&=y^*[-c\sin\theta_m+(h-g)\sin(q\theta_m)].
 \label{eq:P08-imag}
\end{align}
The $n$ Fourier labels give the full spectrum with algebraic multiplicity;
$J_{\MB}$ is normal and unitarily diagonalized by the project's DFT.
In particular
\begin{equation}
 \boxed{\lambda_0^{\MB}=-a<0.}\label{eq:P08-uniform}
\end{equation}
Any linear instability at positive MB synchrony involves a nonuniform
mode. This does not classify a bifurcation as Hopf or steady.
\end{theorem}
\begin{proof}
The preceding differentiation gives \eqref{eq:P08-J}. Apply the
circulant framework of P04 (Gray, \cite[Theorem~3.1]{Gray2006})
to the coefficients $-k,c,-g,-h$, then take real and imaginary parts.
At label zero the value is $y^*(-k+c-g-h)=-My^*=-a$.
\end{proof}

The complete offset-collision checks, using P03/P04 only for their
index geometry, are
\[
\begin{array}{c|l}
q=1&J_{\MB}/y^*=-kI+(c-h)P-gP^{-1}\\
q=n-1&J_{\MB}/y^*=-kI+(c-g)P-hP^{-1}\\
q=n/2\ (n\text{ even})&J_{\MB}/y^*=-kI+cP-SP^{n/2}.
\end{array}
\]
Their imaginary parts are respectively
$y^*(-c+h-g)\sin\theta_m$, $y^*(-c+g-h)\sin\theta_m$, and
$-y^*c\sin\theta_m$. The last real restraint term is $-y^*S(-1)^m$.
Each also follows from the master spectrum, with cancellation retained.
For even $n$, the singleton is
$\lambda_{n/2}^{\MB}=y^*[-k-c-S(-1)^q]$; the uniform value is
unchanged in every collision case. Label conjugacy and semisimplicity
follow from the real circulant matrix and its unitary diagonalization;
a pair may have zero imaginary part or share an eigenvalue with other
labels. No numerical-eigenvalue simplicity is inferred from label counts.

\begin{corollary}[Exact stability and another sufficient global bound]
\label{cor:P08-stability}
Define
\begin{equation}
 K_{n,q}(c,S)=\max_{0\le m<n}[c\cos\theta_m-S\cos(q\theta_m)].
 \label{eq:P08-K}
\end{equation}
At positive MB synchrony, strict stability is equivalent to $k>K_{n,q}$,
and instability to $k<K_{n,q}$. At equality the normal linear system
is Lyapunov stable but not asymptotically stable; nonlinear stability is
undecided here. Hyperbolicity requires every real part to be nonzero,
not merely a nonzero maximum.

Moreover $k>K_{n,q}$ itself implies $M>0$, global boundedness and
dissipativity for nonnegative MB data. This is another sufficient bound,
not a necessary condition or an identification with $k>c$.
\end{corollary}
\begin{proof}
The abscissa is $y^*(K_{n,q}-k)$ and unitary diagonalization gives
$\|e^{tJ_{\MB}}\|_2=e^{t\sa(J_{\MB})}$, proving the linear claims.
For the bound put $\delta=k-K_{n,q}>0$, $V=\sum_i y_i$. Summing the
original equations yields
\begin{equation}
 V'=aV+y^{\mathsf T}[-kI+cP-gP^{-q}-hP^q]y
 \le aV-\delta\|y\|_2^2\le aV-\frac\delta nV^2.
 \label{eq:P08-sum}
\end{equation}
The symmetric part of the displayed matrix has eigenvalues
$-k+c\cos\theta_m-S\cos(q\theta_m)\le-\delta$, and its skew part
contributes zero to the real quadratic form. Logistic comparison bounds
$V$ before any maximal endpoint and proves global continuation.
The simplex $V\le na/\delta$ attracts bounded subsets, and every
strictly larger one absorbs them. Finally $K_{n,q}\ge c-S$ by its
zero label, so $M>0$.
\end{proof}
For example, for $q=2$ the tuple $(9/2,7,8,1,1)$ has $K_{n,2}=6$
for every $n\ge5$: $2+8u-4u^2$ uniquely maximizes at $u=1$ on the
grid. MB is dissipative although $k=7<c=8$. Its own root is $y^*=9/2$,
whereas DM's unique positive root is $x^*=1$. Direct original residuals
are $9/2-8+7/2=0$ for DM after division by $x^*$, and
$9/2-(9/2)=0$ for MB after division by $y^*$. This refutes necessity
of $k>c$. Conversely, $k>c$ need not imply spectral stability when
$S$ is large, so the sufficient bounds have different scopes.

\subsection{Global asymptotic stability in the positive interior}
\begin{theorem}[Strict MB stability implies global interior stability]
\label{thm:P08-global-interior}
Let $n\ge5$, $1\le q<n$, and $a,k,c,g,h>0$, with the canonical
predecessor shift $(Py)_i=y_{i-1}$. Put $S=g+h$ and assume
\[
 \delta:=k-K_{n,q}(c,S)>0,
\]
where $K_{n,q}$ is exactly \eqref{eq:P08-K}, including label zero.
Then $M=k+S-c>0$, and $\bar y=y^*\one$, $y^*=a/M$, is the unique
equilibrium in $\PosOrth$ and is globally asymptotically stable relative
to $\PosOrth$. Every solution starting in $\PosOrth$ exists for all
forward time, remains in a compact subset of $\PosOrth$ depending on
its initial condition, and converges to $\bar y$.
No solution starting with a zero coordinate can converge to $\bar y$.
No assumption $k>c$ or $\gcd(n,q)=1$ is required.
\end{theorem}
\begin{proof}
We use Goh's generalized Lotka--Volterra Lyapunov method
\cite{Goh1977}, with the following direct verification for MB. Write
\begin{equation}
 \dot y=\diag(y)(a\one+By),\qquad B=-kI+cP-gP^{-q}-hP^q.
 \label{eq:P08-global-field}
\end{equation}
The zero label gives $M\ge\delta>0$, hence
$\bar y=(a/M)\one>0$ and $B\bar y=-a\one$. Define
\begin{equation}
 V(y)=\sum_i[y_i-y^*-y^*\log(y_i/y^*)].
 \label{eq:P08-global-V}
\end{equation}
For $z=y-\bar y$, direct differentiation cancels the population factors:
\begin{equation}
 \dot V=z^{\mathsf T}Bz=z^{\mathsf T}B_{\rm sym}z,
 \quad B_{\rm sym}=-kI+\tfrac c2(P+P^{-1})-\tfrac S2(P^q+P^{-q}).
 \label{eq:P08-global-cancellation}
\end{equation}
P04's Fourier framework gives
$\max\spec(B_{\rm sym})=K_{n,q}-k=-\delta$, so
\begin{equation}
 \dot V\le-\delta\|y-\bar y\|_2^2.
 \label{eq:P08-global-decay}
\end{equation}
Each summand of $V$ is nonnegative, vanishes only at $y_i=y^*$,
and diverges as $y_i\downarrow0$ or $y_i\to\infty$.
Thus sublevels are compact subsets of the positive interior. Positivity
from \eqref{eq:P08-coordinate}, decreasing $V$, and ODE continuation
give global trajectories trapped in those sublevels. The strict Lyapunov
inequality then gives stability and convergence to $\bar y$, and excludes
any other positive equilibrium. Invariant coordinate faces prevent
boundary data from converging to $\bar y$.
\end{proof}


This application does not classify $k=K_{n,q}$ or assert a convergence
rate uniform over initial data or sizes. The DM--MB constructions below
are separate model-specific arguments.

\subsection{Exact step-two MB selected-mode rule}
\label{slot:P08-2}
Set $q=2$ and $u_m=\cos\theta_m$. Since $\cos(2\theta)=2u^2-1$,
\begin{align}
 p(u_m):=\Ree\lambda_m^{\MB}/y^*&=-k+S+cu_m-2Su_m^2,
 \label{eq:P08-quadratic}\\
 p(u)&=-k+S+\frac{c^2}{8S}-2S(u-v_{\MB})^2,\nonumber\\
 v_{\MB}&=\frac{c}{4S}=\frac1{\chi_{\MB}},\qquad
 \chi_{\MB}=\frac{4S}{c}>0.\label{eq:P08-vertex}
\end{align}
This is always strictly concave, $p''=-4S<0$. The unconstrained vertex
is $v_{\MB}$; on $[-1,1]$ the continuous maximizer is
$\min(1,v_{\MB})$. Its departure from $1$ need not switch a discrete mode.

Use P05's notation
\[
 N=\lfloor n/2\rfloor,\quad u_j=\cos(2\pi j/n),\quad
 \beta_j=(u_j+u_{j+1})/2,\quad J=\lfloor(n+1)/4\rfloor,
 \quad\alpha_j=1/\beta_j\ (0\le j<J).
\]
P05 proves $\beta_j=\cos(\pi/n)\cos((2j+1)\pi/n)$ and exactly
$\beta_0,\ldots,\beta_{J-1}>0$. Its nearest-grid-cell argument,
midpoint ordering, parity and label bookkeeping are generic finite-grid
facts. Its DM realization theorem is not an MB realization proof.

\begin{theorem}[MB cells, all ties, reachability, and cell stability]
\label{thm:P08-selection}
On $M>0$, the complete MB representative selection is
\begin{equation}
\begin{array}{c|c}
0<\chi_{\MB}<\alpha_0&\{0\}\\
\alpha_{j-1}<\chi_{\MB}<\alpha_j\ (1\le j<J)&\{j\}\\
\chi_{\MB}>\alpha_{J-1}&\{J\}\\
\chi_{\MB}=\alpha_j\ (0\le j<J)&\{j,j+1\}.
\end{array}\label{eq:P08-cells}
\end{equation}
The middle row is empty for $J=1$. Equivalently, closed vertex cells are
$[\beta_0,\infty)$, $[\beta_j,\beta_{j-1}]$ ($1\le j<N$), and
$(-\infty,\beta_{N-1}]$, intersected with $v_{\MB}>0$.
Every listed strict cell and tie is realizable with $M>0$; the exact
reachable representatives are $0,\ldots,J$.

A nonzero selected representative gives $\{j,n-j\}$, while zero is a
singleton. A tie involving zero gives three labels, other reachable
ties four. No nonadjacent or triple representative maximum occurs.
The even Nyquist singleton is excluded. These are label counts, not
numerical-eigenvalue multiplicities or frequency assertions.
On the cell for $j$ strict stability is exactly
\begin{equation}
 \boxed{k>S+cu_j-2Su_j^2.}\label{eq:P08-cell-stability}
\end{equation}
Either selected representative gives the same bound at a tie.
Reverse strict inequality gives instability and equality gives zero
abscissa. In the uniform cell, including a maximal tie with zero,
$M>0$ already makes the maximum $-a<0$.
\end{theorem}
\begin{proof}
Square completion reduces selection to minimizing $|u_j-v_{\MB}|$.
Apply the generic parts of Theorems~\ref{thm:P05-selection} and
\ref{thm:P05-reachable}. An independent sign check is
\[
 p(u_j)-p(u_{j+1})=c(u_j-u_{j+1})(1-\chi_{\MB}\beta_j),
\]
confirming the switch orientation. For MB-specific realizability,
prescribe any $v>0$, take $g=h=1/2$, $c=4v$,
$k>\max(0,c-1)$ and any $a>0$. Then $M>0$ and $v_{\MB}=v$.
Each required strict cell and positive midpoint can be attained; strict
cells are open by continuity. The grid theorem gives the labels,
exclusions and ties. Taking the maximum of the derived quadratic proves
the exact cell-stability test.
\end{proof}

P05's parity table applies unchanged to the cosine grid. At every finite
positive tuple $v_{\MB}>0$, so the formal zero-vertex tie for
$n=4\ell+2$ is unattainable: its two real parts differ by
$2y^*c\sin(\pi/n)>0$. Only the gap divided by $y^*$ is independent
of $S$. At fixed $n$, sufficiently large finite $\chi_{\MB}$ selects
$J$ uniquely in every parity class.

\subsection{Common domain, named branches, and structural comparison}
\label{slot:P08-5}
For DM set $A=k+h$, $\mu=a+c-g$, $b=a-A$, and $D=b^2+4A\mu$.
By Theorem~\ref{thm:P02-positive}, the full common existence set is
\begin{equation}
\begin{split}
 \mathcal C=\{(a,k,c,g,h)>0:\ &M>0,\quad
 [\mu>0\ \text{or}\ (\mu=0,b>0)\\
 &\text{or}\ (\mu<0,b>0,D\ge0)]\}.
\end{split}\label{eq:P08-common}
\end{equation}
It is not entirely an open set. The branch-labelled comparison pairs
a tuple with one of its allowed positive DM roots
$x_\pm=(b\pm\sqrt D)/(2A)$ and the sole MB root $a/M$.
Most constructions use the open domain
\begin{equation}
 \mathcal C_+=\{(a,k,c,g,h)>0:M>0,\ \mu>0\},\qquad x^*=x_+.
 \label{eq:P08-common-simple}
\end{equation}
DM then has exactly one positive synchronous equilibrium; its scalar
uniform derivative is $-x^*\sqrt D/(1+x^*)<0$ by P04. A separate
construction below explicitly follows $x_-$ on a two-positive-root domain.

At a specified positive DM equilibrium put $s=1+x^*$. Quote P03--P05:
\begin{equation}
\begin{gathered}
 d=-kx^*-c(1+2x^*)/s^2,\quad C=c/s,\quad G=gx^*/s^2,\quad H=hx^*,\\
 r_m^{\DM}=d+Cu_m+(G-H)(2u_m^2-1),\\
 \chi_{\DM}=\frac{4(H-G)}C
 =\frac{4x^*[h(1+x^*)^2-g]}{c(1+x^*)},\qquad
 \chi_{\MB}=\frac{4(g+h)}c.
\end{gathered}\label{eq:P08-DM-input}
\end{equation}
DM's parameter can be negative, zero or positive; its positive-vertex
interpretation applies only if $h>g/s^2$. P05's full selection rule uses
$\chi_{\DM}<\alpha_0$ for every uniform case. At fixed $k,c,g,h$,
MB selection is independent of $a$, whereas DM selection can vary with
$a$ through the equilibrium branch.

\begin{proposition}[Structural functional-form distinctions]
\label{prop:P08-structural}
MB real dispersion for $q=2$ is always strictly concave in $u$.
DM dispersion is strictly convex, linear, or strictly concave according
as $h<g/(1+x^*)^2$, $h=g/(1+x^*)^2$, or $h>g/(1+x^*)^2$.
When $h\le g/(1+x^*)^2$, DM uniquely selects zero. All three
curvature regimes occur in $\mathcal C_+$.

MB has invariant coordinate faces. DM has
$F_i^{\DM}|_{x_i=0}=cx_{i-1}$ and propagates positivity from any nonzero
nonnegative initial state by P01. MB's uniform eigenvalue at positive synchrony is
always $-a$. DM's is $-x^*\sqrt D/s$ on $x_+$,
$+x^*\sqrt D/s$ on a distinct positive $x_-$, and zero at a positive
double root.

DM is globally nonnegative and dissipative for the entire canonical
positive class by P01. MB has the sufficient dissipative regimes above,
but global unbounded synchronous growth at $M=0$ and finite-time
synchronous blow-up at $M<0$. Thus the DM functional form is
structurally dissipative over the entire canonical class, whereas MB is
not. The growth/blow-up examples are outside $\mathcal C$.
\end{proposition}
\begin{proof}
The second derivatives are $-4y^*S<0$ and $4(G-H)$, respectively;
use P05 for uniform selection. To realize the three curvature signs,
set $x^*=1$, $c=g=4$, $k=1$, $a=1+h$, with $h=1/2,1,2$.
Then $\mu=a>0$, $M=1+h>0$, $y^*=1$, $G-H=1-h$, and both
original scalar residuals vanish. The face and global statements quote
P01 and the MB proofs above. The uniform DM values result from inserting
P02's roots in \eqref{eq:P04-equilibrium-zero-mode}. No full-system
boundedness is inferred from MB equilibrium existence.
\end{proof}
\begin{remark}[Meaning and limits of matching]
Matching fixes $n,q$, the five positive coefficients, directed offsets
and coefficient roles, and the original time variable. It does not
calibrate equilibrium amplitude, local derivatives, empirical behavior,
or nonlinear response. The successor derivative is positive for DM and
negative for MB at positive synchrony, despite the negative term in both
fields. No biological or quantitative-superiority claim follows.
For a fixed source $z>0$, as receiver $r\downarrow0$, the predecessor
inputs satisfy $cz/(1+r)\to cz$ for DM and $crz\to0$ for MB.
This last assertion is only a limit under the explicitly fixed-source
convention.
\end{remark}

\subsection{Every ordered dominant-mode pair with both models stable}
\begin{theorem}[Independent realization of the two selection parameters]
\label{thm:P08-pair-realization}
For each fixed $n\ge5,q=2$, every ordered pair
$(j_{\DM},j_{\MB})\in\{0,\ldots,J\}^2$ occurs with unique selected
representatives and both equilibria strictly stable on a nonempty open
subset of $\mathcal C_+$. In particular reverse uniform/nonuniform
selection is possible; distinct nonuniform/nonuniform pairs are possible
when $J\ge2$, beginning at $n=7$. This is not a full stability partition.
\end{theorem}
\begin{proof}
Choose positive parameters in the desired strict cells; explicit choices
for representative $j$ are
\[
 z_0=\alpha_0/2,\quad
 z_j=(\alpha_{j-1}+\alpha_j)/2\ (1\le j<J),\quad
 z_J=2\alpha_{J-1}.
\]
When $J=1$ use only the first and last expressions. Put
$\eta=z_{j_{\DM}}$, $\zeta=z_{j_{\MB}}$, and define
\begin{equation}
\begin{gathered}
 c=4,\quad S=\zeta,\quad x=1+\eta/\zeta,\quad s=1+x,\\
 h=\frac{\zeta+s\eta/x}{s^2+1},\quad g=\zeta-h,\\
 k>\max\{c+S,\ (c-g)/(xs)-h\},\qquad
 a=(k+h)x-(c-g)/s.
\end{gathered}\label{eq:P08-pair-construction}
\end{equation}
Here $xs>\eta/\zeta$, so $0<h<S$ and $g,h>0$. The bound on $k$
gives $a>0$ and $M>0$. Further,
$\mu=x[k+h+(c-g)/s]>0$, since $k>S>g$. Substitution into the
\emph{original} vector fields gives
\begin{equation}
 F_i^{\DM}(x\one)=x[a-(k+h)x+(c-g)/s]=0,\qquad
 F_i^{\MB}(y^*\one)=y^*(a-My^*)=0.
 \label{eq:P08-residual-general}
\end{equation}
Algebra gives $\chi_{\MB}=\zeta$ and
$x(hs^2-g)/s=\eta$, hence $\chi_{\DM}=\eta$.
For all MB labels $\Ree\lambda_m^{\MB}/y^*\le-k+c+S<0$.
For DM the formula gives
\begin{equation}
 \frac{r_m^{\DM}}x
 \le-k-\frac{c}{s^2}+\frac{g}{s^2}+h<-k+S<0.
 \label{eq:P08-DM-bound}
\end{equation}
Indeed the two $c$ contributions after division by $x$ are at most
$-c/s^2$, and $|g/s^2-h|\le g/s^2+h$. This proves all claimed
strict inequalities. Openness in all physical parameters, not just this
displayed family, is proved below.
\end{proof}
A rational reverse-selection witness at $n=8$ is
\begin{equation}
 (a,k,c,g,h)=(113/4,9,8,1,1),\quad x^*=3,\quad y^*=113/12.
 \label{eq:P08-reverse-selection-example}
\end{equation}
Here $M=3$, $\mu>0$, $\chi_{\MB}=1$ selects $0$, while
$\chi_{\DM}=45/8>2\sqrt2=\alpha_1$ selects $2$.
The abscissae are $-443/16$ for DM and $-113/4$ for MB.
Original residuals after division by the respective amplitudes are
$113/4-10\cdot3+7/4=0$ and $113/4-3(113/12)=0$.

\subsection{DM stable and MB unstable for every size}
\begin{theorem}[Universal same-tuple opposite stability]
\label{thm:P08-forward}
For each fixed $n\ge5,q=2$, a nonempty open subset of $\mathcal C_+$
has DM strictly stable with unique dominant representative $0$, while
MB is unstable with a unique nonzero dominant representative. A single
rational tuple works for every dimension:
\begin{equation}
 \boxed{(a,k,c,g,h)=(1/4,1/8,1,1,1/8).}
 \label{eq:P08-forward-witness}
\end{equation}
At this tuple the independently computed amplitudes happen to be
$x^*=y^*=1$; amplitude equality is not a comparison assumption.
\end{theorem}
\begin{proof}
All coefficients are positive, $M=1/4$, $\mu=1/4$. At amplitude one
the original DM residual is $1/4-1/8+1/2-1/2-1/8=0$ and the MB
residual is $1/4-1/8+1-1-1/8=0$. P02 ensures the unique positive DM
root. Its coefficients are $d=-7/8,C=1/2,G=1/4,H=1/8$.
Thus $G-H=1/8>0$, zero is strictly dominant and
$r_0^{\DM}=-1/4$; every DM real part is negative.
For MB,
\[
 S=9/8,\qquad p(u)=1+u-\frac94u^2,\qquad
 v_{\MB}=2/9,\qquad\chi_{\MB}=9/2.
\]
The parity formulas give $-1/4<u_J\le1/2$.
In the only negative case $n\equiv3\pmod4$, $n\ge7$ and
$|u_J|=\sin(\pi/(2n))<\pi/14<1/4$. The other cases give the
upper bound, attained at $n=6$. Hence the actual finite label $J$ satisfies
\[
 p(u_J)\ge1-\frac14-\frac94\left(\frac12\right)^2
 =\frac3{16}>0.
\]
MB is unstable, and its negative uniform eigenvalue excludes zero from
the selected set. There is no representative tie: a tie would force
$\beta_j=2/9$, but
$4\beta_j=2\cos(2\pi j/n)+2\cos(2\pi(j+1)/n)$ is an algebraic
integer, being a sum of roots of unity and their inverses. The rational
number $8/9$ is not an integer, whereas every rational algebraic integer
is an integer (apply the rational-root argument to a monic integer
polynomial). This excludes a tie. Strict gaps and branch regularity
give openness by the argument below.
\end{proof}
For $n=5,6,7$ this witness selects MB representative $1$; for $n=8$
it selects representative $2$. Other sizes follow the exact finite-grid
rule and need not retain either label.
The instability proof uses a finite Fourier label rather than only a
positive continuous maximum.

\subsection{Reverse stability: parity obstruction and constructive regions}
For a positive DM equilibrium write $x=x^*$, $s=1+x$,
$u=\cos\theta_m$, $w=-\cos(2\theta_m)$. The DM formula and
the derived MB formula give
\begin{equation}
 \frac{r_m^{\DM}}x=-k+(h-g/s^2)w
       -\frac{c(1+2x-su)}{xs^2},\qquad
 \frac{r_m^{\MB}}{y^*}=-k+Sw+cu.
 \label{eq:P08-modal-comparison}
\end{equation}
Notice $1+2x-su\ge x>0$ for $u\le1$.
\begin{theorem}[Reverse stability on the unique positive branch]
\label{thm:P08-reverse}
On $\mathcal C_+$ with $q=2$, a tuple with DM unstable and MB strictly
stable exists if and only if $n\equiv3\pmod4$. For every such fixed
size there is a nonempty open region realizing this contrast, with both
models uniquely selecting $J$. For all other sizes, MB strict stability
implies DM strict stability on $\mathcal C_+$.
\end{theorem}
\begin{proof}[Proof of the obstruction]
On $\mathcal C_+$ the uniform DM eigenvalue is strictly negative.
If $G-H\ge0$, P05 makes it largest, so DM is strictly stable.
Otherwise $h-g/s^2>0$. A nonnegative DM real part in
\eqref{eq:P08-modal-comparison} forces $w>0$. If its cosine $u\ge0$,
\begin{equation}
 \frac{r_m^{\MB}}{y^*}-\frac{r_m^{\DM}}x
 =g(1+s^{-2})w+cu+\frac{c(1+2x-su)}{xs^2}>0.
 \label{eq:P08-modal-gap}
\end{equation}
That MB mode is therefore positive. Unless $n\equiv3\pmod4$, all
reachable DM representatives have nonnegative cosines by P05. Any
nonnegative DM maximum would contradict MB strict stability. This also
excludes a zero DM abscissa and proves the strict implication.
\end{proof}
\begin{proof}[Proof of the remaining-size realization]
Let $n\equiv3\pmod4$, $\tau=\sin(\pi/(2n))$, so $0<\tau<1/4$,
$u_J=-\tau$, and $W=1-2\tau^2>7/8$.
Also $u_{J-1}=\sin(3\pi/(2n))=3\tau-4\tau^3$, giving
$\beta_{J-1}=\tau W$. Set
\begin{equation}
\begin{gathered}
 h=1,\quad c=\tau,\quad g=\tau^2/4,\quad
 k=1-5\tau^2/2,\quad x=2/\tau,\quad s=1+x,\\
 a=(k+1)x-(\tau-\tau^2/4)/s.
\end{gathered}\label{eq:P08-reverse-family}
\end{equation}
Every coefficient is positive: $k>27/32$, $x>8$, the subtracted
term in $a$ is less than $1/4$, and $(k+1)x>8$.
Further $M=2-\tau-9\tau^2/4>0$, and $\mu>0$ since $c>g$.
The original residuals vanish by \eqref{eq:P08-residual-general}.

For MB, $v_{\MB}=\tau/(4+\tau^2)<\tau/4<\beta_{J-1}$ and
$\beta_J<0$, so it selects $J$ strictly. Its maximum divided by $y^*$ is
\[
 -k+(1+\tau^2/4)W-\tau^2=-\tau^2(1/4+\tau^2/2)<0.
\]
For DM the $J$-mode divided by $x$ is
\[
 \frac{r_J^{\DM}}x=\frac{\tau^2}{2}-\frac{gW}{s^2}
              -\frac{\tau(1+2x+s\tau)}{xs^2}.
\]
Since $1+2x+s\tau\le(2+\tau)s$, $s>x$, and $x=2/\tau$,
the two subtracted terms are less than
\[
 \frac{g+\tau(2+\tau)}{x^2}<\frac{37}{256}\tau^2.
\]
Here $g<1/64$ and $\tau(2+\tau)<9/16$. Thus
$r_J^{\DM}/x>91\tau^2/256>0$.
DM is concave, and its vertex obeys
\[
 v_{\DM}=\frac{\tau}{4x(s-g/s)}<\frac{\tau}{256}
      <\tau W=\beta_{J-1},
\]
since $x>8$ and $s-g/s>8$. So DM also strictly selects $J$.
These are strict canonical common-domain inequalities; openness is
established below.
\end{proof}

\paragraph{Rational reverse-stability witness.}
At $n=7$ a compact exact witness is
\begin{equation}
 (a,k,c,g,h)=(20301/1100,17/20,1/2,1/100,1),\quad
 x^*=10,\quad M=34/25,\quad y^*=20301/1496.
 \label{eq:P08-reverse-rational}
\end{equation}
Original residuals divided by their amplitudes are
$20301/1100-(37/20)10+(49/100)/11=0$ for DM and
$20301/1100-(34/25)(20301/1496)=0$ for MB.
All coefficients and $\mu$ are positive.
The ordered seven-node cosines $u_1,u_2,u_3$ are the three roots of
$8u^3+4u^2-4u-1$: use
$1+2(\cos\theta+\cos2\theta+\cos3\theta)=0$ at nonzero seventh roots
and the Chebyshev identities. Rational isolating intervals are
\begin{equation}
 31/50<u_1<63/100,\quad -23/100<u_2<-11/50,\quad
 -91/100<u_3<-9/10.\label{eq:P08-seven-intervals}
\end{equation}
The endpoint signs of the cubic are respectively $(-,+),(+,-),(-,+)$.
Three disjoint sign changes exhaust its degree; cosine ordering assigns
the labels. Thus no decimal trigonometric value is a proof premise.
MB representatives $1,2,3$ have normalized real parts
$-17/20+u_1/2-101u_2/100$,
$-17/20+u_2/2-101u_3/100$, and
$-17/20+u_3/2-101u_1/100$.
The interval upper bounds are respectively $-3027/10000$,
$-409/10000$, $-19262/10000$, all negative; mode zero is $-34/25$.
DM instead satisfies
\[
 \frac{r_2^{\DM}}{10}
 >-\frac{17}{20}-\frac{21}{2420}-\frac{23}{22000}
      +\frac{12099}{12100}\frac9{10}
 =\frac{9729}{242000}>\frac1{25}>0.
\]
Both select $2$ strictly: $v_{\MB}=25/202<39/200<\beta_1$,
$v_{\DM}=55/48396<39/200<\beta_1$, and $\beta_2<0$.

\paragraph{The full domain requires explicit branch labels.}
The preceding obstruction cannot be extended to all of $\mathcal C$
without qualification. A lower-branch construction gives the exact
existence answer on that larger branch-labelled domain.
\begin{proposition}[Reverse contrast when both DM branches are allowed]
\label{prop:P08-lower}
If $4\mid n$ and MB is strictly stable, every positive DM synchronous
equilibrium is strictly stable. If $4\nmid n$, a nonempty open common
region has MB strictly stable and the specified lower positive DM
branch unstable with unique dominant representative zero.
\end{proposition}
\begin{proof}
When $4\mid n$, MB's label $n/4$ has normalized real part $-k+S$.
Its strict stability forces $k>S$. Bound \eqref{eq:P08-DM-bound},
which did not require $\mu>0$, makes every DM real part negative.

Otherwise put $\rho_n=\max_m[-\cos(2\theta_m)]$. One has
$0<\rho_n<1$: equality to one requires $4m=n$ modulo $2n$, impossible
when $4\nmid n$, and a label nearest $n/4$ has negative double-angle
cosine. Define
\begin{equation}
\begin{gathered}
 \varepsilon=(1-\rho_n)/24,\quad k=1,\quad c=h=\varepsilon,\quad
 g=(1+\varepsilon)^3+2\varepsilon,\\
 a=(1+\varepsilon)(1+2\varepsilon)
                         +\frac{\varepsilon}{1+\varepsilon}.
\end{gathered}\label{eq:P08-lower-family}
\end{equation}
All coefficients are positive and $M=1+g>0$. With $A=1+\varepsilon$,
the two DM roots are exactly
\[
 x_-=\varepsilon,\qquad
 x_+=\varepsilon+\frac{\varepsilon}{(1+\varepsilon)^2}.
\]
Direct substitution of $x_-$ in $a-Ax+(c-g)/(1+x)$ gives zero.
The second is checked by the quadratic factorization
$A(x-x_-)(x-x_+)$ of P02's equilibrium polynomial, or by substitution.
In particular $\mu=-A\varepsilon^2-\varepsilon^2/(1+\varepsilon)<0$,
$b>0$, $D>0$, and both roots are positive. On the named lower branch,
\[
 \lambda_0^{\DM}=\frac{\varepsilon^2}{(1+\varepsilon)^2}>0,
 \qquad G-H=\varepsilon\left(1+
                       \frac{2\varepsilon}{(1+\varepsilon)^2}\right)>0.
\]
Thus zero is uniquely dominant and unstable. For MB,
$S=(1+\varepsilon)^3+3\varepsilon<1+10\varepsilon$, since
$0<\varepsilon<1$, and every normalized real part obeys
\[
 -1+\varepsilon\cos\theta_m-S\cos(2\theta_m)
 \le-1+\varepsilon+S\rho_n
 <-\frac{13}{24}(1-\rho_n)<0.
\]
The MB original residual vanishes at $y^*=a/M$. Distinct simple
positive DM roots and all strict inequalities persist locally on the
specified lower branch. At $n=6$, $\rho_6=1/2$, so
$\varepsilon=1/48$ gives an entirely rational witness.
\end{proof}

\subsection{Mode/stability combinations and physical-parameter openness}
All five requested types of combination have realizations, without a
claim that each occurs in every dimension. Both stable with identical
or different modes follows from Theorem~\ref{thm:P08-pair-realization};
DM stable / MB unstable is universal by Theorem~\ref{thm:P08-forward};
the reverse has the branch and size restrictions just proved.
For both unstable with different dominant modes take $n=8$ and
\begin{equation}
 (a,k,c,g,h)=(16351/1100,1/2,3/2,1/100,1),\quad
 x^*=10,\quad M=1/100,\quad y^*=16351/11.
 \label{eq:P08-both-unstable}
\end{equation}
The original residuals divided by their amplitudes are
$16351/1100-15+(149/100)/11=0$ and
$16351/1100-(1/100)(16351/11)=0$. Coefficients and $\mu$ are positive.
MB's vertex $75/202$ lies strictly between $\sqrt2/4$ and
$(2+\sqrt2)/4$, so MB selects $1$ with normalized maximum
$-1/2+3\sqrt2/4>0$. DM's vertex $55/16132<\sqrt2/4$ selects $2$,
whose normalized maximum is
\[
 -\frac12-\frac{63}{2420}+\frac{12099}{12100}
 =\frac{2867}{6050}>0.
\]
The small value of $M$ is strictly positive. MB's large amplitude is
its own computed equilibrium, not an amplitude imported from DM.

\begin{proposition}[Openness in all five physical coefficients]
\label{prop:P08-open}
Every strict contrast and strict mode pair constructed here persists
on a nonempty open set in $(a,k,c,g,h)\in(0,\infty)^5$ for its stated
fixed size, offset, and DM branch. No openness assertion rests merely
on an interval inside a displayed lower-dimensional family.
\end{proposition}
\begin{proof}
On $\mathcal C_+$, $D>0$ and the unique positive root $x_+$ is smooth
in all five coefficients. Equivalently, the scalar equilibrium function
$B(x)=a-(k+h)x+(c-g)/(1+x)$ has
$B_x=-\sqrt D/(1+x_+)<0$, so the implicit function theorem applies.
MB's branch $a/M$ is smooth for $M>0$. Coefficient positivity,
$M>0$, and $\mu>0$ are strict open conditions. At fixed $n$ there
are finitely many spectral inequalities; each is strict and continuous.
A unique selected representative has a strictly positive gap to every
other representative, while its conjugate labels remain tied.
A sufficiently small intersection of these neighborhoods preserves all
branch, sign, stability and selection assertions.

For the lower-branch family use instead $\mu<0,b>0,D>0$.
Both positive roots are simple and the specified $x_-$ is smooth, with
$B_x=+\sqrt D/(1+x_-)>0$. The same continuity argument preserves
that branch and its strict conclusions. No full-system IFT, later
module, uniform neighborhood radius in $n$, or global branch
identification is used.
\end{proof}

\subsection{Global strengthening of the strict comparison regimes}
\begin{corollary}[Reverse contrast with globally attracting MB]
\label{cor:P08-global-contrast}
Every reverse witness in Theorem~\ref{thm:P08-reverse} and
Proposition~\ref{prop:P08-lower}, and every strict reverse region
preserved by Proposition~\ref{prop:P08-open}, has the following stronger
conclusion: the specified positive DM homogeneous equilibrium is
unstable, while MB's positive equilibrium is globally asymptotically
stable in $\PosOrth$. Likewise the MB equilibrium in every stable
mode-pair realization of Theorem~\ref{thm:P08-pair-realization} is
globally asymptotically stable in $\PosOrth$.
\end{corollary}
\begin{proof}
Each cited construction has strictly stable MB, equivalently
$\delta>0$ by Corollary~\ref{cor:P08-stability}; apply
Theorem~\ref{thm:P08-global-interior}. The precise margins can also
be checked directly. In \eqref{eq:P08-reverse-family}, the maximizing
representative is $J$ and $\cos\theta_J=-\tau$, so
\[
 K_{n,2}=(1+\tau^2/4)(1-2\tau^2)-\tau^2,\qquad
 \delta=\tau^2/4+\tau^4/2>0.
\]
The DM estimate already proved gives $r_J^{\DM}/x>91\tau^2/256>0$.
For \eqref{eq:P08-lower-family}, $4\nmid n$ and
$K_{n,2}\le\varepsilon+S\rho_n$, whence
\[
 \delta>\frac{13}{24}(1-\rho_n)>0,\qquad
 \lambda_0^{\DM}=\frac{\varepsilon^2}{(1+\varepsilon)^2}>0.
\]
Both are the original tuples, with their original DM branch labels.
\end{proof}

\paragraph{Explicit certificates and exclusions.}
For the rational $n=7$ tuple \eqref{eq:P08-reverse-rational},
\[
 K_{7,2}=\tfrac12\cos(4\pi/7)-\tfrac{101}{100}\cos(8\pi/7),
 \qquad \delta>409/10000,\qquad r_2^{\DM}>9729/24200.
\]
These strict bounds follow from the rational isolating intervals
\eqref{eq:P08-seven-intervals}; decimal trigonometry is not a proof premise.
At $n=6$, the lower-branch construction gives the exact tuple
\[
 (a,k,c,g,h)=(61177/56448,1,1/48,122257/110592,1/48),
\]
with $x_-=1/48$, $y^*=5872992/11409601$,
$K_{6,2}=126865/221184$, $\delta=94319/221184$, and
$\lambda_0^{\DM}=1/2401>0$.
The universal witness \eqref{eq:P08-forward-witness} is instead the
\emph{forward} contrast: MB is unstable, with $K_{n,2}-k\ge3/16$.
It does not satisfy the global theorem. Nor does the both-unstable
tuple \eqref{eq:P08-both-unstable}, for which
$\delta=1/2-3\sqrt2/4<0$.

The exact residue-class restrictions retain their original strict-MB
qualification: on $\mathcal C_+$ the reverse contrast exists exactly
when $n\equiv3\pmod4$, and allowing the specified lower branch gives
existence exactly when $4\nmid n$. These statements are not promoted
to an unrestricted classification allowing $k=K_{n,q}$.
P12 remains an unproved transversality module and is not used here.

\subsection{Interpretation, deferred gates, and internal checks}
\label{slot:P08-6}
The main comparison results are a universal forward stability
contrast, both stability directions at $n=7$ on the unique-positive-root
domain, every ordered pair of reachable modes with both models stable,
and exact parity/branch limitations on reverse stability. Identical
abstract reachable labels do not imply identical behavior at the same
physical tuple. The unrestricted reverse-stability hope fails on
$\mathcal C_+$ unless $n\equiv3\pmod4$; even allowing both DM
branches cannot produce it when $4\mid n$.
The rejected phrase ``MB is generically unstable'' is not reinstated:
MB has explicit open stable regions and an exact stability criterion.

For structural growth examples, $(a,k,c,g,h)=(1,1,3,1,1)$ gives
$M=0$, and $(1,1,4,1,1)$ gives $M=-1$. At $s_0=1$ the latter blows
up at $\log2$. DM has its P01 bounded flow for both. Its own positive
root is $1$ or $(\sqrt{33}-1)/4$, respectively, each satisfying the
original residual $x[a-(k+h)x+(c-g)/(1+x)]=0$.
MB has no positive synchronous root in either case; its origin has zero
residual, and the stated scalar trajectories have the exact original
MB derivative. These are not common-equilibrium stability comparisons.

\label{slot:P08-extra-1}
The six-term model is explicitly deferred. Existing source-provenance
records are retained, but no archived six-term equation is imported.
Full MB Hopf, first Lyapunov coefficient, crossing transversality, and
criticality comparisons are deferred as well. The spectra and open
regions are inputs for later P11--P13, not proofs of their hypotheses.
A selected nonuniform pair, positive real part, or candidate imaginary
crossing is insufficient for Hopf. 


% EDITION: citation-based shortening, 4 October 2026.
% Earlier audits refer to the source edition; this revision has not been independently audited.
\section{P09: Transcritical bifurcation at the origin}
\label{sec:P09}

Fix canonical $a,k,c,h>0$, $n\ge5$, and $1\le q<n$, and vary $g$.
Throughout this module put
\begin{equation}
 A=k+h>0,\qquad b=a-A,\qquad g_c=a+c>0,\qquad \mu=g_c-g.
 \label{eq:P09-parameters}
\end{equation}
Here $b$ is a local abbreviation, not an archived comparator coefficient.
Write $F(x,\mu)$ for the original canonical DM field with $g=g_c-\mu$:
\begin{equation}
 F_i(x,\mu)=a x_i-kx_i^2+\frac{c x_{i-1}}{1+x_i}
 -\frac{(g_c-\mu)x_i}{1+x_{i+q}}-h x_i x_{i-q}.
 \label{eq:P09-field}
\end{equation}
For $|\mu|<g_c/2$ all coefficients remain positive. The field is jointly
$C^\infty$ on $\Domain\times(-g_c/2,g_c/2)$,
$\Domain=(-1,\infty)^n$, by its nonvanishing denominators and P01's
local foundation (Proposition~\ref{prop:P01-local}). No global invariance
of $\Domain$ is needed. Neighborhoods may depend on all fixed parameters
and on $n$; no uniform-in-$n$ assertion is made.

The direct proof inputs are P02, P03 and P04. P02's
Proposition~\ref{prop:P02-invariance} and \eqref{eq:P02-scalar} supply
\begin{equation}
 F(r\one,\mu)=f(r,\mu)\one,\qquad
 f(r,\mu)=\frac{r}{1+r}(\mu+br-Ar^2),\qquad r>-1.
 \label{eq:P09-scalar}
\end{equation}
This is the imported exact invariant restriction, not a new global
classification.

\subsection{Critical spectrum and full local equilibrium reduction}
\label{slot:P09-1}

\begin{lemma}[Simple synchronous kernel and transverse gap]
\label{lem:P09-gap}
At $(x,\mu)=(0,0)$,
\begin{equation}
 J_c=D_xF(0,0)=c(P-I),\qquad
 \lambda_m^c=c(e^{-\ii\theta_m}-1),\quad \theta_m=\frac{2\pi m}{n}.
 \label{eq:P09-critical-spectrum}
\end{equation}
Zero is algebraically simple, with real eigenspace
$\operatorname{span}\{\one\}$. Every transverse eigenvalue satisfies
\begin{equation}
 \Ree\lambda_m^c\le-\delta_n\quad(1\le m<n),\qquad
 \delta_n=c\left(1-\cos\frac{2\pi}{n}\right)>0.
 \label{eq:P09-gap}
\end{equation}
No condition on $q$ or $\gcd(n,q)$ is imposed beyond the canonical range.
\end{lemma}
\begin{proof}
P03 gives $D_xF(0,\mu)=\mu I+c(P-I)$ and its simple origin
spectrum. For $m\ne0$, $\cos\theta_m\le\cos(2\pi/n)<1$,
which gives the stated gap, independently of $q$.
\end{proof}

\begin{theorem}[Every sufficiently small equilibrium is synchronous]
\label{thm:P09-reduction}
For every fixed canonical tuple, including $b=0$, there are state and
parameter neighborhoods of zero in which the entire full-system
equilibrium equation is equivalent to
\begin{equation}
 x=r\one+y,\qquad y=0,\qquad r(\mu+br-Ar^2)=0.
 \label{eq:P09-complete-equilibrium}
\end{equation}
In particular, no nearby nonsynchronous equilibrium branch is omitted.
\end{theorem}
\begin{proof}
Set $w^{\mathsf T}=n^{-1}\one^{\mathsf T}$,
$E=\ker w^{\mathsf T}$, $\Pi_\perp=I-\one w^{\mathsf T}$,
and $x=r\one+y$, $y\in E$. The transverse equation
$T(r,y,\mu)=\Pi_\perp F(r\one+y,\mu)=0$ has derivative
$D_yT(0,0,0)=J_c|_E$, invertible by Lemma~\ref{lem:P09-gap}.
The implicit-function theorem \cite[Ch.~1]{GuckenheimerHolmes1983}
gives a unique nearby solution $y=\psi(r,\mu)$. Synchrony already
gives $T(r,0,\mu)=0$, so uniqueness implies
\begin{equation}\psi(r,\mu)\equiv0.\label{eq:P09-psi-zero}\end{equation}
The remaining equation is $f(r,\mu)=0$; since $1+r>0$ it is
exactly \eqref{eq:P09-complete-equilibrium}, including when $b=0$.
\end{proof}

\subsection{Local derivatives and full-vector nondegeneracy}

\begin{proposition}[Transcritical coefficients and parameter orientation]
\label{prop:P09-coefficients}
The local scalar derivatives are
\begin{equation}
 f(0,\mu)=0,\quad f_r(0,0)=0,\quad
 f_{r\mu}(0,0)=1,\quad f_{rr}(0,0)=2b.
 \label{eq:P09-scalar-coefficients}
\end{equation}
Choose $v=\one$, $w^{\mathsf T}=n^{-1}\one^{\mathsf T}$, so
$J_cv=0$, $w^{\mathsf T}J_c=0$, and $w^{\mathsf T}v=1$.
The full vector field satisfies
\begin{equation}
 F(0,\mu)=0,\quad w^{\mathsf T}F_{x\mu}(0,0)v=1,\quad
 w^{\mathsf T}F_{xx}(0,0)[v,v]=2b.
 \label{eq:P09-full-coefficients}
\end{equation}
The quadratic Taylor coefficient includes the factor $1/2$ and is $b$.
With $g$ as parameter the mixed coefficient is instead
$w^{\mathsf T}F_{xg}(0,g_c)v=-1$. Ordinary transcritical
nondegeneracy holds exactly when $b\ne0$.
\end{proposition}
\begin{proof}
Equation~\eqref{eq:P09-scalar} gives $f_r(0,\mu)=\mu$ and
$f_\mu(r,\mu)=r/(1+r)$, whose $r$ derivative at zero is one.
At $\mu=0$, $f(r,0)=br^2-(A+b)r^3+O(r^4)$, so $f_{rr}=2b$.

Independently, every term of the original field vanishes at $x=0$.
Its parameter derivative is $(F_\mu)_i=x_i/(1+x_{i+q})$,
whose state derivative at zero is the identity. Using $g$ instead
gives $-x_i/(1+x_{i+q})$, hence minus the identity.
The second directional derivatives of the five original terms along $v$
are respectively $0,-2k,-2c,2g_c,-2h$. Their sum is
$2(g_c-c-k-h)=2b$ in each component. Averaging by $w$ proves the
full-vector identities without a missing factor of $n$. Also
$F_\mu(0,0)=0$, as required by persistence of the origin branch.
\end{proof}

\subsection{Generic full-system transcritical theorem}
\label{slot:P09-2}

\begin{theorem}[Nondegenerate transcritical bifurcation at $g_c=a+c$]
\label{thm:P09-transcritical}
Suppose $b=a-k-h\ne0$. For sufficiently small $|\mu|$, all full-system
equilibria sufficiently near the origin are exactly
\begin{equation}
 x=0\quad\hbox{and}\quad x=r_*(\mu)\one,\qquad
 r_*(0)=0,\quad r_*(\mu)\ne0\ (\mu\ne0),
 \label{eq:P09-branches}
\end{equation}
where $r_*$ is smooth and
\begin{equation}
 r_*'(0)=-\frac1b,\qquad
 r_*(\mu)=-\frac\mu b+\frac{A\mu^2}{b^3}+O(\mu^3),\qquad
 \frac{dr_*}{dg}(g_c)=\frac1b.
 \label{eq:P09-slopes}
\end{equation}
The two curves have distinct tangent directions in the synchronous
state--parameter plane. The critical eigenspace is one-dimensional and
synchronous; all transverse modes are strictly stable. In $\Domain$,
the full nonlinear stability exchange is
\begin{center}
\begin{tabular}{@{}lll@{}}
\toprule
Side & Origin & Nearby nonzero branch\\
\midrule
$\mu<0$ ($g>g_c$) & Locally exponentially stable & Unstable\\
$\mu>0$ ($g<g_c$) & Unstable & Locally exponentially stable\\
\bottomrule
\end{tabular}
\end{center}
Together with the full-vector nondegeneracy coefficients, these facts
give a nondegenerate full-system local transcritical bifurcation.
\end{theorem}
\begin{proof}
Theorem~\ref{thm:P09-reduction} first makes the equilibrium problem
complete in the full state space. Apart from the persistent origin,
solve $Q(r,\mu)=Ar^2-br-\mu=0$. Since $Q_r(0,0)=-b\ne0$,
the implicit-function theorem gives a unique smooth root $r_*(\mu)$
in a product neighborhood. The identity $Q(0,\mu)=-\mu$ shows that
this root is nonzero exactly for $\mu\ne0$. This proves the full local
list. For fixed $b\ne0$, the other root $b/A$ at $\mu=0$ is excluded
by shrinking the state neighborhood.

Differentiating $Q(r_*(\mu),\mu)=0$ gives
$(2Ar_*-b)r_*'=1$, hence $r_*'(0)=-1/b$.
Differentiating again gives
$2A(r_*')^2+(2Ar_*-b)r_*''=0$, hence $r_*''(0)=2A/b^3$.
Smoothness proves the expansion. In $(r,\mu)$ coordinates the tangents
are $(0,1)$ and $(-1/b,1)$, which are independent. This is transverse
crossing in that plane, not a claim that two curves span the full
$(n+1)$-dimensional ambient space. Since $d\mu/dg=-1$, the physical
$g$ slope has the opposite sign. Increasing $g$ decreases $\mu$.

The standard transcritical and center-manifold framework is \cite[Ch.~3]{GuckenheimerHolmes1983}. Here the center dynamics is explicit:
At $\mu=0$ the synchronous line is smooth, invariant, and tangent to
the entire center eigenspace, so is a local one-dimensional center
manifold. For the extended system
$(\dot x,\dot\mu)=(F(x,\mu),0)$, its product with the parameter
interval is an exact two-dimensional center manifold:
$F_\mu(0,0)=0$, the extended linearization is
$\operatorname{diag}(J_c,0)$, and all remaining eigenvalues are strictly
stable. On this product use
\[
 u=Ar^2-br,\qquad \frac{du}{dr}(0)=-b\ne0.
\]
This is a local smooth coordinate. Direct substitution yields
\begin{equation}
 \dot u=\alpha(r)u(\mu-u),\qquad
 \alpha(r)=\frac{2Ar-b}{(Ar-b)(1+r)},\qquad\alpha(0)=1.
 \label{eq:P09-exact-normal-form}
\end{equation}
Near zero the denominators do not vanish and $\alpha>0$. Thus
$d\tau/dt=\alpha(r)$ gives exactly $du/d\tau=u(\mu-u)$, the
nondegenerate transcritical center dynamics. This is orbital equivalence
of the exact center restriction, not exact decoupling of the full field.

Full stability still requires all Fourier modes. Differentiating first,
\[
 f_r(r,\mu)=\frac{\mu+br-Ar^2}{(1+r)^2}
                  +\frac r{1+r}(b-2Ar).
\]
At the nonzero equilibrium the first numerator vanishes. By P04's
\eqref{eq:P04-synchronous-derivative}, its uniform eigenvalue is
\begin{equation}
 \lambda_0^{\rm nz}=\frac{r_*}{1+r_*}(b-2Ar_*)
 =-\mu\frac{b-2Ar_*}{(1+r_*)(b-Ar_*)}.
 \label{eq:P09-uniform-nonzero}
\end{equation}
The last equality uses $\mu=-r_*(b-Ar_*)$. Choose $|r_*|<1/2$ and
$2A|r_*|<|b|/2$. Both $b-Ar_*$ and $b-2Ar_*$ have the sign of
$b$, so the final fraction is positive. Therefore, exactly,
\begin{equation}
 \operatorname{sgn}\lambda_0^{\rm nz}=-\operatorname{sgn}\mu,
 \qquad\lambda_0^{\rm nz}=-\mu+O(\mu^2).
 \label{eq:P09-uniform-sign}
\end{equation}
The sign argument, rather than the expansion, gives stability.

For the transverse modes use the general P04 coefficients, valid also
at zero or negative $r>-1$:
\begin{align}
 d(r,\mu)&=a-2kr-\frac{cr}{(1+r)^2}-\frac{g_c-\mu}{1+r}-hr,
 \notag\\
 C(r)&=\frac c{1+r},\quad G(r,\mu)=\frac{(g_c-\mu)r}{(1+r)^2},
 \quad H(r)=hr,\notag\\
 \lambda_m(r,\mu)&=d+Ce^{-\ii\theta_m}+Ge^{\ii q\theta_m}
                                      -He^{-\ii q\theta_m}.
 \label{eq:P09-continuous-spectrum}
\end{align}
Every expression is continuous at $(r,\mu)=(0,0)$, where it equals
$\lambda_m^c$. Finitely many labels allow one common neighborhood
with $|\lambda_m(r,\mu)-\lambda_m^c|<\delta_n/2$ for every $m\ne0$.
Thus for some $\eta>0$,
\begin{equation}
 |r|+|\mu|<\eta\quad\Longrightarrow\quad
 \Ree\lambda_m(r,\mu)<-\delta_n/2\quad(m\ne0).
 \label{eq:P09-persistent-gap}
\end{equation}
No mode-selection theorem is used.

At the origin $\lambda_0=\mu$, and every other real part is
$\mu-c(1-\cos\theta_m)<-\delta_n/2$ if $|\mu|<\delta_n/2$.
On the nonzero branch use \eqref{eq:P09-uniform-sign} and
\eqref{eq:P09-persistent-gap}. A Hurwitz Jacobian gives local exponential
stability of the smooth nonlinear equilibrium; a positive eigenvalue
gives nonlinear instability, as in P04's nonlinear implications.
Every unstable equilibrium here has exactly one unstable real direction.
This proves the table and the full-system theorem, without using P05,
P07 or P08.
\end{proof}

\begin{corollary}[Physical positive-branch orientation]
\label{cor:P09-physical}
For sufficiently small nonzero $\mu$, $r_*(\mu)>0$ exactly when
$\operatorname{sgn}\mu=-\operatorname{sgn}b$.
If $a<k+h$, this branch is positive and locally exponentially stable
for $\mu>0$, equivalently $g<g_c$, while the origin is unstable.
If $a>k+h$, it is positive and unstable for $\mu<0$, equivalently
$g>g_c$, while the origin is locally exponentially stable.
Origin instability for $\mu>0$ is also visible relative to $\Orth$;
for $\mu<0$ exponential stability restricts to that state space.
\end{corollary}
\begin{proof}
Use the exact relation $\mu=-r_*(b-Ar_*)$ and the sign of $b-Ar_*$,
then the full stability table. The unstable origin eigenvector
$v=\one$ is strictly positive. Explicitly, at any fixed small $\mu>0$,
$f(r,\mu)/r\to\mu$ as $r\downarrow0$. Choose $\varepsilon>0$ such
that $f(r,\mu)\ge(\mu/2)r$ on $0<r\le\varepsilon$.
Arbitrarily small positive synchronous states reach $r=\varepsilon$
in finite time. Their full-system lifts remain nonnegative, proving
relative instability. For $\mu<0$, full local exponential stability
and P01's nonnegative invariance give the restriction. An unstable
positive branch point is interior to $\Orth$; small positive synchronous
perturbations also show its relative instability. P09's local branch
is not any additional positive root away from zero allowed by P02
when $b>0$.
\end{proof}

\begin{proposition}[Generic critical origin in the open domain]
\label{prop:P09-generic-critical}
For $\mu=0$, $b\ne0$, the origin is nonlinearly unstable in $\Domain$.
For $b>0$ this is visible from nonnegative synchronous states.
No full relative-orthant stability conclusion is asserted at
$\mu=0$, $b<0$.
\end{proposition}
\begin{proof}
The exact invariant-line dynamics is
\begin{equation}
 \dot r=\frac{r^2(b-Ar)}{1+r}.\label{eq:P09-generic-critical}
\end{equation}
Choose $0<\varepsilon<1/2$ with $A\varepsilon<|b|/2$.
For $0<|r|\le\varepsilon$ the sign is that of $b$, and the magnitude
is at least $|b|r^2/3$. For $b>0$ take $r(0)>0$; for $b<0$ take
$r(0)<0$. On that side $s=|r|$ satisfies
$\dot s\ge(|b|/3)s^2$ until $s=\varepsilon$, which must occur in
finite time. Arbitrarily small data therefore exit a fixed neighborhood.
The exact synchronous lift stays in $\Domain$ up to that exit.
When $b<0$ these initial states are negative, so this proof does not
establish relative-orthant instability.
\end{proof}

\subsection{The cubic-degenerate case and its precise scope}
\label{slot:P09-3}

\begin{proposition}[Degenerate geometry and full nonlinear stability]
\label{prop:P09-degenerate}
Let $b=0$, that is $a=A=k+h$. Ordinary transcritical quadratic
nondegeneracy fails. All sufficiently small full-system equilibria are
synchronous, with the following exact local geometry and stability.
\begin{enumerate}[label=(\roman*),leftmargin=*]
\item For $\mu<0$, only the origin exists locally; it is locally exponentially stable.
\item At $\mu=0$, only the origin exists locally. It is a triple zero
of the scalar field and is locally asymptotically stable in the full
open domain $\Domain$, but not locally exponentially stable.
\item For $\mu>0$, there are precisely the origin and
\begin{equation}
 x=r_\pm\one,\qquad r_\pm=\pm\sqrt{\mu/A}.
 \label{eq:P09-square-root}
\end{equation}
For sufficiently small $\mu>0$, both nonzero equilibria belong to
$\Domain$ and are locally exponentially stable there. The origin is
unstable, also relative to $\Orth$. Only $r_+\one$ is physically
nonnegative.
\end{enumerate}
The two nonzero local roots occur for $g<g_c$; for $g>g_c$ only the
stable origin remains locally. This is a cubic-degenerate origin with
square-root nonzero branches. No pitchfork classification or
two-parameter unfolding is asserted.
\end{proposition}
\begin{proof}
Theorem~\ref{thm:P09-reduction} does not require $b\ne0$. Now
$f(r,\mu)=r(\mu-Ar^2)/(1+r)$, and the nonzero equation $Ar^2=\mu$
gives the complete list. For $0<\mu<A$, the negative root exceeds
$-1$; shrink the interval further to put both roots within the
reduction and spectral-gap neighborhoods. At $\mu=0$,
\begin{equation}
 f(r,0)=-\frac{Ar^3}{1+r},\qquad f_{rrr}(0,0)=-6A\ne0.
 \label{eq:P09-cubic}
\end{equation}
The scalar zero is therefore triple. At either nonzero root,
\begin{equation}
 \lambda_0=f_r(r_\pm,\mu)=-\frac{2Ar_\pm^2}{1+r_\pm}<0.
 \label{eq:P09-degenerate-eigenvalue}
\end{equation}
The proof of \eqref{eq:P09-persistent-gap} never used $b\ne0$;
it supplies strict stability of all other Fourier modes. Thus both
full equilibria are locally exponentially stable. At the origin
$\lambda_0=\mu$ gives the off-critical stability; the positive
synchronous escape argument above proves relative instability for
$\mu>0$.

At the fixed value $\mu=0$, the smooth invariant synchronous line
passes through zero and is tangent to the whole one-dimensional center
subspace. All complementary eigenvalues have real part at most
$-\delta_n$, and there is no unstable subspace. On this exact center
manifold, $r^2/2$ has derivative $-Ar^4/(1+r)<0$ for small $r\ne0$,
so its origin is locally asymptotically stable. These verify the
hypotheses of the finite-dimensional center-manifold stability
reduction principle: a smooth local system with center and strictly
stable spectra inherits local asymptotic stability from its center
restriction. No parameter direction is included in this fixed-parameter
stability statement.

This applies the center-manifold stability reduction principle of \cite[Ch.~3]{GuckenheimerHolmes1983}; the displayed gap and exact invariant restriction verify its hypotheses.

Convergence is not exponential. Each sufficiently small nonzero
synchronous solution keeps its sign, tends to zero, and obeys
\[
 \frac{d}{dt}(r^{-2})=\frac{2A}{1+r(t)}\longrightarrow2A.
\]
It follows that $|r(t)|\sim(2At)^{-1/2}$ as $t\to\infty$.
These exact full-system solutions exclude a local exponential stability
estimate. No common asymptotic rate for every nonsynchronous trajectory
is asserted.
\end{proof}

\begin{proposition}[Codimension and the meeting of the algebraic loci]
\label{prop:P09-codimension}
For fixed admissible $(n,q)$, the degenerate origin locus is a smooth
codimension-two submanifold of $(0,\infty)^5$, given by
\begin{equation}
 a-k-h=0,\qquad g-a-c=0.\label{eq:P09-degenerate-locus}
\end{equation}
The branch-intersection and P02 repeated-root loci meet there.
This is an algebraic/spectral codimension statement only.
\end{proposition}
\begin{proof}
In coordinates $(a,k,c,g,h)$ the gradients are
\[
 (1,-1,0,0,-1),\qquad(-1,0,-1,1,0).
\]
They are independent: the $g$ coordinate forces the coefficient of
the second vector to be zero, and the $k$ coordinate then forces the
first to be zero. The regular level-set theorem proves codimension two.
The locus is nonempty: any $k,c,h>0$ give $a=k+h$, $g=a+c>0$.
The center eigenvalue remains simple, not a double-zero spectral
 degeneracy.

P02's Proposition~\ref{prop:P02-loci} gives
\[
 g_f=a+c+\frac{(a-A)^2}{4A},\qquad r_f=\frac{a-A}{2A}.
\]
Hence $g_f=g_c$ exactly at $a=A$, where $r_f=0$.
The repeated root is positive only for $a>A$; the meeting point is
the origin, not a positive nonzero equilibrium. Only this algebraic
overlap is recorded. It is not a saddle-node proof; that remains P10.
\end{proof}


% EDITION: citation-based shortening, 4 October 2026.
% Earlier audits refer to the source edition; this revision has not been independently audited.
\section{P10: Saddle-node bifurcation and true bistable region}
\label{sec:P10}

Fix $a,k,c,h>0$, $n\ge5$, $1\le q<n$, and vary $g>0$ in the
canonical DM field \eqref{eq:canonical-dm}. Put
\begin{equation}
 A=k+h,\quad b=a-A>0,\quad g_c=a+c,\quad
 g_f=g_c+\frac{b^2}{4A},\quad r_f=\frac b{2A},\quad s_f=1+r_f.
 \label{eq:P10-data}
\end{equation}
Here $r$ is the synchronous amplitude and $b$ is a local abbreviation.
\emph{Notation reconciliation:} P02's roots $s_-,s_+$ are the amplitudes
called $r_-,r_+$ here. P02's repeated-root amplitude $s_f$ in
\eqref{eq:P02-gf} is our $r_f$; P10's $s_f$ denotes its denominator
$1+r_f$, following P03/P04's denominator convention. No earlier notation
or file is changed. All neighborhoods below may depend on the fixed
parameters and on $(n,q)$; no uniform spectral gap as $n\to\infty$ is claimed.

The direct inputs are P02's invariant restriction and root formulas,
P03's original-field Jacobian and origin stability, P04's complete
Fourier spectrum, and P09's origin branch organization. P01 supplies
smoothness, uniqueness and global nonnegative existence. We also use
the real implicit-function theorem, the usual smooth local stable/unstable
manifold theorem, and linearized stability/instability. These standard
local results introduce no new project-module dependency or literature work.

\subsection{Exact fold geometry}
\label{slot:P10-1}

P02 gives, for $r>-1$,
\begin{equation}
 F(r\one,g)=f(r,g)\one,\qquad
 f(r,g)=ar-Ar^2+\frac{(c-g)r}{1+r}
       =\frac r{1+r}(\mu+br-Ar^2),\quad \mu=g_c-g.
 \label{eq:P10-scalar}
\end{equation}
The nonzero equilibrium equation is $Ar^2-br-\mu=0$.
Completing the square, or using its discriminant, verifies independently
\begin{equation}
 g_f=c+\frac{(a+A)^2}{4A},\qquad
 s_f=1+\frac{a-A}{2A}=\frac{a+A}{2A},\qquad
 g_f-g_c=\frac{b^2}{4A}>0.
 \label{eq:P10-fold-identities}
\end{equation}
Define the state and parameter coordinates
\begin{equation}
 u=r-r_f,\qquad \nu=g_f-g.
 \label{eq:P10-coordinates}
\end{equation}
Then $\mu=\nu-b^2/(4A)$, $br_f-Ar_f^2=b^2/(4A)$, and
$b-2Ar_f=0$. Expanding in $u$ therefore proves, with no approximation,
\begin{equation}
 \boxed{\mu+br-Ar^2=\nu-Au^2,\qquad
 \dot u=\frac{r_f+u}{s_f+u}(\nu-Au^2).}
 \label{eq:P10-exact-fold}
\end{equation}
Choose $|u|<r_f/2$. Both amplitude and denominator are positive.
For $\nu<0$ there is no nonzero synchronous equilibrium anywhere on
$r>-1$; in particular none near $r_f$. For $\nu=0$ there is one
double positive scalar equilibrium $r_f$. For small $\nu>0$ there
are two distinct nearby positive roots
\begin{equation}
 \boxed{r_\pm=r_f\pm\sqrt{\nu/A}=s_\pm\ \text{of P02}.}
 \label{eq:P10-local-roots}
\end{equation}
Both are positive exactly when $0<\nu<b^2/(4A)$ within this real
root regime. The origin always exists, but is excluded from the
chosen fold state neighborhood.

\subsection{Complete fold spectrum and the real transverse reduction}

For a synchronous state let $s=1+r$ and import P03/P04:
\[
 d=a-2kr-\frac{cr}{s^2}-\frac gs-hr,\qquad
 C=\frac cs,\quad G=\frac{gr}{s^2},\quad H=hr.
\]
\begin{align}
 J(r,g)&=dI+CP+GP^{-q}-HP^q,\notag\\
 \lambda_m&=d+Ce^{-\ii\theta_m}+Ge^{\ii q\theta_m}-He^{-\ii q\theta_m},
 \label{eq:P10-spectrum}
\end{align}
where $\theta_m=2\pi m/n$ and $0\le m<n$.
Colliding offsets are added, as in P03; none is removed.
At a positive equilibrium P04 gives
\begin{equation}
 \lambda_0=f_r(r,g)=\frac r{1+r}(b-2Ar).
 \label{eq:P10-uniform}
\end{equation}
Thus $\lambda_0^f=0$, also obtained by differentiating
\eqref{eq:P10-exact-fold} at $(u,\nu)=(0,0)$.

\begin{lemma}[Automatic transverse stability at every positive fold]
\label{lem:P10-gap}
At the fold, with subscript $f$ denoting evaluation there,
\begin{equation}
 \boxed{G_f-H_f=r_f\left(k+\frac{4A^2c}{(a+A)^2}\right)>0.}
 \label{eq:P10-Lf}
\end{equation}
The exact complex fold spectrum is \eqref{eq:P10-spectrum}, with
$d_f=-C_f-G_f+H_f$. In particular,
\begin{align}
 \Ree\lambda_m^f
 &=C_f(\cos\theta_m-1)+(G_f-H_f)(\cos(q\theta_m)-1),
 \label{eq:P10-fold-real}\\
 \Imm\lambda_m^f&=-C_f\sin\theta_m+(G_f+H_f)\sin(q\theta_m).
 \label{eq:P10-fold-imag}
\end{align}
For every $m\ne0$,
\begin{equation}
 \boxed{\Ree\lambda_m^f\le-\delta_f<0,\qquad
 \delta_f=\frac c{s_f}\left(1-\cos\frac{2\pi}{n}\right)>0.}
 \label{eq:P10-gap}
\end{equation}
The only eigenvalue on the imaginary axis is zero. It has algebraic
and geometric multiplicity one and real eigenspace
$\operatorname{span}\{\one\}$.
\end{lemma}
\begin{proof}
Using \eqref{eq:P10-fold-identities}, directly
\begin{equation}
 \frac{g_f}{s_f^2}-h
 =\frac{c+A s_f^2}{s_f^2}-h
 =k+\frac c{s_f^2}=k+\frac{4A^2c}{(a+A)^2}>0.
 \label{eq:P10-effective}
\end{equation}
Multiplication by $r_f>0$ proves \eqref{eq:P10-Lf}.
Since $0=\lambda_0^f=d_f+C_f+G_f-H_f$, taking real and imaginary
parts of the complete P04 formula proves the two identities.
For $1\le m<n$, set $j=\min(m,n-m)$. The angle
$2\pi j/n$ belongs to $[2\pi/n,\pi]$, so
$\cos\theta_m\le\cos(2\pi/n)<1$. The first cosine difference
is strictly negative, while the second is nonpositive. This proves
the margin. For even $n$ the real singleton $m=n/2$ is included;
for odd $n$ the same argument ends before $\pi$. Collisions
$q=1,n-1,n/2$ and $\gcd(n,q)>1$ change neither argument.
Even if $q\theta_m$ is a multiple of $2\pi$, the first term is
strictly negative. P04's unitary diagonalization lists exactly $n$
eigenvalues with algebraic multiplicity. Only label zero can vanish,
so its multiplicities are one and its eigenvector is $\one$.
No simplicity of the other numerical eigenvalues is asserted.
\end{proof}

\begin{proposition}[Full-system local equilibrium completeness]
\label{prop:P10-reduction}
In a state neighborhood of $r_f\one$ and parameter interval about
$g_f$, every full-system equilibrium is synchronous. The entire local
equilibrium equation reduces exactly to $\nu-Au^2=0$.
\end{proposition}
\begin{proof}
Use $x=r\one+y$ with $E=\{y:\one^{\mathsf T}y=0\}$ and
$\Pi_\perp=I-\one\one^{\mathsf T}/n$. By Lemma~\ref{lem:P10-gap},
the transverse derivative of $\Pi_\perp F(r\one+y,g)$ at the fold
is the invertible map $J_f|_E$. Apply the implicit-function and
synchrony-uniqueness argument of Theorem~\ref{thm:P09-reduction}
\cite[Ch.~1]{GuckenheimerHolmes1983}. It gives $y=0$ for every
nearby equilibrium. Since $r>0$, the remaining equation is
$\nu-Au^2=0$ by \eqref{eq:P10-exact-fold}.
\end{proof}

\subsection{Full-system nondegeneracy and exact center dynamics}

\begin{proposition}[Original-vector-field saddle-node coefficients]
\label{prop:P10-coefficients}
With the above $v,w$, one has $J_fv=0$, $w^{\mathsf T}J_f=0$,
$w^{\mathsf T}v=1$, and
\begin{align}
 \alpha_{\rm SN}:=w^{\mathsf T}F_g(r_f\one,g_f)
 &=-\frac{r_f}{s_f}<0,\label{eq:P10-alpha}\\
 w^{\mathsf T}F_{xx}(r_f\one,g_f)[v,v]
 &=-\frac{2Ar_f}{s_f}\ne0,\label{eq:P10-contraction}\\
 \beta_{\rm SN}:=\tfrac12w^{\mathsf T}F_{xx}[v,v]
 &=-\frac{Ar_f}{s_f}<0.\label{eq:P10-beta}
\end{align}
Derivatives hold all other physical coefficients fixed. With $\nu$
as parameter the parameter coefficient has the opposite sign,
$+r_f/s_f$.
\end{proposition}
\begin{proof}
The row and column sums of $J_f$ are both $\lambda_0^f=0$, giving
the right and left nullvectors and their normalization. In the original
vector field, $(F_g)_i=-x_i/(1+x_{i+q})$, hence its average at the
fold is $-r_f/s_f$.

Independently differentiate each original summand twice along
$x=(r_f+t)\one$ at fixed $g=g_f$. The five componentwise second
derivatives, retaining all terms even at coincident offsets, are
\[
 0,\qquad -2k,\qquad -\frac{2c}{s_f^3},\qquad
 \frac{2g_f}{s_f^3},\qquad -2h.
\]
Their sum is $-2A+2(g_f-c)/s_f^3$. Since $g_f-c=A s_f^2$,
this equals $-2A+2A/s_f=-2Ar_f/s_f$ in every component.
Averaging produces no extra factor of $n$. The Taylor coefficient
of $u^2$ is half this full contraction; the derivatives themselves
contain no Taylor factorial.
\end{proof}

\begin{theorem}[Nondegenerate full-system positive saddle-node]
\label{thm:P10-saddle-node}
For every fixed canonical $a>k+h$, $c,k,h>0$ and admissible $(n,q)$,
the equilibrium $r_f\one$ at
\[
 \boxed{g=g_f=a+c+\frac{(a-k-h)^2}{4(k+h)}}
\]
is a nondegenerate full-system saddle-node. It has a simple uniform
zero eigenvalue and strictly stable transverse spectrum. Locally there
are no equilibria for $g>g_f$, one double equilibrium for $g=g_f$,
and precisely the two positive equilibria \eqref{eq:P10-local-roots}
for $g<g_f$ sufficiently close. The exact center restriction is
orbitally equivalent, by a positive change of time, to
\begin{equation}
 \boxed{\frac{du}{d\tau}=\nu-Au^2,\qquad \frac{d\nu}{d\tau}=0.}
 \label{eq:P10-normal-form}
\end{equation}
\end{theorem}
\begin{proof}
The field is smooth near the positive fold. Lemma~\ref{lem:P10-gap}
and Proposition~\ref{prop:P10-coefficients} verify the simple-zero
and nonzero parameter/quadratic coefficients of the saddle-node theorem
\cite{KuznetsovFold2006,GuckenheimerHolmes1983};
Proposition~\ref{prop:P10-reduction} gives local equilibrium completeness.
For the parameter extension, the linearization in $(u,y,\nu)$ is
\begin{equation}
 \dot u=(r_f/s_f)\nu,\qquad \dot y=(J_f|_E)y,\qquad\dot\nu=0.
 \label{eq:P10-extended-linear}
\end{equation}
Thus the invariant surface $y=0$ is tangent to the whole generalized
center subspace and is an exact center manifold. On it,
\eqref{eq:P10-exact-fold} and the positive factor
$(r_f+u)/(s_f+u)$ for $|u|<r_f/2$ give
\eqref{eq:P10-normal-form} by $d\tau/dt=(r_f+u)/(s_f+u)$.
The roots and their sides follow immediately. At $\nu=0$, negative
$u$ escapes, so the fold is not locally asymptotically stable.
\end{proof}

\subsection{Both branches and all Fourier modes}
\label{slot:P10-2}

\begin{proposition}[Exact uniform signs throughout the two-root interval]
\label{prop:P10-branch-signs}
For every $g_c<g<g_f$, define $D$ as in P02 and $\Delta=\sqrt D$.
Then
\begin{align}
 0<\Delta&=\sqrt{b^2+4A(a+c-g)}=2\sqrt{A(g_f-g)}<b,\notag\\
 r_\pm&=\frac{b\pm\Delta}{2A},\qquad
 b-2Ar_\pm=\mp\Delta,\label{eq:P10-discriminant}\\
 \boxed{\lambda_0(r_\pm)=\mp\frac{r_\pm}{1+r_\pm}\Delta.}
 \label{eq:P10-branch-eigenvalues}
\end{align}
Thus the lower branch is always unstable in the uniform direction,
whereas the upper uniform eigenvalue is always strictly negative.
Any upper-branch loss of stability within this open interval must
involve a nonuniform mode.
\end{proposition}
\begin{proof}
P02's discriminant factorization is also obtained directly by inserting
$g_f=g_c+b^2/(4A)$. The strict parameter inequalities give $0<\Delta<b$,
so both roots are positive. Substitution into $b-2Ar_\pm$ cancels
$b$ and leaves $\mp\Delta$. Equation~\eqref{eq:P10-uniform}
gives the exact signs, and the positive lower eigenvalue implies
nonlinear instability. This last implication needs no hyperbolicity
of the other lower-branch eigenvalues.
\end{proof}

\begin{corollary}[Near-fold full branch stability]
\label{cor:P10-near-stability}
There exists $\varepsilon_{\rm sp}>0$ such that both branches for
$0\le\nu<\varepsilon_{\rm sp}$ satisfy
\[
 \Ree\lambda_m(r_\pm,g_f-\nu)<-\delta_f/2\quad(m\ne0).
\]
For $0<\nu<\min(\varepsilon_{\rm sp},b^2/(4A))$, the upper
branch is Hurwitz and locally exponentially stable. The lower branch
is a hyperbolic saddle with exactly one simple positive real eigenvalue
and $n-1$ stable real dimensions. Its local stable manifold has dimension
$n-1$ and codimension one, and its unstable manifold has dimension one.
\end{corollary}
\begin{proof}
Every labeled eigenvalue in \eqref{eq:P10-spectrum} is continuous
in $(r,g)$ near $(r_f,g_f)$. There are finitely many labels.
A common neighborhood makes each transverse real part differ from
its fold value by less than $\delta_f/2$. Both continuous roots
$r_\pm(\nu)$ tend to $r_f$ as $\nu\downarrow0$; smoothness in
$\nu$ at zero is not needed. Shrink $\varepsilon_{\rm sp}$ so both
branches lie in that neighborhood and remain positive. Combine the
margin with \eqref{eq:P10-branch-eigenvalues}. Only label zero has
positive real part on the lower branch, so its eigenvalue is simple.
Real conjugate-pair blocks and real singleton modes together give
$n-1$ stable dimensions, even if stable numerical eigenvalues repeat.
P04's nonlinear stability implications and the smooth local
stable/unstable manifold theorem apply.
\end{proof}

\subsection{The nonempty open bistability set}
\label{slot:P10-3}

\begin{theorem}[Stable coexistence and openness in physical parameters]
\label{thm:P10-bistability}
For each fixed $a,k,c,h>0$ with $a>k+h$ and each admissible $(n,q)$,
there is $\varepsilon>0$ with
\begin{equation}
 (g_f-\varepsilon,g_f)\subset(g_c,g_f)
 \label{eq:P10-bistable-interval}
\end{equation}
such that throughout this interval the full system has the locally
exponentially stable origin, a positive index-one saddle $r_-\one$,
and the locally exponentially stable positive equilibrium $r_+\one$.
At every strict point of a sufficiently small such interval this
configuration persists on an open neighborhood in $(0,\infty)^5$
with coordinates $(a,k,c,g,h)$ and fixed $(n,q)$.
Here ``bistability'' means coexistence of at least two locally stable
equilibria; it does not mean that these are the only attractors.
\end{theorem}
\begin{proof}
At the origin P03 gives
\[
 \Ree\lambda_m(0)=g_c-g-c(1-\cos\theta_m)\le g_c-g<0
 \quad(g_c<g<g_f).
\]
Thus the origin is Hurwitz and locally exponentially stable in
$\Domain$, and also relative to $\Orth$ using P01 invariance.
Take $0<\varepsilon<\min(\varepsilon_{\rm sp},b^2/(4A))$ and
apply Corollary~\ref{cor:P10-near-stability}. The interval is nonempty
and lies strictly above $g_c$.

At a strict witness all three Jacobians are invertible and hyperbolic.
The implicit-function theorem continues each equilibrium smoothly
under changes of all five physical coefficients. The origin stays
identically zero. Alternatively the positive roots are explicitly
smooth because $D>0$; uniqueness identifies them with the IFT
continuations. Their separation and strict positivity persist after
shrinking the parameter neighborhood. Spectral continuity preserves
both Hurwitz properties and the lower saddle's unstable dimension.
The inequalities $a>A$, $g_c<g<g_f$ and positivity of every coefficient
also persist. Such neighborhoods can depend on $(n,q)$; no common
all-dimension persistence neighborhood for the index-one configuration
is asserted.
\end{proof}

\begin{proposition}[Exact coexistence set and a simple sufficient criterion]
\label{prop:P10-bistability-set}
For fixed $(n,q)$ define the open two-positive-root domain
\[
 \mathcal D=\{\param\in(0,\infty)^5:a>k+h,\ g_c<g<g_f\}.
\]
For $\param\in\mathcal D$, evaluate \eqref{eq:P10-spectrum} at its
upper root and define
\begin{equation}
 \boxed{\mathcal B_{n,q}=
 \left\{\param\in\mathcal D:
 \max_{0\le m<n}\Ree\lambda_m(r_+,g)<0\right\}.}
 \label{eq:P10-B}
\end{equation}
This set is open and nonempty. It is exactly the set in $\mathcal D$
where the origin and upper synchronous equilibrium are both locally
exponentially stable. Every point has an unstable lower positive
branch, although its unstable dimension is asserted to be one only
in the near-fold neighborhood above. A sufficient condition for
membership in $\mathcal B_{n,q}$ is
\begin{equation}
 \boxed{G\ge H\quad\Longleftrightarrow\quad
 \frac{g}{(1+r_+)^2}\ge h.}
 \label{eq:P10-sufficient}
\end{equation}
This condition allows equality.
\end{proposition}
\begin{proof}
The positive upper root and all its labeled eigenvalues are continuous
(and smooth before taking the finite maximum) on $\mathcal D$.
The strict finite set of inequalities defines an open subset of
$(0,\infty)^5$, hence of $\R^5$. Nonemptiness follows from
Theorem~\ref{thm:P10-bistability}. The origin is automatically Hurwitz
on $\mathcal D$. A Hurwitz Jacobian gives nonlinear local exponential
stability. Conversely, a smooth equilibrium satisfying a local
exponential estimate has a Jacobian semigroup satisfying that estimate
by differentiation of the flow at the equilibrium, so all eigenvalues
have strictly negative real part. Thus the definition is exact for
exponential stable coexistence, not a classification of possible
nonhyperbolic asymptotically stable boundary points or other attractors.

Subtracting the uniform row sum in \eqref{eq:P10-spectrum} gives,
at either branch and in fact at any synchronous state,
\begin{equation}
 \Ree\lambda_m=\lambda_0+C(\cos\theta_m-1)
                         +(G-H)(\cos(q\theta_m)-1).
 \label{eq:P10-difference}
\end{equation}
On the upper branch $\lambda_0<0$ and $C>0$. If $G-H\ge0$,
every mode has negative real part, including $m=0$; for $m\ne0$
the $C$ contribution is strictly negative. Since $r_+>0$, dividing
$G-H$ by it gives the physical-coefficient condition. The lower
uniform sign follows from Proposition~\ref{prop:P10-branch-signs}.
\end{proof}

One can give an explicit, dimension-independent sufficient interval
for the two stable equilibria (not necessarily for the saddle's index).
Set
\begin{equation}
 M=g_f-hs_f^2=c+ks_f^2>0,\qquad
 t_*:=\min\left\{\frac{r_f}{2},\frac{M}{4hs_f},
                         \sqrt{\frac{M}{2(A+h)}}\right\}>0.
 \label{eq:P10-explicit-margin}
\end{equation}
For $0<\nu<A t_*^2$, put $t=\sqrt{\nu/A}<t_*$. Then
$g=g_f-At^2$, $r_+=r_f+t$, and directly
\begin{equation}
 g-h(1+r_+)^2=M-2hs_ft-(A+h)t^2>0.
 \label{eq:P10-explicit-condition}
\end{equation}
Indeed each of the last two terms is strictly less than $M/2$.
Also $At_*^2\le b^2/(16A)<g_f-g_c$, so this whole interval lies
in $\mathcal D$ and in $\mathcal B_{n,q}$ for every admissible
$(n,q)$. This supplies a quantitative alternative to the continuity
argument for stable coexistence. Equation~\eqref{eq:P10-effective}
also proves directly that $G>H$ persists on the upper branch near
the fold. Intersecting this interval with that of
Corollary~\ref{cor:P10-near-stability} retains the index-one conclusion.

\begin{remark}[An exact obstruction to whole-interval stability]
\label{rem:P10-counterexample}
Take
\begin{equation}
 n=6,\quad q=1,\qquad
 (a,k,c,g,h)=\left(6,\frac1{10},\frac1{10},\frac{101}{10},\frac9{10}\right).
 \label{eq:P10-counterexample}
\end{equation}
Here $A=1$, $b=5$, $g_c=61/10$, $g_f=247/20$, $\mu=-4$,
and the equilibrium polynomial is $(r-1)(r-4)$.
At $r_+=4$ the original-field equilibrium residual is exactly zero,
and
\[
 d=-\frac{109}{250},\quad C=\frac1{50},\quad
 G=\frac{202}{125},\quad H=\frac{18}{5},\quad
 G-H=-\frac{248}{125}.
\]
Consequently $\lambda_0=-12/5<0$, but the real mode $m=3$ has
\begin{equation}
 \lambda_3=d-C-G+H=\frac{191}{125}>0.
 \label{eq:P10-counterexample-mode}
\end{equation}
The origin is stable and the upper equilibrium unstable at this
strict interior point of the two-root interval. This refutes universal
upper-branch stability, hence universal origin/upper stable coexistence
on that interval. At the lower root $r_-=1$, one has
$\lambda_0=3/2$ and $\Ree\lambda_1=\Ree\lambda_5=53/80>0$;
the other three real parts are negative. Thus its unstable dimension
is three here, also precluding a whole-interval index-one claim.
This makes no assertion about other attractors there.
It is a static full-spectrum counterexample, not a theorem about a
crossing, a Hopf bifurcation or a periodic branch.
\end{remark}

\subsection{Synchronous threshold and the precise local basin statement}
\label{slot:P10-4}

\begin{proposition}[Exact basin threshold on the nonnegative synchronous ray]
\label{prop:P10-ray}
For every $g_c<g<g_f$, the scalar flow on the ray satisfies
\[
 \begin{array}{c|ccc}
 r&(0,r_-)&(r_-,r_+)&(r_+,\infty)\\ \hline
 f(r,g)&-&+&-
 \end{array}
\]
If $0<r(0)<r_-$, then $r(t)\to0$. If $r(0)>r_-$, then
$r(t)\to r_+$; this includes the constant solution $r(0)=r_+$.
The data $r(0)=r_-$ and $r(0)=0$ give their respective constant
solutions. These are exact full-system solutions by synchrony,
even when the upper equilibrium is transversely unstable.
\end{proposition}
\begin{proof}
Factor \eqref{eq:P10-scalar} as
\begin{equation}
 f(r,g)=-\frac{Ar}{1+r}(r-r_-)(r-r_+).
 \label{eq:P10-ray-factor}
\end{equation}
The prefactor is negative for $r>0$, giving the three signs.
Uniqueness prevents crossing an equilibrium. P01 and P02 give global
nonnegative existence and boundedness. More specifically solutions
below $r_-$ decrease and are bounded below by zero; solutions between
the roots increase and are bounded above by $r_+$; those above $r_+$
decrease and are bounded below by it. Each therefore has a limit.
If that limit were not a zero of $f$, continuity would bound $|f|$
away from zero near the limit, contradicting convergence. The monotone
interval and the initial strict inequality select the claimed zero.
No nonsynchronous basin conclusion follows from this argument alone.
\end{proof}

\begin{proposition}[Unstable synchrony and a local separating stable disk]
\label{prop:P10-local-basins}
At a near-fold parameter tuple covered by
Corollary~\ref{cor:P10-near-stability} and
Theorem~\ref{thm:P10-bistability}, let $p=r_-\one$.
The local stable manifold is a smooth embedded $(n-1)$-disk tangent
to $E$. Its germ locally separates a sufficiently small neighborhood
into two sides. The local unstable manifold is exactly a segment of
the synchronous line through $p$. Its two forward branches tend,
respectively, to the origin and to $r_+\one$.

In fact one may choose a neighborhood $V\subset\PosOrth$ of $p$
such that its local stable disk $S$ has
\begin{equation}
 V\setminus S=V_-\mathbin{\dot\cup}V_+,\qquad
 V_-\subset\mathcal A(0),\quad V_+\subset\mathcal A(r_+\one),
 \label{eq:P10-local-partition}
\end{equation}
where the two sets are connected and $\mathcal A(z)$ denotes the
basin of the equilibrium $z$ relative to $\Orth$. Every point of
$S$ tends to $p$. Thus $S$ is the common basin boundary \emph{relative
to this neighborhood}. No assertion identifies the full global basin
boundary or globally continued stable manifold.
\end{proposition}
\begin{proof}
The stable/unstable manifold theorem \cite[Ch.~1]{GuckenheimerHolmes1983}
applies to the hyperbolic index-one saddle. On a small synchronous
interval, every point converges exponentially to $p$ in negative time;
uniqueness of the local unstable manifold therefore identifies it with
that interval. Proposition~\ref{prop:P10-ray} supplies its two forward limits.

To identify the basin sides, write the stable disk as $s=\phi(y)$ in
$x=p+s\one+y$, $y\in E$, and set $\xi=s-\phi(y)$, $z=y$.
Invariance and smoothness give
\begin{equation}
 \dot\xi=\xi(\lambda+\alpha(\xi,z)),\qquad
 \dot z=Bz+R(\xi,z),\quad \lambda=\lambda_0(p)>0,\quad B=J(p)|_E,
 \label{eq:P10-saddle-chart}
\end{equation}
with $\alpha(0,0)=0$, $R(\xi,0)=0$, and $D_zR(0,0)=0$.
Normality and the stable gap give $z^{\mathsf T}Bz\le-\gamma\|z\|^2$.
Shrink the neighborhood until $|\alpha|<\lambda/2$ and
$\|R\|\le(\gamma/2)\|z\|$. Then $\|z\|$ decreases at least
exponentially and every nonzero $|\xi|$ grows at rate at least $\lambda/2$.

Choose small exit points $(\pm\rho,0)$ on the two unstable branches.
Each belongs to its respective sink's open basin. Hence for sufficiently
small $\sigma>0$ both caps $\{(\pm\rho,z):\|z\|\le\sigma\}$
remain in those basins. In the box $|\xi|<\rho$, $\|z\|<\sigma$,
a point with $\xi\ne0$ exits through the cap of its sign and converges
to that sink; points with $\xi=0$ converge to $p$. The two connected
half-boxes and the stable disk give exactly \eqref{eq:P10-local-partition}.
This argument establishes the boundary only relative to that neighborhood.
\end{proof}

\paragraph{Global scope limit.}
The ray threshold is global only within the invariant synchronous ray.
The preceding full-space statement concerns one sufficiently small
neighborhood of the lower saddle at a fixed near-fold tuple. It neither
identifies the complete global boundary between the sinks nor gives a
partition of the whole orthant. Other equilibria, periodic attractors,
or other invariant attracting sets elsewhere are not excluded. In
particular ``bistability'' here does not assert exactly two global attractors.

\subsection{Connection to P09, degenerate limit, and the P11 boundary}

For $b>0$, as $g\downarrow g_c$ from above,
\begin{equation}
 r_-\longrightarrow0,\qquad r_+\longrightarrow b/A,
 \qquad \frac{dr_-}{dg}=\frac1\Delta\longrightarrow\frac1b.
 \label{eq:P10-transcritical-connection}
\end{equation}
Thus P09's small positive branch on $g>g_c$ is precisely P02/P10's
lower branch, unstable in the uniform direction throughout
$(g_c,g_f)$. It meets the upper branch at $g_f$. No upper stability
all the way to $g_c$ follows; Remark~\ref{rem:P10-counterexample}
shows why this would be false in general.

For fixed $k,c,h>0$, as $b=a-A\downarrow0$,
\begin{equation}
 r_f=\frac b{2A}\to0,\qquad g_f-g_c=\frac{b^2}{4A}\to0,
 \qquad \alpha_{\rm SN},\beta_{\rm SN}\to0.
 \label{eq:P10-degenerate-limit}
\end{equation}
The limiting locus is P09's $a=A$, $g=a+c$, where the center field
is $-Ar^3/(1+r)$, its cubic derivative is $-6A$, and the center
zero remains simple. P09 supplies the codimension-two algebraic/spectral
statement. The P10 nondegeneracy theorem excludes this limit.
There is no two-parameter unfolding or cusp classification here.

The exact facts available for later P11 work are the uniform signs
on the entire open two-root interval, the complete fold spectrum,
and automatic all-mode upper stability sufficiently near the fold.
A subsequent loss of upper stability inside this interval must be
nonuniform. P10 establishes no Hopf existence or frequency, critical
pair simplicity, crossing transversality, Lyapunov coefficient, or
periodic branch. The static counterexample does not begin those tasks.
The frozen direct edge into P10 remains P02/P03/P04/P09, with P01
inherited; recording potential later inputs does not alter the frozen DAG.


% EDITION: citation-based shortening, 4 October 2026.
% Earlier audits refer to the source edition; this revision has not been independently audited.
\section{P11: Hopf spectral conditions and spectral simplicity}
\label{sec:P11}
We prove spectral assertions for the canonical demand-modulated model at
positive synchrony. Direct inputs are P02, P04 and P05; P07 is used only
for the five-node specialization. P03's coefficients are inherited through
P04. P08--P10 and P12 or later results supply no mathematical premise.
In particular, an imaginary pair here is only a \emph{spectral Hopf
candidate}. No parameter crossing, nonlinear nondegeneracy, periodic
orbit or propagation theorem is asserted. The physical path is frozen for
P12 without differentiating along it. The optional realizability result
below concerns spectral boundaries, not first instability.

\subsection{Setting, representative classes and frequency}
\label{slot:P11-1}
Fix $q=2$, $n\ge5$, $a,k,c,g,h>0$, and a specified positive synchronous
equilibrium $x^*\one$. Set
\begin{equation}
 s=1+x^*,\qquad C=\frac cs>0,\qquad
 G=\frac{gx^*}{s^2}>0,\qquad H=hx^*>0.
 \label{eq:P11-effective}
\end{equation}
Use P04's verified spectrum and unitary diagonalization:
\begin{align}
 \lambda_m&=d+Ce^{-\ii\theta_m}+Ge^{2\ii\theta_m}
                       -He^{-2\ii\theta_m},\quad \theta_m=2\pi m/n,
                       \label{eq:P11-input}\\
 r_m&=d+C\cos\theta_m+(G-H)\cos2\theta_m,\nonumber\\
 \omega_m&=-C\sin\theta_m+(G+H)\sin2\theta_m,\nonumber\\
 J&=\Fourier\diag(\lambda_m)\Fourier^*,\qquad
 \lambda_{n-m}=\overline{\lambda_m}\quad(1\le m<n).\nonumber
\end{align}
Here $J=DF^{\DM}(x^*\one)$; $\omega_m$ is a signed imaginary part.
The reserved positive frequency is chosen separately when needed.

Let $N=\lfloor n/2\rfloor$, $u_j=\cos(2\pi j/n)$ and
$L_j$ be P05's representative label classes: $L_0=\{0\}$,
$L_N=\{n/2\}$ if $n$ is even, and otherwise
$L_j=\{j,n-j\}$. Write
\[
 \mathcal S=\argmax_{0\le j\le N}r_j
\]
for the \emph{complete} selected representative set, including every tie.
The selected Fourier labels are $\bigcup_{j\in\mathcal S}L_j$.
This notation distinguishes $\mathcal S$ from the scalar $G+H$.

\begin{lemma}[Reachable nonuniform representatives are genuine pairs]
\label{lem:P11-pairs}
Every reachable nonuniform representative satisfies
$1\le j\le j_{\max}:=\lfloor(n+1)/4\rfloor<n/2$.
Consequently $0<\theta_j<\pi$, $\sin\theta_j>0$, and
$L_j=\{j,n-j\}$ has two different labels. No even-$n$ Nyquist
singleton is included.
\end{lemma}
\begin{proof}
P05 Theorem~\ref{thm:P05-reachable} supplies the reachable list.
Since $(n+1)/4<n/2$ for $n>1$, its nonzero members satisfy the
claimed strict bounds. The angle and sine statements follow.
This counts Fourier labels only, not numerical eigenvalue multiplicities.
\end{proof}

\begin{proposition}[Exact signed frequency and its zero locus]
\label{prop:P11-frequency}
For any genuine nonuniform pair representative $0<j<n/2$,
\begin{equation}
 \boxed{\omega_j=\sin\theta_j\,[2(G+H)u_j-C].}
 \label{eq:P11-frequency}
\end{equation}
Thus $\omega_j=0$ if and only if $2(G+H)u_j=C$.
If $u_j\le0$, then $\omega_j<0$ automatically. If $u_j>0$,
the zero-frequency locus is exactly
\begin{equation}
 G+H=\frac{C}{2u_j}.
 \label{eq:P11-zero-frequency}
\end{equation}
For positive $u_j$ the sign of $\omega_j$ is the sign of
$G+H-C/(2u_j)$. These are spectral statements, not bifurcation classifications.
\end{proposition}
\begin{proof}
Insert $\sin2\theta_j=2\sin\theta_j\cos\theta_j$ in
\eqref{eq:P11-input}. The factor $\sin\theta_j$ is strictly positive.
If $u_j\le0$, the bracket is at most $-C<0$; otherwise divide by
$2u_j>0$. The Nyquist singleton is outside the proposition's scope.
\end{proof}

\subsection{Strict candidates, simplicity and the stable complement}
\label{slot:P11-2}
\begin{definition}[Strict nonuniform spectral Hopf candidate]
\label{def:P11-candidate}
A positive synchronous equilibrium of the canonical $q=2$ model is a
\emph{strict nonuniform spectral Hopf candidate in representative $j$}
if $j\ne0$, $\mathcal S=\{j\}$, $r_j=0$, and $\omega_j\ne0$.
Uniqueness refers to representatives and includes the exclusion of all
maximal ties. This definition does not assert a Hopf bifurcation.
\end{definition}

\begin{theorem}[General strict spectral criterion]
\label{thm:P11-strict}
At a strict candidate, the entire imaginary-axis spectrum is exactly
\begin{equation}
 \lambda_j=\ii\omega_j,\qquad
 \lambda_{n-j}=-\ii\omega_j,\qquad \omega_j\ne0.
 \label{eq:P11-critical-pair}
\end{equation}
Each value is algebraically and geometrically simple over $\C$ and the
two values are distinct. Every other Fourier label has $r_\ell<0$.
The number
\begin{equation}
 \eta=-\max_{\ell\notin L_j}r_\ell
 \label{eq:P11-gap}
\end{equation}
is strictly positive, and all noncritical eigenvalues satisfy
$\Ree\lambda_\ell\le-\eta$. The real critical space has dimension
two and its real invariant complement has dimension $n-2$.
Moreover
\begin{equation}
 0,\ 2\ii\omega_j\notin\spec(J),\qquad
 J^{-1}\ \hbox{and}\ (2\ii\omega_j I-J)^{-1}\ \hbox{exist}.
 \label{eq:P11-resolvent-gates}
\end{equation}
\end{theorem}
\begin{proof}
Strict representative selection leaves precisely labels $j,n-j$ at
real part zero; their values are distinct because $\omega_j\ne0$.
No other label can share either value. P04's multiplicity formula
therefore proves simplicity. The finite complement has maximum
$-\eta<0$ and real dimension $n-2$. Neither $0$ nor $2\ii\omega_j$
equals $\pm\ii\omega_j$ or a stable eigenvalue, proving the inverse gates.
These are the spectral hypotheses used in the usual Hopf theorem
\cite{KuznetsovHopf2006}; crossing and nonlinear nondegeneracy remain separate.
\end{proof}

\begin{proposition}[Exact real and adjoint normalization]
\label{prop:P11-vectors}
At the strict candidate, choose
\begin{equation}
 v_H=v^{(j)},\qquad p_H=v_H,\qquad
 e_c=\sqrt2\,\Ree v_H,\quad e_s=\sqrt2\,\Imm v_H.
 \label{eq:P11-vectors}
\end{equation}
Under the project's Hermitian pairing $\inner{p}{v}=p^*v$,
\begin{equation}
 Jv_H=\ii\omega_jv_H,\qquad
 J^*p_H=-\ii\omega_jp_H,\qquad \inner{p_H}{v_H}=1.
 \label{eq:P11-adjoint}
\end{equation}
The real vectors $e_c,e_s$ are orthonormal and span the real critical
plane. In the ordered basis $(e_c,e_s)$ the matrix is
\begin{equation}
 \left.J\right|_{\rm crit}=
 \begin{pmatrix}0&\omega_j\\-\omega_j&0\end{pmatrix}.
 \label{eq:P11-real-block}
\end{equation}
In its coefficient plane the signed angular velocity is $-\omega_j$.
\end{proposition}
\begin{proof}
Unitary diagonalization gives the adjoint equation with conjugate
eigenvalue and $p_H=v_H$. Orthogonality of $v_H$ and $\bar v_H$
gives $\|\Ree v_H\|^2=\|\Imm v_H\|^2=1/2$ and their orthogonality.
Taking real and imaginary parts of $Jv_H=\ii\omega_jv_H$ gives
$Je_c=-\omega_je_s$ and $Je_s=\omega_je_c$, hence the stated block
and orientation. If a positive frequency is required when $\omega_j<0$,
use the conjugate label.
\end{proof}

\begin{remark}[Ties and zero frequency are excluded separately]
\label{rem:P11-degeneracies}
If two different representative classes have $r_j=r_t=0$, their
labels give additional imaginary-axis values or larger multiplicities.
Either possibility fails the single-simple-pair-with-stable-complement
hypothesis. A tie of real parts alone is not equality of numerical
eigenvalues. In particular choosing one member of a maximal tie does
not satisfy Definition~\ref{def:P11-candidate}.

If a genuine pair has $r_j=\omega_j=0$, both labels produce the
same real zero eigenvalue. P04 makes it semisimple, with complex
algebraic and geometric multiplicity exactly
$|\{m:\lambda_m=0\}|\ge2$. Under strict selection of that
representative its multiplicity is exactly two and its real kernel
has dimension two. Without strict selection other zero labels must
also be counted. This is ineligible for the simple Hopf package.
No named multiple-criticality or steady-state bifurcation is classified.
The two-dimensional real plane in Theorem~\ref{thm:P11-strict}, by
contrast, corresponds to two distinct simple complex eigenvalues.
\end{remark}

\subsection{Regular spectral surfaces and physical admissibility}
Fix a reachable nonuniform representative $1\le j\le j_{\max}$
throughout this subsection, so $0<\theta_j<\pi$.
Use P02's nonzero equilibrium equation in branch coordinates:
\begin{equation}
 a=(k+h)x-\frac{c-g}{1+x}>0,\qquad x,k,c,g,h>0.
 \label{eq:P11-chart}
\end{equation}
This is exactly P05's open chart domain $\mathcal U_{\rm eq}$.
Combining P04's diagonal coefficient with the equilibrium equation gives
\[
 d=a-2kx-\frac{cx}{(1+x)^2}-\frac g{1+x}-hx
   =-kx-\frac{c(1+2x)}{(1+x)^2}.
\]
For fixed representative $j$, put $v_j=\cos2\theta_j=2u_j^2-1$;
this scalar $v_j$ is not a Fourier vector $v^{(j)}$. Define
\begin{equation}
 \begin{split}
 R_j(x,k,c,g,h)={}&-kx-\frac{c(1+2x)}{(1+x)^2}
                 +\frac{c}{1+x}u_j\\
                &+x\left(\frac{g}{(1+x)^2}-h\right)v_j.
 \end{split}
 \label{eq:P11-R}
\end{equation}
Substitution in P04's real dispersion shows exactly $R_j=r_j$.

\begin{proposition}[Regular spectral boundary and excluded frequency locus]
\label{prop:P11-surface}
At fixed $x,c,g,h$,
\begin{equation}
 \frac{\partial R_j}{\partial k}=-x<0,
 \qquad
 k_H=\frac1x\left[-\frac{c(1+2x)}{(1+x)^2}
       +\frac{c}{1+x}u_j
       +x\left(\frac{g}{(1+x)^2}-h\right)v_j\right].
 \label{eq:P11-kH}
\end{equation}
Thus $R_j=0$ is a regular codimension-one level set wherever it
meets $\mathcal U_{\rm eq}$, explicitly the graph $k=k_H$.
A physical point requires $k_H>0$ and
$a_H=(k_H+h)x-(c-g)/(1+x)>0$. Within a strict nonuniform
P05 cell, the portions with $\omega_j\ne0$ are strict candidates,
and this regular surface transfers locally to physical parameter space
along the specified positive branch.

For $u_j>0$, any nonempty intersection $R_j=\omega_j=0$ is
locally a regular codimension-two level set in branch coordinates.
It is codimension one inside $R_j=0$. The same statement transfers
to physical parameters wherever P05's concave-chart hypotheses hold.
For $u_j\le0$ this intersection is empty.
\end{proposition}
\begin{proof}
In \eqref{eq:P11-R} only $-kx$ depends on $k$, proving the
partial derivative and every sign in the solved formula. Strict
selection and $\omega_j\ne0$ are open conditions.
P05 Theorem~\ref{thm:P05-selection} implies $H-G>0$ at any
nonuniform selected representative. On precisely this concave part
of the positive chart, P05 supplies
\[
 \Phi(x,k,c,g,h)=(a,k,c,g,h),\qquad
 \det D\Phi=k+h+\frac{c-g}{(1+x)^2}
     =k+\left(h-\frac{g}{(1+x)^2}\right)
         +\frac{c}{(1+x)^2}>0.
\]
Its local diffeomorphism therefore transfers the regular level set.
No global injectivity or topology after forgetting the branch is claimed.
The derivative $\partial_k R_j$ holds $x,c,g,h$ fixed and lets $a$
follow \eqref{eq:P11-chart}; it is not a derivative along the later
fixed-$a$ continuation path.

For the additional frequency equation, at fixed branch coordinates
\[
 \partial_k\omega_j=0,\qquad
 \partial_g\omega_j=\frac{2xu_j\sin\theta_j}{(1+x)^2}>0
 \quad(u_j>0).
\]
The Jacobian of $(R_j,\omega_j)$ in columns $(k,g)$ has determinant
$-2x^2u_j\sin\theta_j/(1+x)^2\ne0$. The regular level set theorem
proves the local codimension assertion, conditional on existence and
physical positivity; the local diffeomorphism applies under the same
concavity hypotheses just stated. Proposition~\ref{prop:P11-frequency}
excludes all $u_j\le0$ intersections. No unfolding is inferred.
\end{proof}

\begin{proposition}[Exact obstruction and cell realizability criterion]
\label{prop:P11-realizability}
Fix a reachable nonuniform representative $1\le j\le j_{\max}$.
Let $\beta_j=(u_j+u_{j+1})/2$ be P05's midpoint, and put
\begin{equation}
 w_j=1-2u_j^2,\quad L_j^{\rm vert}=\max(0,\beta_j),\qquad
 B_j=\frac{w_j}{4(1-u_j)}.
 \label{eq:P11-realizability-data}
\end{equation}
Its strict P05 cell contains a physical strict spectral Hopf candidate
if and only if
\begin{equation}
 \boxed{w_j>0\quad\hbox{and}\quad L_j^{\rm vert}<B_j.}
 \label{eq:P11-realizability-test}
\end{equation}
More precisely, for a prescribed positive vertex
$z=C/[4(H-G)]$ in that strict cell and prescribed $x>0$, such a
candidate can be realized by some positive physical coefficients if
and only if
\begin{equation}
 \frac{w_j}{4z}>1-u_j+\frac{x}{1+x}.
 \label{eq:P11-fixed-vertex}
\end{equation}
Consequently not every reachable nonuniform cell supports a candidate.
For example $n=8,j=1$ has $w_j=0$ and never reaches $r_j=0$.
For $n=12,j=2$, $u_j=1/2$, $L_j^{\rm vert}=B_j=1/4$;
this cell also has no physical candidate despite $w_j>0$.
\end{proposition}
\begin{proof}
Write $D_c=H-G>0$ and $C=c/(1+x)$ within the selected cell.
At a candidate, \eqref{eq:P11-R} reads
\begin{equation}
 kx=D_cw_j-C\left(1-u_j+\frac{x}{1+x}\right),\qquad
 \frac{kx}{C}=\frac{w_j}{4z}-\left(1-u_j+\frac{x}{1+x}\right).
 \label{eq:P11-k-obstruction}
\end{equation}
Because $k,x,C>0$ and $u_j<1$, this forces $w_j>0$ and
\eqref{eq:P11-fixed-vertex}. It further forces $z<B_j$.
P05's strict cell is
$L_j^{\rm vert}<z<\beta_{j-1}$, where $\beta_{j-1}>0$;
$\beta_j$ exists since $j_{\max}<N$. Thus necessity of
\eqref{eq:P11-realizability-test} follows.

Conversely, if \eqref{eq:P11-realizability-test} holds, choose
\[
 L_j^{\rm vert}<z<\min(\beta_{j-1},B_j).
\]
Both upper bounds exceed the lower bound. Choose $x>0$ sufficiently
small that \eqref{eq:P11-fixed-vertex} holds. For fixed $x,z$
satisfying that inequality, the following construction proves the more
precise sufficiency as well. Choose any $C>0$, set $D_c=C/(4z)$,
and define $k>0$ by \eqref{eq:P11-k-obstruction}. For a free $G>0$,
set
\begin{equation}
 c=C(1+x),\qquad g=\frac{G(1+x)^2}{x},\qquad
 h=\frac{G+D_c}{x},\qquad
 a=kx+D_c-C+G\left(1+\frac{1+x}{x}\right).
 \label{eq:P11-realization}
\end{equation}
These equations give $H-G=D_c$ and exactly
$a=(k+h)x-(c-g)/(1+x)$. Increasing $G$ makes $a>0$
without changing $z$, $k$ or $r_j=0$. The signed frequency bracket
is $2(D_c+2G)u_j-C$. For $u_j\le0$ it is automatically negative;
for $u_j>0$ it vanishes for at most one $G$, which can be avoided
while retaining positivity. Hence a strict candidate is obtained.
The two examples follow by exact substitution of their cosines and
midpoints. They also show why P05's realization of dominance alone
cannot imply realization of a spectral boundary.
\end{proof}

\subsection{An exact five-node physical candidate}
Use the P07 constants and spectra, without a new P07 derivation:
\[
 \phi=\frac{1+\sqrt5}{2},\quad
 \sigma=\frac{\sqrt{10-2\sqrt5}}4>0,\quad
 u_1=\frac1{2\phi},\quad u_2=-\frac\phi2.
\]
P07 shows that $\{1,4\}$ is the only potentially dominant nonuniform
pair and that $\omega_1=\sigma[(G+H)-\phi C]$.

\begin{theorem}[Exact five-node spectral Hopf candidate]
\label{thm:P11-five}
For $n=5,q=2$ set
\begin{equation}
 \begin{gathered}
 x_H=1,\quad c=2,\quad g_H=16,\quad h=6,\\
 k_H=\frac{3\phi-4}{2}=\frac{3\sqrt5-5}{4},\qquad
 a_H=\frac{3\phi+22}{2}=\frac{47+3\sqrt5}{4}.
 \end{gathered}
 \label{eq:P11-tuple}
\end{equation}
This is a positive physical tuple with a positive synchronous equilibrium
$x_H\one$. It is a strict spectral candidate in representative $1$.
Its full labeled spectrum is
\begin{align}
 \lambda_0&=-\frac{3\phi+1}{2},\label{eq:P11-five-spectrum}\\
 \lambda_1&=\ii\omega_H,\qquad \lambda_4=-\ii\omega_H,
       &\omega_H&=\sigma(10-\phi)>0,\nonumber\\
 \lambda_2&=-\frac{3\sqrt5}{2}-\ii\sigma(1+10\phi),
       &\lambda_3&=\overline{\lambda_2}.\nonumber
\end{align}
The critical values are each algebraically and geometrically simple.
The ordering is $r_1=0>r_0>r_2$, with exact noncritical gap
\begin{equation}
 \eta_H=\frac{3\phi+1}{2}>0.
 \label{eq:P11-five-gap}
\end{equation}
The equilibrium is the upper of two distinct positive synchronous roots.
The other root is
\begin{equation}
 x_- =\frac{3(2-\phi)}{3\phi+8}\in(0,1).
 \label{eq:P11-lower-root}
\end{equation}
\end{theorem}
\begin{proof}
Since $5/3<\sqrt5<3$, $k_H>0$ and $a_H>0$;
the remaining three coefficients are positive. In the original
canonical field, each component at $\one$ is exactly
\[
 F_i^{\DM}(\one)=a_H-k_H+\frac22-\frac{16}2-6
     =a_H-k_H-13=0.
\]
This verifies the original model residual before using any spectral data.
Direct substitution gives
\begin{equation}
 \begin{gathered}
 s=2,\quad C=1,\quad G=4,\quad H=6,\quad T=G-H=-2,\quad S=G+H=10,\\
 d=-k_H-\frac32=\frac{1-3\phi}{2}.
 \end{gathered}
 \label{eq:P11-five-effective}
\end{equation}
The raw P04 expression independently gives
$d=a_H-2k_H-1/2-8-6=-k_H-3/2$.
P07's exact real-part formulas now yield
\begin{align*}
 r_1&=\frac{1-3\phi}{2}+\frac{1}{2\phi}+\phi=0,\\
 r_0&=\frac{1-3\phi}{2}+1-2=-\frac{3\phi+1}{2},\\
 r_2&=\frac{1-3\phi}{2}-\frac\phi2-\frac1\phi
       =\frac{3-6\phi}{2}=-\frac{3\sqrt5}{2}.
\end{align*}
Here $\phi^{-1}=\phi-1$ and $2\phi-1=\sqrt5$.
Their noncritical difference is
$r_0-r_2=(3\phi-4)/2=k_H>0$. Equivalently,
$-\phi^2<-2<-\phi^{-2}$, so P07's full ordering table gives
the same strict ordering. In particular pair $\{2,3\}$ is excluded.

For the imaginary parts, P07 gives
$\omega_1=\sigma(10-\phi)$ and
$\omega_2=-\sigma(1+10\phi)$. The exact bounds $1<\phi<2$
imply $10-\phi>8>0$, so the critical frequency is nonzero without
decimal evidence. P04 conjugacy fixes the signs at labels $4,3$.
This proves every entry of \eqref{eq:P11-five-spectrum}.
The general theorem proves simplicity and the inverse gates, and the
strict noncritical ordering gives \eqref{eq:P11-five-gap}.

For the branch identification, use only P02's abbreviations and polynomial:
\begin{equation}
 A=k_H+h=\frac{3\phi+8}{2},\quad b=a_H-A=7,\quad
 \mu=a_H+2-16=\frac{3\phi-6}{2}<0.
 \label{eq:P11-branch-data}
\end{equation}
Because $1$ is a root and $A-b=\mu$, that polynomial factors as
\[
 Q(X)=AX^2-bX-\mu=A(X-1)(X-x_-),\qquad x_-=-\mu/A.
\]
The displayed $x_-$ is positive because $\phi<2$, and is less than
one since $3(2-\phi)<3\phi+8$. Thus the discriminant is
$A^2(1-x_-)^2>0$. Together with $b>0$ and $\mu<0$ this checks
P02's two-positive-root regime fully; those two signs alone would
not suffice without the positive discriminant. The larger root is
exactly $x_H=1$. No P10 theorem is used.
\end{proof}

\subsection{Frozen continuation path and exact spectral data}
\label{slot:P11-3}
\begin{proposition}[Local branch, path, vectors and resolvents for P12/P13]
\label{prop:P11-data}
Hold $a=a_H,k=k_H,c=2,h=6$ fixed and freeze the physical path
\begin{equation}
 \boxed{\pi(g)=(a_H,k_H,2,g,6),\qquad g_H=16.}
 \label{eq:P11-path}
\end{equation}
There is a neighborhood of $16$ with a unique smooth positive
synchronous root near $1$, denoted $x_+(g)$, with $x_+(16)=1$.
It continues the upper root. With the project's node ordering $0,\ldots,4$,
\begin{equation}
 \begin{gathered}
 v_H=\frac1{\sqrt5}(1,e^{2\pi\ii/5},e^{4\pi\ii/5},
                         e^{6\pi\ii/5},e^{8\pi\ii/5})^{\mathsf T},\\
 p_H=v_H,\qquad \overline{v_H}=v^{(4)},\qquad
 J_Hv_H=\ii\omega_Hv_H,\quad J_H^*p_H=-\ii\omega_Hp_H,
 \quad p_H^*v_H=1.
 \end{gathered}
 \label{eq:P11-five-vectors}
\end{equation}
The real orthonormal basis $e_c=\sqrt2\Ree v_H$,
$e_s=\sqrt2\Imm v_H$ has block
$\left(\begin{smallmatrix}0&\omega_H\\-\omega_H&0\end{smallmatrix}\right)$.
Both required inverse operators exist and are exactly
\begin{equation}
 \begin{split}
 J_H^{-1}&=\sum_{m=0}^4\lambda_m^{-1}v^{(m)}v^{(m)*},\\
 (2\ii\omega_H I-J_H)^{-1}
   &=\sum_{m=0}^4(2\ii\omega_H-\lambda_m)^{-1}v^{(m)}v^{(m)*}.
 \end{split}
 \label{eq:P11-five-resolvents}
\end{equation}
In particular the mode-0 and mode-2 denominators useful for later work are
\begin{equation}
 \lambda_0=-\eta_H\ne0,\qquad
 2\ii\omega_H-\lambda_2
       =\frac{3\sqrt5}{2}+\ii\sigma(21+8\phi)\ne0.
 \label{eq:P11-denominators}
\end{equation}
\end{proposition}
\begin{proof}
Let $B(X,g)=a_H-(k_H+6)X+(2-g)/(1+X)$.
P02's scalar field is $f(X,g)=XB(X,g)$. P04's uniform identity
$\lambda_0=\partial_X f$ at synchrony, together with $B(1,16)=0$,
gives $B_X(1,16)=\lambda_0=-\eta_H<0$.
The scalar implicit-function theorem supplies the unique smooth root
near $1$. Positivity persists in a sufficiently small neighborhood,
and P02's distinct roots and their strict ordering persist there, so
this is the locally continued upper root, not uniqueness of all positive
roots. The parameter tuple stays positive for $g$ near $16$.
No derivative of $x_+(g)$ is calculated.

Proposition~\ref{prop:P11-vectors} gives all vector identities and
the real block with its stated sign. The spectral theorem and
Theorem~\ref{thm:P11-strict} give the two sums, all of whose denominators
are nonzero. Finally,
$2(10-\phi)+(1+10\phi)=21+8\phi$ verifies the second denominator
in \eqref{eq:P11-denominators}; its real part is positive.
No tensor, nonlinear Fourier selection rule or Taylor-factor convention
is required for these spectral identities.
\end{proof}

\begin{proposition}[Spectral hypothesis ledger at the five-node point]
\label{prop:P11-ledger}
The following records the established spectral inputs and the remaining P12/P13 work.
\begin{enumerate}[label=H\arabic*.,leftmargin=3em]
 \item All five physical coefficients in \eqref{eq:P11-tuple} are positive.
 \item $x_H=1$ is an exact positive synchronous equilibrium of the original DM field.
 \item The path varying $g$ in \eqref{eq:P11-path} has the smooth local
       upper branch $x_+(g)$ through $x_+(16)=1$.
 \item Labels $1,4$ give $\pm\ii\omega_H$ with $\omega_H=\sigma(10-\phi)>0$.
 \item Each critical complex eigenvalue is algebraically and geometrically simple.
 \item Every other eigenvalue has real part at most $-\eta_H<0$,
       where $\eta_H=(3\phi+1)/2$.
 \item $0\notin\spec(J_H)$.
 \item $2\ii\omega_H\notin\spec(J_H)$.
 \item Exact normalized right and adjoint vectors are \eqref{eq:P11-five-vectors};
       the pairing and real-plane orientation are fixed.
 \item Spectral transversality is \textbf{not established}; no state or
       eigenvalue derivative along \eqref{eq:P11-path} has been computed.
 \item The first Lyapunov coefficient and nonlinear nondegeneracy are
       \textbf{not established}; no tensor or normal-form coefficient is computed.
 \item No periodic orbit, periodic-branch stability or rotating wave is established.
\end{enumerate}
\end{proposition}
\begin{proof}
Items H1--H2 and H4--H6 are Theorem~\ref{thm:P11-five}; H3 and
H9 are Proposition~\ref{prop:P11-data}; H7--H8 follow from
Theorem~\ref{thm:P11-strict}. H10--H12 explicitly delimit unproved
later conclusions, rather than adding assumptions already certified here.
\end{proof}


% STATUS: conditional framework; outstanding model-specific work is explicit.
\section{P12: Hopf transversality (unfinished application)}
\label{sec:P12}
The source has no derivative calculation. The formulas below specify
the standard reduction; the crossing on P11's path remains to be verified.
\subsection{Equilibrium derivative}\label{slot:P12-1}
For $B(x,g)=a-(k+h)x+(c-g)/(1+x)=0$ with $a,k,c,h$ fixed,
implicit differentiation gives
\[
 B_x=-(k+h)-\frac{c-g}{(1+x)^2},\qquad
 x_g=\frac{1}{(1+x)B_x}\quad(B_x\ne0).
\]
The implicit-function theorem applies locally at P11's simple upper root;
it fails in this parameterization at a scalar fold
\cite[Ch.~1]{GuckenheimerHolmes1983}.
\subsection{Total crossing derivative}\label{slot:P12-2}
Differentiate P04's coefficients along the moving equilibrium:
\[
 \frac{d\lambda_m}{dg}
 =d_g+C_g e^{-\ii\theta_m}+G_g e^{2\ii\theta_m}-H_g e^{-2\ii\theta_m}.
\]
Here every subscript denotes a total derivative, including $x_g$.
\subsection{Outstanding nondegeneracy}\label{slot:P12-3}
Evaluate and certify $\Ree(d\lambda_j/dg)\ne0$ at P11's exact tuple,
with the orientation of increasing $g$ explicit. This is the crossing
hypothesis in \cite{KuznetsovHopf2006}; the spectral conditions alone
do not prove it.


% STATUS: conditional framework; outstanding model-specific work is explicit.
\section{P13: First Lyapunov coefficient (unfinished application)}
\label{sec:P13}
No model-specific derivative tensors or sign calculation were supplied.
\subsection{Derivative tensors}\label{slot:P13-1}
At the candidate, use the Taylor convention
$F(x_H+z)=Jz+\frac12\mathcal B(z,z)+\frac16\mathcal C(z,z,z)+O(\|z\|^4)$,
where $\mathcal B=D^2F$ and $\mathcal C=D^3F$.
Both tensors must be differentiated from the DM field, including its
rational terms and coincident offsets; DM is not a quadratic model.
\subsection{Fourier and resolvent inputs}\label{slot:P13-2}
Use P11's $Jv=\ii\omega v$, $J^*p=-\ii\omega p$, $p^*v=1$,
$\omega>0$, and its verified inverse gates. Cyclic equivariance restricts
products of modes to their sum modulo $n$; retain normalization factors
and possible aliasing when evaluating the actual tensors.
\subsection{Standard coefficient}\label{slot:P13-3}
In this convention the standard formula is \cite{KuznetsovHopf2006}
\begin{align*}
 l_1=\frac1{2\omega}\Ree\bigl\{p^*\bigl[
 &\mathcal C(v,v,\bar v)-2\mathcal B(v,J^{-1}\mathcal B(v,\bar v))\\
 &+\mathcal B(\bar v,(2\ii\omega I-J)^{-1}\mathcal B(v,v))\bigr]\bigr\}.
\end{align*}
The formula replaces a derivation of general normal-form theory; the DM
evaluation and rigorous sign or nonvanishing certificate remain required.
\subsection{Hopf conclusion}\label{slot:P13-4}
Only after P12's nonzero crossing and $l_1\ne0$ have been verified can
the nondegenerate Hopf theorem be applied. Branch stability additionally
uses P11's stable complement. No periodic branch is established here.


% STATUS: conditional framework; outstanding model-specific work is explicit.
\section{P14: Rotating waves (unfinished application)}
\label{sec:P14}
\subsection{Exact symmetry}\label{slot:P14-1}
Once P12--P13 verify a nondegenerate Hopf bifurcation, use equivariant
Hopf theory \cite[Theorem~3.7]{Stewart2024}. The homogeneous DM field
commutes with $P$; the critical Fourier column obeys
$Pv^{(m)}=e^{-2\pi\ii m/n}v^{(m)}$. With $m$ chosen to carry
$+\ii\omega$, the expected branch symmetry is
$Px(t)=x(t-mT/n)$, with the branch's actual period $T$.
This is a conditional application, not an existence proof from a cosine ansatz.
\subsection{Clusters and orientation}\label{slot:P14-2}
The character kernel has size $\gcd(n,m)$, giving $n/\gcd(n,m)$
phase classes. Reconcile the exact time-shift sign with the chosen
eigenvalue and predecessor convention before stating propagation direction.
\subsection{Outstanding verification}\label{slot:P14-3}
Verify the relevant axial subgroup and branch hypotheses and distinguish
exact nonlinear symmetry from its leading Fourier approximation.
Generic ring results do not establish those DM hypotheses or exact
rotating waves after heterogeneous perturbation.


% STATUS: conditional framework; outstanding model-specific work is explicit.
\section{P15: Heterogeneous persistence (unfinished application)}
\label{sec:P15}
The heterogeneous coefficient assignment has not been specified in the source.
\subsection{Local equilibrium}\label{slot:P15-1}
For a specified smooth family $F(x,\eps)$ with $F(x_0,0)=0$ and
invertible $J=D_xF(x_0,0)$, the implicit-function theorem gives a
unique nearby equilibrium $x^*(\eps)$ with the family's regularity
\cite[Ch.~1]{GuckenheimerHolmes1983}. For an analytic family, use the
analytic version stated in \cite[\S3.2]{GreenbaumLiOverton2020}.
\subsection{Positivity and stability}\label{slot:P15-2}
If $x_0>0$, positivity persists locally by continuity. If $J$ is
Hurwitz, local exponential stability persists by spectral continuity.
Invertibility alone does not imply stability; a nonzero-frequency Hopf
equilibrium can have invertible $J$.
\subsection{Remaining scope}\label{slot:P15-3}
Specify where every perturbed coefficient enters, its regularity, and an
admissible parameter neighborhood. Any quantitative radius remains to be
estimated. Optional periodic-orbit persistence requires hyperbolicity
modulo phase, via a transverse Poincare map
\cite[Ch.~1]{GuckenheimerHolmes1983}; equilibrium persistence is insufficient.


% STATUS: conditional framework; outstanding model-specific work is explicit.
\section{P16: Static sensitivity (unfinished application)}
\label{sec:P16}
\subsection{Physical forcing}\label{slot:P16-1}
Under P15's hypotheses, differentiation gives
$\xi=x^{*\prime}(0)=-J^{-1}b$, where $b=F_\eps(x_0,0)$.
The vector $b$ must be computed from the chosen heterogeneous assignment.
\subsection{Fourier inverse}\label{slot:P16-2}
P04 and circulant inversion \cite[Theorem~3.1]{Gray2006} give
\[
 \widehat\xi_m=-\widehat b_m/\lambda_m,\qquad
 \|\xi\|_2\le\frac{\|b\|_2}{\min_m|\lambda_m|},\qquad
 (J^{-1})_{j\ell}=\frac1n\sum_m\frac{e^{\ii\theta_m(j-\ell)}}{\lambda_m}.
\]
\subsection{Localization}\label{slot:P16-3}
Graph-distance decay may be obtained from
\cite[Theorem~3.4]{BenziRazouk2007} only after verifying its spectral
continuum and analytic-neighborhood hypotheses for $f(z)=1/z$,
including exclusion of zero. Pointwise invertibility does not supply
uniform decay constants. Use the graph-distance convention of that theorem;
any directed or wraparound refinement needs its own argument.
\subsection{Remainder}\label{slot:P16-4}
The physical forcing, neighborhood, remainder constants and any uniformity
in $n$ remain uncomputed. A nonzero imaginary Hopf pair need not make
the static inverse singular.


% STATUS: conditional framework; outstanding model-specific work is explicit.
\section{P17: Linear transient and harmonic response}
\label{sec:P17}
The source is an outline. The following standard linear corollaries require
no further model assumptions; quantitative nonlinear response remains open.
\subsection{Transient}\label{slot:P17-1}
P04's single Fourier framework \cite[Theorem~3.1]{Gray2006} gives
\[
 e^{tJ}=Q\diag(e^{t\lambda_m})Q^*,\qquad
 \|e^{tJ}\|_2=e^{t\sa(J)}\quad(t\ge0).
\]
\subsection{Leading modes}\label{slot:P17-2}
Late-time dominance requires nonzero initial projection onto the leading
eigenspace; retain every tie. Euclidean norm decay is not coordinatewise
monotonicity or a propagation theorem.
\subsection{Explicit harmonic extension}\label{slot:P17-3}
For the added linear forcing $\dot z=Jz+be^{\ii\Omega t}$, substitution
gives $z_p=e^{\ii\Omega t}(\ii\Omega I-J)^{-1}b$ whenever the inverse
exists. If $J$ is Hurwitz, every solution approaches this harmonic response.
Real forcing and response are obtained by taking real parts.
\subsection{Unfinished nonlinear scope}\label{slot:P17-4}
Forcing conventions, observables, projected resonance amplitudes and
nonlinear approximation bounds remain to be specified and proved.


% STATUS: conditional framework; outstanding model-specific work is explicit.
\section{P18: Stability radius and Hopf shift (unfinished application)}
\label{sec:P18}
\subsection{Moving Jacobian}\label{slot:P18-1}
For P15's specified family set
$E(\eps)=D_xF(x^*(\eps),\eps)-J_0$. Bounds must include both direct
coefficient variation and equilibrium motion; these estimates are absent
from the supplied outline.
\subsection{Sufficient stability radius}\label{slot:P18-2}
If $J_0$ is normal and Hurwitz, put $\gamma=-\sa(J_0)>0$.
The normal-matrix perturbation inclusion
\cite[\S1 and Theorem~$4_\eps$]{EmbreeTrefethen2001} guarantees
stability for $\|E(\eps)\|_2<\gamma$, even when the perturbed matrix
is nonnormal. Turning this into a physical parameter radius requires
an explicit bound on $E$; it is not a sharp structured-radius assertion.
\subsection{Simple eigenvalue and threshold shift}\label{slot:P18-3}
Separately, at a simple transverse Hopf point of an analytic family,
\cite[Theorem~1]{GreenbaumLiOverton2020} and the implicit-function
theorem give
\[
 g_H'(0)=-\frac{\Ree[p^*J_\eps v]}{\Ree[p^*J_gv]},\qquad
 p^*v=1,\quad\Ree[p^*J_gv]\ne0.
\]
Here $J_\eps,J_g$ differentiate the Jacobian evaluated at the moving
equilibrium. These are not derivatives with the equilibrium frozen.
\subsection{Remaining Hopf conditions}\label{slot:P18-4}
Verify the family's regularity, remaining spectral gap, crossing and
first-Lyapunov-coefficient nondegeneracy. A spectral threshold shift alone
does not establish a nonlinear Hopf bifurcation.


\clearpage
\begin{thebibliography}{10}

\bibitem{Gray2006}
Robert~M. Gray.
\newblock Toeplitz and circulant matrices: A review.
\newblock {\em Foundations and Trends in Communications and Information
  Theory}, 2(3):155--239, 2006.
\newblock \url{https://ee.stanford.edu/~gray/toeplitz.pdf}.

\bibitem{GuckenheimerHolmes1983}
John Guckenheimer and Philip Holmes.
\newblock {\em Nonlinear Oscillations, Dynamical Systems, and Bifurcations of
  Vector Fields}, volume~42 of {\em Applied Mathematical Sciences}.
\newblock Springer, New York, 1983.
\newblock \url{https://doi.org/10.1007/978-1-4612-1140-2}.

\bibitem{vanKekemSterk2018}
Dirk~L. van Kekem and Alef~E. Sterk.
\newblock Travelling waves and their bifurcations in the {Lorenz-96} model.
\newblock {\em Physica D: Nonlinear Phenomena}, 367:38--60, 2018.
\newblock \url{https://arxiv.org/abs/1704.05442v3}.

\bibitem{Goh1977}
B.~S. Goh.
\newblock Global stability in many-species systems.
\newblock {\em The American Naturalist}, 111(977):135--143, 1977.
\newblock \url{https://doi.org/10.1086/283144}.

\bibitem{KuznetsovFold2006}
Yuri~A. Kuznetsov.
\newblock Saddle-node bifurcation.
\newblock {\em Scholarpedia}, 1(10):1859, 2006.
\newblock \url{https://doi.org/10.4249/scholarpedia.1859}.

\bibitem{KuznetsovHopf2006}
Yuri~A. Kuznetsov.
\newblock Andronov--hopf bifurcation.
\newblock {\em Scholarpedia}, 1(10):1858, 2006.
\newblock \url{https://doi.org/10.4249/scholarpedia.1858}.

\bibitem{Stewart2024}
Ian Stewart.
\newblock {Hopf} bifurcation and phase patterns in symmetric ring networks,
  2024.
\newblock \url{https://arxiv.org/html/2403.15814v1}.

\bibitem{GreenbaumLiOverton2020}
Anne Greenbaum, Ren-Cang Li, and Michael~L. Overton.
\newblock First-order perturbation theory for eigenvalues and eigenvectors.
\newblock {\em SIAM Review}, 62(2):463--482, 2020.
\newblock \url{https://arxiv.org/abs/1903.00785}.

\bibitem{BenziRazouk2007}
Michele Benzi and Nader Razouk.
\newblock Decay bounds and {$O(n)$} algorithms for approximating functions of
  sparse matrices.
\newblock {\em Electronic Transactions on Numerical Analysis}, 28:16--39, 2007.
\newblock \url{https://www.cs.emory.edu/~benzi/Web_papers/benzi_razouk.pdf}.

\bibitem{EmbreeTrefethen2001}
Mark Embree and Lloyd~N. Trefethen.
\newblock Generalizing eigenvalue theorems to pseudospectra theorems.
\newblock {\em SIAM Journal on Scientific Computing}, 23(2):583--590, 2001.
\newblock
  \url{https://people.maths.ox.ac.uk/~trefethen/publication/PDF/2001_94.pdf}.

\end{thebibliography}

\end{document}
